\chapter{Functions and graphs}
\label{vol02:ch02}

\epigraph{Look at the unknown. And try to think of a familiar
problem having the same or a similar unknown.}{George P\'olya, on the
discipline of restating a problem before solving it. The engineering
corollary applies to a graph: read the frame before reading the curve.}

\chapmeta{Half-life: foundational. Archetypes: scaling (log and exponential scales as the formal handling of order-of-magnitude reasoning from Volume I). Exercise target: 30. Project track: analysis only. Hazard class: none. Reader path default: standard, with sections explicitly tagged. Named cases: Space Shuttle Challenger (1986), where the pre-launch O-ring temperature plot omitted the incident-free flights and so erased the visual signal of cold sensitivity; briefly revisited from Volume I, the Mars Climate Orbiter (1999) and Air France Flight 447 (2009) where displayed quantities crossed an interface without their scale or unit.}

\noindent
A function is the formal record of how one quantity moves under
another. A graph is the picture of that record. An engineer reads the
picture before trusting it: which axis is linear, which is
logarithmic, where the origin sits, what units travel with the marks,
what scale the eye is being asked to compress. The chapter trains
that reading discipline alongside the working vocabulary of the
elementary functions: polynomials, rational functions, exponentials,
logarithms, and the trigonometric pair.

\section{Functions as mappings}\pathtag{core}

A \engterm{function} is a rule that assigns to each element of one
set, the \engterm{domain}, exactly one element of another, the
\engterm{codomain}. We write $f: D \to C$ to mean that $f$ takes an
input from $D$ and returns an output in $C$. The single-output
discipline is what makes a function useful: a procedure that returns
two values for the same input is called a \engterm{relation} and
needs different machinery.

Three pieces of information define a function: its domain, its
codomain, and the rule. A formula like $f(x) = \sqrt{x}$ is incomplete
on its own. The expression accepts non-negative real inputs and
returns non-negative real outputs by convention; an engineer
encountering it in a textbook supplies that domain mentally. In a
table, dataset, or program, the convention does not save the reader.
The domain has to be stated.

The \engterm{image} of $f$ is the set of values $f$ actually returns
as its input runs over the domain. The image is a subset of the
codomain; the two are equal only when $f$ is \engterm{surjective}
(``onto''). A function is \engterm{injective} (``one-to-one'') when
no two domain values share an output. A function that is both is
\engterm{bijective}, and a bijective function has an inverse, which
we develop in Section~2.5.

\subsection{Domain restriction is engineering, not mathematics}

A real engineering quantity rarely runs over all of $\R$. A pipe
diameter cannot be negative; a pH value lives in $[0, 14]$; a
temperature in kelvin cannot be below zero; a probability is bounded
in $[0, 1]$; an angle of attack on a wing has a stall ceiling beyond
which the model the function encodes does not apply. The domain of an
engineering function carries the physics of the problem.

A common defect in draft engineering reports is a function written as
if the domain were $\R$ when it is not. Extrapolating a polymer
viscosity formula calibrated between $20\,\degC$ and $80\,\degC$ down
to $-30\,\degC$ produces a number; the number is a calculation, not a
viscosity. The rule of carrying domain alongside formula, in the
same way Volume I asked the reader to carry units, will save the
reader the same kind of error.

The domain of an engineering function carries three kinds of
constraint at once, and the working engineer keeps all three in view.
The \emph{algebraic} constraint excludes inputs at which the formula
is undefined: a denominator zero, a negative argument under a
square root, a non-positive argument to a logarithm. The
\emph{physical} constraint excludes inputs the quantity cannot take:
a negative absolute temperature, a mass fraction above one, a
probability outside $[0, 1]$. The \emph{validity} constraint, the one
draft reports most often omit, excludes inputs outside the range over
which the model was calibrated or derived, even when the formula
returns a finite, physically plausible number there. The first two
constraints announce themselves; the third is silent, and the silence
is what makes extrapolation the most expensive of the three errors.

\paragraph{Worked example: the three domain constraints on one rule.}
The Antoine equation for the saturation vapour pressure of a liquid
is $\log_{10} p^{*} = A - B/(C + T)$, with $T$ a temperature in
$\si{\degreeCelsius}$ and $A$, $B$, $C$ tabulated constants. For
water with $A = 8.07131$, $B = 1730.63$, $C = 233.426$, three
domains apply. The algebraic domain excludes $T = -C =
-233.426\,\si{\degreeCelsius}$, where the denominator vanishes; that
point is far outside any liquid range, so the algebraic constraint
never binds in practice. The physical domain runs from the freezing
point to the critical point of water, since the liquid does not exist
outside that band. The validity domain is the narrowest and the one
that matters: this Antoine constant set is fitted to data over
$1\,\si{\degreeCelsius}$ to $100\,\si{\degreeCelsius}$, and a second
constant set covers $99\,\si{\degreeCelsius}$ to
$374\,\si{\degreeCelsius}$. An engineer who applies the low-range
constants at $150\,\si{\degreeCelsius}$ gets a number that is finite,
positive, and wrong by several percent, with no symptom in the
arithmetic. Carrying the validity domain alongside the formula is the
only defence, because the formula will not complain.

\paragraph{Worked example: domain of a composed expression.}
Consider $f(x) = \ln(4 - x^{2})$. The natural logarithm requires a
positive argument, so the domain is the set of $x$ for which
$4 - x^{2} > 0$, that is, $x^{2} < 4$, that is, $-2 < x < 2$. The
domain is the open interval $(-2, 2)$; the endpoints are excluded
because the argument reaches zero there and the logarithm diverges to
$-\infty$. The image runs from $-\infty$ (at the endpoints) up to
$\ln 4 \approx 1.386$ (at $x = 0$, where the argument is largest), so
the image is $(-\infty, \ln 4]$. Reading the domain off the inner
expression before evaluating the outer one is the habit that prevents
a program from calling $\ln$ on a negative number and returning a
not-a-number value that then propagates silently through the rest of
a calculation.

\paragraph{Worked example: domain and image of a real-world rule.}
A throttle position sensor reports its angle as a voltage $V$
proportional to $\sin\theta$ over the range $\theta \in [10^{\circ},
80^{\circ}]$. Three quantities are needed before the engineer can
use this function. The domain is the closed interval $[10^{\circ},
80^{\circ}]$, set by the mechanical end-stops of the throttle.
The codomain is the voltage range the sensor's output amplifier can
reach, $[0, 5]\,\si{\volt}$. The image, the values the sensor
actually returns over the domain, is $[V_{0} \sin 10^{\circ},
V_{0} \sin 80^{\circ}]$ for some calibration constant $V_{0}$;
if $V_{0} = 5\,\si{\volt}$, the image is approximately $[0.87,
4.92]\,\si{\volt}$. The image is a strict subset of the codomain
because the sensor never reads exactly zero (the geometry forbids
$\theta = 0$) and never exactly $5\,\si{\volt}$ (the geometry forbids
$\theta = 90^{\circ}$). An engine controller that treats $0\,\si
{\volt}$ as ``throttle closed'' is treating a value outside the
image as if it were inside, which means the controller will read
``throttle closed'' under exactly one condition: when the sensor has
failed and the line has gone to ground. The image-versus-codomain
distinction is the algebraic statement of that diagnostic.

% Figure: the function zoo plotted over a shared horizontal range
% so the reader can compare scaling behaviours on the same axes. A
% linear function, a quadratic, an exponential, a logarithm, and a
% sine all share x in [0, 6].
\begin{figure}[htbp]
\centering
\begin{tikzpicture}[scale=1.05,
  ax/.style={->, draw=engblack, thick},
  tick/.style={draw=engblack!60, thin},
  lin/.style={draw=engblack, thick},
  quad/.style={draw=engblue, thick},
  expo/.style={draw=engred, thick},
  loga/.style={draw=engamber, thick},
  trig/.style={draw=enggray, thick, dashed},
  ann/.style={font=\footnotesize, align=left},
]
% --- axes ---
\draw[ax] (0, 0) -- (6.5, 0) node[below right, ann] {$x$};
\draw[ax] (0, 0) -- (0, 5.0) node[left, ann] {$y$};
\foreach \x in {1, 2, 3, 4, 5, 6} {
  \draw[tick] (\x, -0.05) -- (\x, 0.05);
  \node[ann, below, yshift=-2pt] at (\x, 0) {\x};
}
\foreach \y in {1, 2, 3, 4} {
  \draw[tick] (-0.05, \y) -- (0.05, \y);
  \node[ann, left, xshift=-2pt] at (0, \y) {\y};
}

% --- linear: y = 0.5 x ---
\draw[lin, domain=0:6, samples=2, smooth] plot (\x, {0.5*\x});
\node[ann, anchor=south west, draw=engblack, fill=white, inner sep=2pt]
  at (5.2, 2.8) {$y = x/2$};

% --- quadratic: y = 0.12 x^2 ---
\draw[quad, domain=0:6, samples=40, smooth] plot (\x, {0.12*\x*\x});
\node[ann, anchor=south, draw=engblue, fill=white, inner sep=2pt]
  at (5.5, 4.4) {$y = x^{2}/8$};

% --- exponential: y = 0.2 e^{0.6 x} ---
\draw[expo, domain=0:5.5, samples=40, smooth] plot (\x, {0.2*exp(0.6*\x)});
\node[ann, anchor=west, draw=engred, fill=white, inner sep=2pt]
  at (5.5, 4.85) {$y = 0.2 e^{0.6 x}$};

% --- logarithm: y = ln(x + 1) ---
\draw[loga, domain=0:6, samples=40, smooth]
  plot (\x, {ln(\x + 1)});
\node[ann, anchor=south, draw=engamber, fill=white, inner sep=2pt]
  at (5.4, 2.1) {$y = \ln(x + 1)$};

% --- sinusoid: y = 1.5 + sin(x*180/pi) ---
\draw[trig, domain=0:6, samples=60, smooth]
  plot (\x, {1.5 + sin(\x * 57.2958)});
\node[ann, anchor=west, draw=enggray, fill=white, inner sep=2pt]
  at (0.3, 0.5) {$y = 1.5 + \sin x$};

\end{tikzpicture}
\caption[The elementary-function zoo on a shared range]{Five
elementary functions plotted on the same horizontal range $x \in [0,
6]$ so the reader can compare scaling behaviours on the same eyes.
The linear function $y = x/2$ rises by a constant slope. The
quadratic $y = x^{2}/8$ stays below the linear until $x = 4$, then
overtakes it. The exponential $0.2 e^{0.6 x}$ begins below both,
crosses the linear near $x = 3$, and outruns the quadratic by $x =
5$. The logarithm $\ln(x + 1)$ rises rapidly near the origin and
flattens. The sinusoid oscillates between $0.5$ and $2.5$ around
its mean $1.5$. The diagnostic discipline: knowing which of the five
shapes a dataset most resembles is the first move before fitting any
model.}
\label{fig:vol02:ch02:function-zoo}
\end{figure}


\Cref{fig:vol02:ch02:function-zoo} sets out the five elementary
families that the rest of the chapter develops, plotted on a single
horizontal range so the eye can compare them. The linear function
keeps a constant slope. The quadratic starts below the linear and
overtakes it at $x = 4$. The exponential starts below both and
outruns the quadratic by $x = 5$. The logarithm rises steeply near
the origin and flattens. The sinusoid stays bounded. Most engineering
data resembles one of these five shapes once the right axis transform
has been applied, and recognising the shape is the engineer's first
diagnostic move before any fitting begins.

\begin{archetype}[Scaling]
A function records how an output changes when an input changes. The
engineer asks a sharper question: how does the output change when the
input is scaled by a factor of two, or ten, or a thousand? Linear
functions double when their input doubles; quadratics quadruple;
cubic terms increase eightfold; exponentials are unrecognisable
between consecutive doublings; logarithms barely move. Recognising
the scaling behaviour of a function on sight is the engineer's first
diagnostic when reading data, sizing equipment, or checking a
colleague's argument. The Volume I order-of-magnitude habit becomes,
in Volume II, the habit of identifying the dominant scaling term.
\end{archetype}

\subsection{Composition, restriction, extension}

Two functions $f: A \to B$ and $g: B \to C$ \engterm{compose} as
$g \circ f: A \to C$, defined by $(g \circ f)(x) = g(f(x))$. The
composition is associative; it is rarely commutative. Engineers
encounter composition every time a sensor reading is converted to a
displayed quantity: the photodiode's photon flux first becomes a
current $i = f(\Phi)$, then a voltage $V = g(i)$, then a digital
count $n = h(V)$. The whole chain is $h \circ g \circ f$, and the
chain's behaviour is controlled by the link with the worst
linearity, the largest noise, or the narrowest dynamic range.

\engterm{Restriction} narrows a function's domain (the same rule, on
fewer inputs); \engterm{extension} broadens it. Engineering work is
mostly restriction: a manufacturer's datasheet curve is a function on
the temperature range the manufacturer characterised, and the engineer
restricts further to the range the device will see in service.
Extension demands evidence that the rule still applies; the absence
of that evidence is what makes extrapolation dangerous.

% Figure: schematic of composition and inverse. Three sets and two
% maps; the composition runs around the triangle and the inverse
% runs back along the diagonal.
\begin{figure}[htbp]
\centering
\begin{tikzpicture}[
  set/.style={draw=engblack, thick, ellipse,
    minimum width=2.0cm, minimum height=2.8cm},
  el/.style={circle, fill=engblack, inner sep=1.3pt},
  arrf/.style={->, thick, draw=engblack!85},
  arrg/.style={->, thick, draw=engblue},
  arrcomp/.style={->, very thick, draw=engred, dashed},
  ann/.style={font=\footnotesize, align=center},
  lbl/.style={font=\small\itshape, align=center},
]

% Three sets across the page
\node[set, label=above:{$A$}] (A) at (0, 0) {};
\node[set, label=above:{$B$}] (B) at (4.5, 0) {};
\node[set, label=above:{$C$}] (C) at (9.0, 0) {};

% Three points in A
\node[el] (a1) at (0, 0.7) {};
\node[el] (a2) at (0, 0) {};
\node[el] (a3) at (0, -0.7) {};

% Three points in B
\node[el] (b1) at (4.5, 0.7) {};
\node[el] (b2) at (4.5, 0) {};
\node[el] (b3) at (4.5, -0.7) {};

% Three points in C
\node[el] (c1) at (9.0, 0.7) {};
\node[el] (c2) at (9.0, 0) {};
\node[el] (c3) at (9.0, -0.7) {};

% f maps A to B
\draw[arrf] (a1) -- (b1);
\draw[arrf] (a2) -- (b2);
\draw[arrf] (a3) -- (b3);

% g maps B to C
\draw[arrg] (b1) -- (c1);
\draw[arrg] (b2) -- (c2);
\draw[arrg] (b3) -- (c3);

% Labels for f and g
\node[lbl] at (2.25, 1.6) {$f: A \to B$};
\node[lbl] at (6.75, 1.6) {$g: B \to C$};

% g o f composition: long arrow underneath
\draw[arrcomp] (a1) to[out=-30, in=-150] (c1);
\node[lbl, draw=engred, fill=white, inner sep=2pt]
  at (4.5, -2.0) {$g \circ f: A \to C$};

% f-inverse (return arrow above)
\draw[arrf, dotted, draw=engblack!60] (b2) to[out=110, in=70] (a2);
\node[ann] at (2.25, 1.0) {$f^{-1}$};

\end{tikzpicture}
\caption[Composition and inverse as directed maps between sets]{The
composition $g \circ f$ sends an element of $A$ to its image in $B$
under $f$ and then to its image in $C$ under $g$. The dashed red
arrow is the composed map read as a single function from $A$ to $C$.
The dotted arrow above shows the inverse of $f$ from $B$ back to
$A$, which exists exactly when $f$ is bijective. Composition is
associative ($h \circ (g \circ f) = (h \circ g) \circ f$) and is
rarely commutative. The inverse of the composition reverses both the
order and the individual inversions: $(g \circ f)^{-1} = f^{-1}
\circ g^{-1}$.}
\label{fig:vol02:ch02:composition-inverse}
\end{figure}


\Cref{fig:vol02:ch02:composition-inverse} draws the composition and
the inverse as directed maps between three sets. The composition
$g \circ f$ is read by following the arrows from $A$ to $B$ to $C$;
the inverse $f^{-1}$ is the reversal of the arrows on the single
map $f$ and exists only when $f$ pairs each output with exactly one
input. The picture is short and the algebra long, but the picture
fixes the order: $g \circ f$ reads ``apply $f$ then $g$,''
right-to-left in the symbol order and left-to-right in the diagram.
A standing source of error in engineering code is reading the symbol
order the wrong way, with $f$ and $g$ swapped; the diagram is the
shortest fix.

\paragraph{Worked example: a sensor-to-display chain.}
A pressure transducer outputs a current $i$ proportional to the
applied pressure $p$ by $i = f(p) = a p + i_{0}$, with sensitivity
$a = 16\,\si{\micro\ampere\per\pascal}$ and zero offset
$i_{0} = 4\,\si{\milli\ampere}$ (the standard 4--20\,mA loop). An
amplifier converts the current to a voltage by $V = g(i) = R\,i$
with $R = 250\,\si{\ohm}$. An ADC then digitises by
$n = h(V) = \lfloor V/V_{\text{LSB}} \rfloor$ with
$V_{\text{LSB}} = 5\,\si{\volt}/4096$.
The composition $h \circ g \circ f$ takes the engineering
quantity (pressure in pascals) to the displayed integer count
$n$, and is what the operator's screen reads from. Undoing the
display to recover the pressure runs the chain backwards:
$p = f^{-1}(g^{-1}(h^{-1}(n))) = (h^{-1}(n) \cdot 1/(Ra)) -
i_{0}/a$, where the inverses are taken in reverse order of the
forward composition. The order matters: applying $f^{-1}$ to the
voltage and $g^{-1}$ to the pressure would produce a number with
the wrong units and the wrong magnitude. The calibration sheet in
the field carries the forward chain explicitly so the inverse can
be applied without ambiguity; this is the same metrology
discipline Volume I Chapter 3 named for the calibration chain.

\subsection{Symmetry: even, odd, neither}

A function $f$ is \engterm{even} if $f(-x) = f(x)$ for every $x$ in
the domain (which is symmetric about zero), and \engterm{odd} if
$f(-x) = -f(x)$ on the same kind of domain. Graphically, an even
function is the mirror image of itself across the $y$-axis; an odd
function is its own point reflection through the origin. The
polynomials $x^{2k}$ are even, $x^{2k+1}$ are odd, the cosine is
even, the sine is odd. The product of two even functions is even;
the product of two odd functions is even; the product of an even
and an odd is odd. The sum of an even and an odd function is
neither.

% Figure: two panels showing the visual definition of even and odd
% functions; an even function is a mirror image across the y-axis, an
% odd function is a point reflection through the origin.
\begin{figure}[htbp]
\centering
\resizebox{\linewidth}{!}{%
\begin{tikzpicture}[scale=0.95,
  ax/.style={->, draw=engblack, thick},
  tick/.style={draw=engblack!60, thin},
  curve/.style={draw=engblue, thick},
  oddc/.style={draw=engred, thick},
  ann/.style={font=\footnotesize, align=center},
]

% --- Left panel: even, y = x^2 - 1 ---
\begin{scope}[xshift=0cm]
  \draw[ax] (-2.6, 0) -- (2.6, 0) node[right, ann] {$x$};
  \draw[ax] (0, -1.5) -- (0, 3.5) node[above, ann] {$y$};
  \foreach \x in {-2, -1, 1, 2} {
    \draw[tick] (\x, -0.05) -- (\x, 0.05);
  }
  \foreach \y in {-1, 1, 2, 3} {
    \draw[tick] (-0.05, \y) -- (0.05, \y);
  }
  \draw[curve, domain=-1.85:1.85, samples=40, smooth]
    plot (\x, {\x*\x - 1});
  \node[ann, anchor=south] at (1.65, 1.8) {$y = x^{2} - 1$};
  % Mirror cue: vertical dashed line
  \draw[dashed, draw=engblack!50] (0, -1.5) -- (0, 3.5);
  \node[ann, anchor=north, fill=white, inner sep=2pt]
    at (0, -1.7) {even: $f(-x) = f(x)$};
\end{scope}

% --- Right panel: odd, y = x^3 - 3x ---
\begin{scope}[xshift=8cm]
  \draw[ax] (-2.6, 0) -- (2.6, 0) node[right, ann] {$x$};
  \draw[ax] (0, -2.8) -- (0, 2.8) node[above, ann] {$y$};
  \foreach \x in {-2, -1, 1, 2} {
    \draw[tick] (\x, -0.05) -- (\x, 0.05);
  }
  \foreach \y in {-2, -1, 1, 2} {
    \draw[tick] (-0.05, \y) -- (0.05, \y);
  }
  \draw[oddc, domain=-2.05:2.05, samples=80, smooth]
    plot (\x, {\x*\x*\x - 3*\x});
  \node[ann, anchor=south] at (-1.5, 1.6) {$y = x^{3} - 3x$};
  % Symmetry cue: small inversion arrow through origin
  \draw[dashed, draw=engblack!50] (-2.4, -2.4) -- (2.4, 2.4);
  \node[ann, anchor=north, fill=white, inner sep=2pt]
    at (0, -3.0) {odd: $f(-x) = -f(x)$};
\end{scope}

\end{tikzpicture}%
}%
\caption[Even and odd symmetry by visual inspection]{Even functions
satisfy $f(-x) = f(x)$ and graph as a left-right mirror image across
the $y$-axis (left panel, the parabola $y = x^{2} - 1$). Odd
functions satisfy $f(-x) = -f(x)$ and graph as a point reflection
through the origin (right panel, the cubic $y = x^{3} - 3x$). An
arbitrary function decomposes uniquely into its even and odd parts,
$f(x) = \tfrac{1}{2}[f(x) + f(-x)] + \tfrac{1}{2}[f(x) - f(-x)]$, a
decomposition Volume II Chapter~13 lifts to the Fourier cosine and
sine series. Recognising parity on sight halves the integration work
for any symmetric or antisymmetric quantity (mean, kinetic energy,
overlap integral, autocorrelation).}
\label{fig:vol02:ch02:even-odd-symmetry}
\end{figure}


\Cref{fig:vol02:ch02:even-odd-symmetry} draws the two parities for
the parabola $y = x^{2} - 1$ (even, mirror across the $y$-axis) and
the cubic $y = x^{3} - 3x$ (odd, point reflection through the
origin). Any function with symmetric domain decomposes uniquely
into its even and odd parts,
\[
  f(x) = \underbrace{\tfrac{1}{2}\bigl[f(x) + f(-x)\bigr]}_{
    \text{even part}}
    + \underbrace{\tfrac{1}{2}\bigl[f(x) - f(-x)\bigr]}_{
    \text{odd part}},
\]
a decomposition Volume II Chapter 13 lifts to the Fourier cosine
and sine series. The engineering payoff is bookkeeping: a symmetric
integration $\int_{-a}^{a} f(x)\,\dd x$ vanishes for the odd part
and doubles for the even part, so identifying parity halves the
integration work for any symmetric or antisymmetric quantity (a
mean, a moment, a kinetic energy, an overlap integral, an
autocorrelation lag).

\begin{mastery}
The even-odd split is the first instance of a decomposition idea the
rest of the volume develops into one of its central tools. The
even-odd decomposition splits a function into two pieces under the
single symmetry $x \mapsto -x$; the Fourier decomposition (Volume II
Chapter 13) splits a periodic function into infinitely many pieces
under the family of symmetries that the sinusoids of every harmonic
frequency carry. The even part of a real periodic signal becomes its
cosine series, the odd part its sine series, so the parity split is
the two-term seed of the full harmonic expansion. The payoff
generalises too. A symmetric integral kills the odd part and doubles
the even part, and the orthogonality of the harmonics performs the
same bookkeeping term by term, which is why a Fourier coefficient is
computed by integrating the signal against a single harmonic and
reading off how much of that harmonic the signal contains. An
engineer who sees the even-odd split as the first harmonic
decomposition arrives at the Fourier series already knowing what it
is for: writing a complicated periodic function as a sum of simple
symmetric pieces, each of which the engineer can reason about alone.
\end{mastery}

\subsection{Transformations: shifts, scalings, reflections}

The elementary functions arrive as a small base set, and most of the
functions an engineer meets are those base functions shifted, scaled,
and reflected. Four operations generate the family from a base curve
$y = f(x)$, and reading them correctly is the difference between
sketching a transformed curve in seconds and grinding through a table
of points.

A \engterm{vertical shift} replaces $f(x)$ by $f(x) + c$, moving the
whole graph up by $c$ (down when $c < 0$). A \engterm{horizontal
shift} replaces $f(x)$ by $f(x - a)$, moving the graph right by $a$.
The horizontal case carries the trap: the inside operation moves the
graph the opposite way to the sign it appears to carry, because
$f(x - a)$ reaches at $x = a$ the value the base reached at the
origin. A \engterm{vertical scaling} replaces $f(x)$ by $k f(x)$,
stretching the graph away from the horizontal axis by the factor $k$
(compressing when $|k| < 1$, reflecting across the horizontal axis
when $k < 0$). A \engterm{horizontal scaling} replaces $f(x)$ by
$f(bx)$, compressing the graph toward the vertical axis by the factor
$b$ (the inside trap again: $b > 1$ compresses, $b < 1$ stretches).

% Figure: the four rigid and scaling transformations of a base curve.
% A single base parabola y = x^2 is shifted, reflected, and scaled to
% show how the four operations move the graph. The point is that the
% engineer reads a transformed curve by undoing the transformations
% in the right order.
\begin{figure}[htbp]
\centering
\begin{tikzpicture}[scale=0.95,
  ax/.style={->, draw=engblack, thick},
  tick/.style={draw=engblack!60, thin},
  base/.style={draw=engblack, thick},
  shiftc/.style={draw=engblue, thick},
  scalec/.style={draw=engred, thick, dashed},
  refl/.style={draw=engamber, thick, dash dot},
  ann/.style={font=\footnotesize, align=left},
]
% --- axes centred ---
\draw[ax] (-3.3, 0) -- (3.5, 0) node[below right, ann] {$x$};
\draw[ax] (0, -2.3) -- (0, 4.3) node[left, ann] {$y$};
\foreach \x in {-3, -2, -1, 1, 2, 3} {
  \draw[tick] (\x, -0.05) -- (\x, 0.05);
  \node[ann, below, yshift=-2pt] at (\x, 0) {$\x$};
}
\foreach \y in {-2, -1, 1, 2, 3, 4} {
  \draw[tick] (-0.05, \y) -- (0.05, \y);
  \node[ann, left, xshift=-2pt] at (0, \y) {$\y$};
}
% --- base: y = x^2 ---
\draw[base, domain=-2:2, samples=40, smooth] plot (\x, {\x*\x});
\node[ann, anchor=south west, draw=engblack, fill=white, inner sep=2pt]
  at (1.4, 3.6) {$y = x^{2}$};
% --- horizontal shift: y = (x+1.5)^2 ---
\draw[shiftc, domain=-3.3:0.5, samples=40, smooth]
  plot (\x, {(\x+1.5)*(\x+1.5)});
\node[ann, anchor=east, draw=engblue, fill=white, inner sep=2pt]
  at (-2.0, 3.0) {$y = (x+1.5)^{2}$};
% --- vertical scale + shift down: y = 0.5 x^2 - 1.5 ---
\draw[scalec, domain=-2.8:2.8, samples=40, smooth]
  plot (\x, {0.5*\x*\x - 1.5});
\node[ann, anchor=west, draw=engred, fill=white, inner sep=2pt]
  at (1.6, 1.4) {$y = \tfrac{1}{2}x^{2} - \tfrac{3}{2}$};
% --- reflection: y = -0.6 x^2 + 3 ---
\draw[refl, domain=-2.4:2.4, samples=40, smooth]
  plot (\x, {-0.6*\x*\x + 3});
\node[ann, anchor=south, draw=engamber, fill=white, inner sep=2pt]
  at (-2.6, 1.2) {$y = -\tfrac{3}{5}x^{2}+3$};
\end{tikzpicture}
\caption[The four transformations of a base curve]{One base parabola
$y = x^{2}$ (solid black) and three transforms of it. The blue curve
$y = (x + 1.5)^{2}$ shifts the base left by $1.5$ (the inside
operation moves the graph the opposite way to its sign). The red
dashed curve $y = \tfrac{1}{2}x^{2} - \tfrac{3}{2}$ scales the base
vertically by $\tfrac{1}{2}$ and shifts it down by $\tfrac{3}{2}$
(outside operations act in the intuitive direction). The amber
dash-dotted curve $y = -\tfrac{3}{5}x^{2} + 3$ reflects the base
across the horizontal and shifts it up by $3$. The reading rule:
inside operations on $x$ act horizontally and against their sign,
outside operations on the whole expression act vertically and with
their sign.}
\label{fig:vol02:ch02:transformations}
\end{figure}


\Cref{fig:vol02:ch02:transformations} draws one base parabola and
three transforms. The general second-degree curve $y = a(x - h)^{2}
+ k$ is the base $y = x^{2}$ scaled vertically by $a$ and shifted to
the vertex $(h, k)$; the engineer who reads $h$, $k$, and the sign of
$a$ off the equation knows the vertex, the opening direction, and the
width without plotting. The same parsing applies across the function
families. The sinusoid $A \sin(\omega t + \varphi) + D$ is the base
$\sin t$ scaled vertically by the amplitude $A$, scaled horizontally
by the angular frequency $\omega$, shifted horizontally by the phase
$\varphi$, and shifted vertically by the offset $D$. The exponential
$A e^{k(t - t_{0})}$ is the base $e^{t}$ scaled and shifted. Reading
a transformed curve is reading the four parameters off its standard
form.

The order of operations matters when scaling and shifting combine,
because the operations do not commute. The form $f(b(x - a))$ shifts
then scales, while $f(bx - c)$ with $c = ab$ describes the same curve
parsed as scale then shift. The reliable habit is to write the
transformed argument in the factored form $b(x - a)$, so the shift
$a$ and the scale $b$ read off directly and the order ambiguity does
not arise. A controller that applies a gain and an offset to a sensor
signal is performing exactly this transformation chain, and reading
the chain in the wrong order swaps the calibration offset with the
gain, an error that scales every reading rather than biasing it.

\paragraph{Worked example: parsing a damped sinusoid.}
A vibration sensor returns $y(t) = 3 + 2 e^{-t/4} \sin(6t - \pi/3)$
in millimetres, with $t$ in seconds. Parsing the transformations
reads the physics off the algebra. The constant $3$ is a vertical
offset: the signal oscillates about a mean of $3\,\si{\milli\meter}$,
which is the static deflection the structure carries under its own
weight. The factor $2 e^{-t/4}$ is a time-varying amplitude that
starts at $2\,\si{\milli\meter}$ and decays with time constant
$4\,\si{\second}$, so the oscillation has fallen to $1/e$ of its
initial size after $4\,\si{\second}$ and has fallen to about $5\,\%$
of it after $12\,\si{\second}$ (three time constants). The argument $6t - \pi/3 = 6(t - \pi/18)$ shows
angular frequency $\omega = 6\,\si{\radian\per\second}$, hence
ordinary frequency $f = 6/(2\pi) \approx 0.95\,\si{\hertz}$ and
period $T \approx 1.05\,\si{\second}$, with a phase shift of
$\pi/18 \approx 0.17\,\si{\second}$ in time. The engineer reads
mean, decay rate, frequency, and phase off the standard form without
plotting a single point, which is the working payoff of carrying the
transformation vocabulary.

\subsection{Fitting a function family and reading its parameters}

Transformation vocabulary runs in reverse when an engineer meets data
rather than an equation. The forward task takes a standard form and
draws the curve; the inverse task takes a curve and recovers the
standard form. The recovery has two stages that the working engineer
keeps separate. The first stage is choosing the \emph{family}, the
functional form whose shape the physics dictates: a saturating
approach, an exponential decay, a power law, a sinusoid. The second
stage is fitting the \emph{parameters} of the chosen family to the
data. The order is not negotiable, because a parameter fit is only as
honest as the family it sits inside, and a family chosen to flatter
the data rather than the physics produces parameters that mean
nothing.

% Figure: fitting a function family to data and reading its parameters.
% A saturating exponential approach y = y_inf (1 - e^{-t/tau}) fitted to
% scattered measurements; the three readable parameters (asymptote,
% time constant, initial slope) annotated on the curve.
\begin{figure}[htbp]
\centering
\resizebox{0.86\linewidth}{!}{%
\begin{tikzpicture}[scale=1.0,
  ax/.style={->, draw=engblack, thick},
  tick/.style={draw=engblack!60, thin},
  curve/.style={draw=engblue, very thick},
  asy/.style={draw=engblack!55, thick, dashed},
  tang/.style={draw=engred, thick, dash pattern=on 3pt off 2pt},
  pt/.style={circle, fill=engred, inner sep=1.1pt},
  ann/.style={font=\footnotesize, align=center},
]
  \draw[ax] (0, 0) -- (7.6, 0) node[below right, ann] {$t$ (s)};
  \draw[ax] (0, 0) -- (0, 4.6) node[left, ann] {$y$ (\si{\degreeCelsius})};
  \foreach \x in {1, 2, 3, 4, 5, 6, 7} {
    \draw[tick] (\x, -0.05) -- (\x, 0.05);
    \node[ann, below, yshift=-2pt] at (\x, 0) {\x};
  }
  \foreach \y in {1, 2, 3, 4} {
    \draw[tick] (-0.05, \y) -- (0.05, \y);
    \node[ann, left, xshift=-2pt] at (0, \y) {\y};
  }
  % Asymptote y_inf = 4
  \draw[asy] (0, 4) -- (7.4, 4);
  \node[ann, anchor=south east, fill=white, inner sep=1pt]
    at (7.4, 4.02) {$y_{\infty}$ (asymptote)};
  % Fitted curve y = 4 (1 - e^{-t/1.8})
  \draw[curve, domain=0:7.4, samples=60, smooth]
    plot (\x, {4*(1 - exp(-\x/1.8))});
  % Initial-slope tangent: slope = y_inf/tau = 4/1.8 = 2.22
  \draw[tang, domain=0:1.95, samples=2] plot (\x, {(4/1.8)*\x});
  \node[ann, anchor=west, text=engred] at (1.55, 3.5)
    {initial slope\\$y_{\infty}/\tau$};
  % Time-constant marker: at t = tau, y = y_inf(1 - 1/e) = 0.632 y_inf
  \draw[tick] ({1.8}, -0.05) -- ({1.8}, 0.05);
  \node[ann, below, yshift=-2pt, text=engblue] at ({1.8}, 0) {$\tau$};
  \draw[engblack!45, thin, dotted] ({1.8}, 0) -- ({1.8}, {4*0.632});
  \draw[engblack!45, thin, dotted] (0, {4*0.632}) -- ({1.8}, {4*0.632});
  \node[ann, anchor=west, fill=white, inner sep=1pt]
    at ({1.95}, {4*0.632}) {$0.63\,y_{\infty}$ at $t=\tau$};
  % Scattered data points
  \foreach \x/\y in {0.3/0.55, 0.7/1.25, 1.1/1.75, 1.6/2.30, 2.1/2.78,
    2.7/3.10, 3.3/3.42, 4.0/3.62, 4.8/3.80, 5.6/3.88, 6.5/3.95} {
    \node[pt] at (\x, \y) {};
  }
\end{tikzpicture}%
}%
\caption[Fitting a function family and reading its parameters]{A
first-order thermal step response fitted to scattered measurements
by the saturating-exponential family $y(t) = y_{\infty}(1 -
e^{-t/\tau})$. Three parameters read off the fitted curve without
further computation. The horizontal asymptote $y_{\infty}$ is the
steady-state value the response approaches. The time constant $\tau$
is the time at which the response has covered $1 - 1/e \approx 63\,\%$
of the gap to the asymptote. The initial slope $y_{\infty}/\tau$ is
the tangent at the origin, the rate at which the response would reach
the asymptote if it kept its starting pace. Choosing the family
first, then reading its parameters off the curve, is the working
order: the family encodes the physics (a single dominant relaxation
mechanism), and the parameters quantify it.}
\label{fig:vol02:ch02:family-fit}
\end{figure}


\Cref{fig:vol02:ch02:family-fit} works the recovery for a first-order
thermal step response. A thermocouple is dropped into a stirred bath
and its reading logged as it warms toward the bath temperature. The
physics is a single dominant relaxation, heat flowing across one
thermal resistance into one thermal mass, so the family is the
saturating exponential $y(t) = y_{\infty}(1 - e^{-t/\tau})$ with two
parameters. The asymptote $y_{\infty}$ reads off the curve as the
level the data approaches, here $4\,\si{\degreeCelsius}$ above the
start. The time constant $\tau$ reads off as the time at which the
response has covered $1 - 1/e \approx 63\,\%$ of the gap to the
asymptote, here about $1.8\,\si{\second}$. The initial slope
$y_{\infty}/\tau$ reads off as the tangent at the origin, the rate
the response would hold if it kept its starting pace. Three numbers,
each read off a single feature of the curve, and the model is
complete. The engineer who carries the transformation vocabulary
reads these parameters the way a musician reads a key signature,
before playing a note.

\paragraph{Worked example: choosing the family before the fit.}
A cooling-fin temperature is logged against distance $x$ from the
heated base, giving a column of $(x, T)$ pairs that fall on a smooth
declining curve. Three families could be drawn through the points: a
straight line $T = a - bx$, a power law $T = a x^{-n}$, and a decaying
exponential $T = T_{\infty} + (T_{0} - T_{\infty}) e^{-mx}$. The
straight line is wrong because it predicts a negative temperature at
finite distance, which the physics forbids. The power law is wrong
because it diverges at $x = 0$, where the data is finite and largest.
The exponential is right because the fin equation, a second-order
linear differential equation with constant coefficients (Volume II
Chapter 5), has exponential solutions, and the boundary condition at
the fin tip sets the offset $T_{\infty}$. Choosing the exponential
family before fitting is choosing the physics; the parameters
$m$ and $T_{\infty}$ then carry meaning ($m$ is the inverse decay
length, $1/m$ the distance over which the excess temperature falls by
$1/e$). The same data fitted to the wrong family would return numbers
that fit the points and explain nothing. Family first, parameters
second, every time.

\section{Polynomials and rational functions}\pathtag{core}

A \engterm{polynomial} of degree $n$ is a finite sum of integer-power
terms,

\[
p(x) = a_{n} x^{n} + a_{n-1} x^{n-1} + \cdots + a_{1} x + a_{0},
\]

with $a_{n} \neq 0$. Polynomials are the simplest non-trivial
functions and the easiest to compute. They arise in engineering as
fits to smooth data, as interpolants between tabulated values, as
power-series approximations of more complicated functions, and as the
characteristic functions of linear systems.

The qualitative behaviour of a polynomial is governed by its degree
and its leading coefficient. As $|x|$ becomes large, $p(x)$ behaves
like $a_{n} x^{n}$; the lower-order terms become negligible. A cubic
with positive leading coefficient runs from $-\infty$ at the left to
$+\infty$ at the right, passing through zero somewhere; a quartic
with positive leading coefficient is bounded below and unbounded
above. The engineer who has internalised these shapes can sketch a
polynomial's broad behaviour from its degree and leading coefficient
alone.

\subsection{Roots and the fundamental theorem}

The \engterm{roots} of $p$ are the values of $x$ for which $p(x) = 0$.
The fundamental theorem of algebra, which we state without proof,
guarantees that a polynomial of degree $n$ with complex coefficients
has exactly $n$ roots in the complex plane, counted with multiplicity.
For real polynomials of practical engineering size, complex roots
arrive in conjugate pairs.

The reader has met root-finding for quadratics through the formula

\[
x = \frac{-b \pm \sqrt{b^{2} - 4ac}}{2a}.
\]

The discriminant $b^{2} - 4ac$ tells the engineer the qualitative
character of the answer before computation: positive yields two real
roots (the quadratic crosses the axis twice), zero a single repeated
root (the quadratic touches the axis), negative two complex
conjugates (the quadratic does not cross). For higher-degree
polynomials, closed-form root expressions exist up to degree four
(due to Cardano and Ferrari) but are rarely used in engineering;
numerical root-finders treated in Volume II Chapter 16 are the
working tool.

\paragraph{Worked example: sketching a cubic from its roots.}
A cubic $p(x) = x^{3} - 7x + 6$ is to be sketched without a plotting
tool. The integer-root candidates are the divisors of the constant
term $6$: $\pm 1, \pm 2, \pm 3, \pm 6$. Testing, $p(1) = 1 - 7 + 6 =
0$, so $x = 1$ is a root and $(x - 1)$ divides $p$. Synthetic
division gives $p(x) = (x - 1)(x^{2} + x - 6) = (x - 1)(x + 3)
(x - 2)$. The three roots are $x = -3, 1, 2$. The leading coefficient
is positive, so the curve runs from $-\infty$ at the far left to
$+\infty$ at the far right. Between the roots the sign alternates: $p
< 0$ on $(-\infty, -3)$, $p > 0$ on $(-3, 1)$, $p < 0$ on $(1, 2)$,
$p > 0$ on $(2, \infty)$. The curve crosses upward at $x = -3$, peaks
in $(-3, 1)$, crosses downward at $x = 1$, troughs in $(1, 2)$, and
crosses upward at $x = 2$. Three roots, one sign pattern, and the
leading-coefficient rule produce a usable sketch with no derivative
calculation and no plotted points.

\paragraph{Worked example: a repeated root changes the crossing.}
Compare $p(x) = (x - 2)^{2}(x + 1)$ with the previous example. The
root at $x = 2$ has multiplicity two, the root at $x = -1$
multiplicity one. At a root of even multiplicity the curve touches
the axis and turns back without crossing; at a root of odd
multiplicity it crosses. So this cubic crosses upward through
$x = -1$, rises, and at $x = 2$ comes down to touch the axis and
bounces back up rather than passing through. The multiplicity of a
root is the local-shape information the bare root list omits, and it
is what distinguishes a contact point (a double root, a tangency, a
critically damped response) from a transversal crossing (a simple
root, a sign change, an underdamped overshoot). Reading multiplicity
is reading the local behaviour of the system the polynomial models.

\begin{estimation}
\textbf{Drag on a thrown object.}

\smallskip
\textit{Estimate, before reading on.}

\smallskip
A baseball ($D \approx 0.075 \,\si{\meter}$, $m \approx 0.145
\,\si{\kilogram}$) is thrown at $v = 40 \,\si{\meter\per\second}$
through air ($\rho \approx 1.2 \,\si{\kilogram\per\meter\cubed}$,
$\mu \approx 1.8 \times 10^{-5} \,\si{\pascal\second}$). The
aerodynamic drag is approximately $F = \tfrac{1}{2} C_{D} \rho v^{2}
A$ with $A = \pi D^{2}/4$ and $C_{D} \approx 0.35$ at the relevant
Reynolds number. Estimate the deceleration the drag produces. Is the
drag a quadratic function of speed, a linear function, or something
in between?

The drag scales as $v^{2}$: at this Reynolds number ($\mathrm{Re}
\approx 2 \times 10^{5}$, computed from $\rho v D / \mu$), inertial
effects dominate and $C_{D}$ is roughly constant. Substituting,

\[
F \approx \tfrac{1}{2} (0.35)(1.2)(40^{2})\bigl(\tfrac{\pi (0.075)^{2}}
{4}\bigr) \approx 1.5 \,\si{\newton},
\]

so the deceleration $a = F/m \approx 10 \,\si{\meter\per\second
\squared}$, comparable to gravity. A baseball decelerates noticeably
across a $20 \,\si{\meter}$ throw, and the deceleration is quadratic
in the speed. Doubling the speed quadruples the drag and the
deceleration. The engineer who plots ``ball speed at the plate'' as a
linear function of ``ball speed at release'' is missing the physics.

The estimate ignores spin (the Magnus force), the seam-roughness
transition that drops $C_{D}$ at higher Reynolds numbers (the drag
crisis), and altitude effects on $\rho$. Each is a refinement Volume
III Chapter 12 will treat.
\end{estimation}

\subsection{Rational functions}

A \engterm{rational function} is a ratio of polynomials, $r(x) =
p(x)/q(x)$, with $q$ not identically zero. The roots of $q$ are
\engterm{poles} of $r$: values where $r$ blows up. The engineer reads
a rational function's behaviour at three places: at its zeros (where
$p = 0$ and $q \neq 0$), at its poles, and at infinity.

The behaviour at infinity follows the leading-term ratio. For
$r(x) = (3x^{2} + x + 1)/(x^{2} - 4)$, at large $|x|$ the lower-order
terms vanish and $r(x) \to 3$. The horizontal line $y = 3$ is the
\engterm{horizontal asymptote}. For $r(x) = (x^{3} + 1)/(x^{2} + 1)$,
the leading-term ratio is $x$, so the function approaches the slanted
line $y = x$ at infinity, an \engterm{oblique asymptote}.

The behaviour near a pole follows the local algebra. Near a simple
pole at $x = x_{0}$, $r(x) \sim A/(x - x_{0})$ for some constant $A$,
and $r$ goes to $\pm \infty$ depending on the side from which $x$
approaches $x_{0}$. Near a double pole, $r \sim A/(x - x_{0})^{2}$
and goes to $+\infty$ from both sides if $A > 0$. The engineer who
sketches a rational function as a sequence of asymptotes, sign
changes, and pole behaviours can produce a usable picture without
plotting a single point.

% Figure: a rational function sketched from its features alone, with
% a vertical asymptote, a horizontal asymptote, a zero, and the sign
% behaviour across the pole. The point is that the engineer draws the
% curve from asymptotes and signs without plotting points.
\begin{figure}[htbp]
\centering
\begin{tikzpicture}[scale=0.95,
  ax/.style={->, draw=engblack, thick},
  tick/.style={draw=engblack!60, thin},
  asy/.style={draw=enggray, thin, dashed},
  curve/.style={draw=engblue, thick},
  ann/.style={font=\footnotesize, align=left},
]
% --- axes ---
\draw[ax] (-4.3, 0) -- (4.5, 0) node[below right, ann] {$x$};
\draw[ax] (0, -3.3) -- (0, 4.3) node[left, ann] {$y$};
\foreach \x in {-4, -3, -2, -1, 1, 2, 3, 4} {
  \draw[tick] (\x, -0.05) -- (\x, 0.05);
  \node[ann, below, yshift=-2pt] at (\x, 0) {$\x$};
}
\foreach \y in {-3, -2, -1, 1, 2, 3, 4} {
  \draw[tick] (-0.05, \y) -- (0.05, \y);
  \node[ann, left, xshift=-2pt] at (0, \y) {$\y$};
}
% r(x) = (x - 1)/(x + 2): vertical asymptote x=-2, horizontal y=1,
% zero at x=1.
% Vertical asymptote at x = -2
\draw[asy] (-2, -3.2) -- (-2, 4.2);
\node[ann, anchor=south west, color=enggray] at (-2, 4.0) {$x=-2$};
% Horizontal asymptote at y = 1
\draw[asy] (-4.2, 1) -- (4.4, 1);
\node[ann, anchor=south east, color=enggray] at (4.4, 1) {$y=1$};
% Left branch (x < -2): goes to +inf near asymptote, to 1 from above
% at -inf
\draw[curve, domain=-4.2:-2.25, samples=60, smooth]
  plot (\x, {(\x-1)/(\x+2)});
% Right branch (x > -2): comes from -inf, crosses zero at x=1,
% approaches 1 from below
\draw[curve, domain=-1.55:4.4, samples=80, smooth]
  plot (\x, {(\x-1)/(\x+2)});
% mark the zero
\node[circle, fill=engblue, inner sep=1.2pt] at (1, 0) {};
\node[ann, anchor=south west] at (1.05, 0.1) {zero $x=1$};
\node[ann, anchor=west, draw=engblue, fill=white, inner sep=2pt]
  at (1.8, 2.6) {$r(x) = \dfrac{x-1}{x+2}$};
\end{tikzpicture}
\caption[A rational function sketched from its features]{The rational
function $r(x) = (x - 1)/(x + 2)$ drawn from four features and no
plotted points. The vertical asymptote at $x = -2$ (the pole, where
the denominator vanishes) splits the plane into two branches. The
horizontal asymptote $y = 1$ (the ratio of leading coefficients) is
the limit at large $|x|$. The single zero at $x = 1$ is where the
numerator vanishes. The sign across the simple pole flips: the right
branch climbs from $-\infty$ just right of $x = -2$, through the zero
at $x = 1$, up toward $y = 1$ from below; the left branch falls from
$y = 1$ at the far left down toward $+\infty$ just left of the pole.
Four features fix the whole curve.}
\label{fig:vol02:ch02:rational-asymptotes}
\end{figure}


\Cref{fig:vol02:ch02:rational-asymptotes} draws a rational function
from its features alone. The pole sets the vertical asymptote, the
leading-coefficient ratio sets the horizontal asymptote, the
numerator zero sets the axis crossing, and the sign behaviour across
the simple pole completes the sketch. An engineer who has read those
four features has the curve; the plotted points add nothing the
features did not already fix.

\paragraph{Worked example: enzyme kinetics as a rational function.}
The Michaelis--Menten rate law for an enzyme-catalysed reaction is
$v(S) = V_{\max} S/(K_{m} + S)$, a rational function of the substrate
concentration $S$. Its three behaviours read off the algebra. At
$S \ll K_{m}$ the denominator is approximately $K_{m}$ and $v \approx
(V_{\max}/K_{m}) S$, linear in $S$: the enzyme is far from saturation
and the rate rises in proportion to substrate. At $S = K_{m}$ the
rate is exactly $V_{\max}/2$, which is the operational definition of
the Michaelis constant $K_{m}$ as the substrate concentration at
half-maximal rate. At $S \gg K_{m}$ the rate approaches the
horizontal asymptote $v \to V_{\max}$, the saturation plateau where
every enzyme site is occupied and adding substrate buys nothing. The
function has no real pole in the physical domain $S \geq 0$ because
$K_{m} > 0$, so the curve is the smooth rising-then-saturating shape
that the engineer recognises as a binding curve. The same rational
form, under different names, describes a Langmuir adsorption
isotherm, a saturating sensor, and a first-order receptor binding
curve; recognising the shape is recognising the saturation mechanism.

Rational functions are the transfer functions of linear time-invariant
systems, the response curves of resonant circuits, the gain curves of
filters, the kinetic equations of enzyme reactions. We will return to
their phasor and Laplace analyses in Volume II Chapter 3 and to their
control-system interpretation in Volume IX.

\begin{mastery}
A rational function $r = p/q$ with $\deg p < \deg q$ and distinct
poles decomposes into partial fractions, a sum of simple terms
$A_{i}/(x - x_{i})$ one per pole. The decomposition is the algebraic
engine behind the inverse Laplace transform (Volume II Chapter 3),
because each simple term inverts to a single exponential mode. The
poles $x_{i}$ are the system's natural frequencies; a pole with
positive real part is an unstable mode, one with negative real part a
decaying mode, one on the imaginary axis a sustained oscillation. The
whole stability theory of linear systems reads off the location of
the poles of a rational transfer function, which is why the engineer
who can find the poles of a rational function on sight has the first
tool of control engineering already in hand. Volume IX develops the
pole-placement and root-locus methods that this observation seeds.
\end{mastery}

\paragraph{Worked example: a low-pass filter's magnitude response.}
The first-order low-pass filter has the transfer function
$H(s) = 1/(1 + s/\omega_{c})$ with $s = \mathrm{j}\omega$ and cutoff
frequency $\omega_{c}$. The magnitude $|H(\omega)| = 1/\sqrt{1 +
(\omega/\omega_{c})^{2}}$ is a rational function of $\omega^{2}$
with no real poles. Its three behaviours are: at $\omega \ll
\omega_{c}$, $|H| \to 1$ (the passband); at $\omega = \omega_{c}$,
$|H| = 1/\sqrt{2} \approx 0.707$ ($-3\,\si{\deci\bel}$, the cutoff);
at $\omega \gg \omega_{c}$, $|H| \approx \omega_{c}/\omega$ (the
stopband). The last is the engineer's working result: each decade
above cutoff drops the magnitude by a decade, that is,
$-20\,\si{\deci\bel}$ per decade, a slope of $-1$ on the log-log
Bode plot. The Design exercise on the Bode plot is the worked
instance of this calculation.

\subsection{Piecewise functions, steps, and saturation}

Real engineering systems do not stay in the same algebraic form
across their full operating range. A function defined by different
rules on different sub-intervals of its domain is \engterm{piecewise};
the working engineer meets piecewise functions every day, in the form
of saturation, deadbands, hysteresis, regime transitions, and
switched circuit elements.

% Figure: a piecewise assembly. Three panels: Heaviside step,
% signum function, and a piecewise linear ramp-saturate (a soft
% saturator used in controller design).
\begin{figure}[htbp]
\centering
\resizebox{\linewidth}{!}{%
\begin{tikzpicture}[scale=0.92,
  ax/.style={->, draw=engblack, thick},
  tick/.style={draw=engblack!60, thin},
  curve/.style={draw=engblue, very thick},
  hole/.style={draw=engblue, thick, fill=white},
  filled/.style={fill=engblue, draw=engblue, thick},
  ann/.style={font=\footnotesize, align=center},
]

% --- Left: Heaviside step ---
\begin{scope}[xshift=0cm]
  \draw[ax] (-2.5, 0) -- (2.5, 0) node[right, ann] {$x$};
  \draw[ax] (0, -0.5) -- (0, 2.0) node[above, ann] {$H(x)$};
  \draw[tick] (-2, -0.05) -- (-2, 0.05);
  \draw[tick] (2, -0.05) -- (2, 0.05);
  \draw[tick] (-0.05, 1) -- (0.05, 1);
  \node[ann, left] at (-0.1, 1) {$1$};
  \draw[curve] (-2.4, 0) -- (-0.02, 0);
  \draw[curve] (0.02, 1) -- (2.4, 1);
  \node[hole, circle, inner sep=1.3pt] at (0, 0) {};
  \node[filled, circle, inner sep=1.3pt] at (0, 1) {};
  \node[ann, anchor=north] at (0, -0.7) {Heaviside step $H(x)$};
\end{scope}

% --- Middle: signum ---
\begin{scope}[xshift=6cm]
  \draw[ax] (-2.5, 0) -- (2.5, 0) node[right, ann] {$x$};
  \draw[ax] (0, -1.5) -- (0, 1.7) node[above, ann] {$\sgn(x)$};
  \draw[tick] (-2, -0.05) -- (-2, 0.05);
  \draw[tick] (2, -0.05) -- (2, 0.05);
  \draw[tick] (-0.05, 1) -- (0.05, 1);
  \draw[tick] (-0.05, -1) -- (0.05, -1);
  \node[ann, anchor=west] at (0.1, 1.0) {$+1$};
  \node[ann, anchor=west] at (0.1, -1.0) {$-1$};
  \draw[curve] (-2.4, -1) -- (-0.02, -1);
  \draw[curve] (0.02, 1) -- (2.4, 1);
  \node[filled, circle, inner sep=1.3pt] at (0, 0) {};
  \node[hole, circle, inner sep=1.3pt] at (0, 1) {};
  \node[hole, circle, inner sep=1.3pt] at (0, -1) {};
  \node[ann, anchor=north] at (0, -2.0) {signum $\sgn(x)$};
\end{scope}

% --- Right: ramp-saturate (deadband + soft saturator) ---
\begin{scope}[xshift=12cm]
  \draw[ax] (-2.5, 0) -- (2.5, 0) node[right, ann] {$x$};
  \draw[ax] (0, -1.5) -- (0, 1.7) node[above, ann] {$y$};
  \draw[tick] (1, -0.05) -- (1, 0.05);
  \draw[tick] (-1, -0.05) -- (-1, 0.05);
  \draw[tick] (-0.05, 1) -- (0.05, 1);
  \draw[tick] (-0.05, -1) -- (0.05, -1);
  % Saturate: y = -1 for x < -1; y = x for -1 <= x <= 1; y = 1 for x > 1
  \draw[curve] (-2.4, -1) -- (-1, -1) -- (1, 1) -- (2.4, 1);
  \node[ann, anchor=north] at (0, -2.0)
    {soft saturator\\(piecewise linear)};
\end{scope}

\end{tikzpicture}%
}%
\caption[Piecewise definitions assembled from elementary pieces]{Three
piecewise functions the working engineer meets daily. The Heaviside
step $H(x)$, by convention $0$ on the left and $1$ on the right
(value at zero a matter of convention; signal-processing conventions
choose $1/2$ for Fourier work and $1$ for control work). The signum
function $\sgn(x)$, the odd cousin of the step. The soft saturator,
a piecewise linear function that ramps through a central window
$[-1, 1]$ and clips outside it, the canonical model of an op-amp
output stage hitting the supply rails or a controller's actuator
limit. The filled circles mark the value the function actually takes
at a discontinuity; the open circles mark the limit the function
does not take. Sloppy notation about which value the function takes
at the discontinuity is a recurring source of error in published
control diagrams.}
\label{fig:vol02:ch02:piecewise-step}
\end{figure}


\Cref{fig:vol02:ch02:piecewise-step} draws the three piecewise
functions the working engineer meets most often. The
\engterm{Heaviside step} $H(x)$ takes the value zero on the negative
side and one on the positive side, and is the canonical model for a
switched input or an applied force. The \engterm{signum} function
$\sgn(x)$ returns the sign of its argument, takes the value zero at
zero, and is the canonical model for a Coulomb-friction force or a
sign-detection comparator. The \engterm{soft saturator}, a
piecewise-linear function that ramps through a central window and
clips outside it, is the canonical model for an actuator's
operating range, an op-amp output stage's rail limit, or the
linear region of any sensor before saturation.

The notation discipline for a piecewise definition is to state, for
each sub-interval, the inclusion of the endpoint. Writing
$f(x) = 1$ for $x \geq 0$ and $f(x) = 0$ for $x < 0$ is one
convention for the Heaviside step. Writing the alternative
$f(0) = 1/2$ is the Fourier-analysis convention because it makes
the inversion formula clean at the discontinuity. Two conventions,
both valid for their own purposes; the engineer specifies which is
meant when the choice could matter and the reader cannot guess.

\paragraph{Worked example: an op-amp output stage.}
A non-inverting amplifier with gain $A = 11$ and supply rails
$\pm 5\,\si{\volt}$ has the input-output relation
\[
  V_{\text{out}}(V_{\text{in}}) = \begin{cases}
    -5\,\si{\volt} & V_{\text{in}} < -5/11\,\si{\volt}, \\
    11\,V_{\text{in}} & -5/11\,\si{\volt} \leq V_{\text{in}}
      \leq 5/11\,\si{\volt}, \\
    +5\,\si{\volt} & V_{\text{in}} > +5/11\,\si{\volt}.
  \end{cases}
\]
The function is continuous (the cases match at the breakpoints)
but not differentiable at the breakpoints; the slope drops from
$11$ to zero on either side. A controller designed against the
linear gain alone will be in error outside the central window. The
piecewise model is the engineering correction; the linear model is
the small-signal approximation valid for $|V_{\text{in}}| < 5/11\,
\si{\volt}$.

\subsection{Dead band, hysteresis, and the limit of the function model}

Two nonlinearities beyond saturation recur often enough to name. A
\engterm{dead band} is a window of input near zero that produces no
output: a control valve that does not move until the actuator force
exceeds static friction, a gear train with backlash that takes up
slack before it transmits motion, a comparator with a threshold below
which it ignores its input. A dead band is still a single-valued
function, one output for each input, so the chapter's machinery
applies without change. The piecewise model carries a flat central
segment, a linear region on each side, and (often) saturation at the
extremes.

% Figure: two nonlinear input-output characteristics that a single-valued
% function cannot capture. Left: a dead-band-plus-saturation static
% characteristic (still a function). Right: a hysteresis loop, where the
% output depends on the direction of travel and so is NOT a function of
% the instantaneous input.
\begin{figure}[htbp]
\centering
\resizebox{\linewidth}{!}{%
\begin{tikzpicture}[scale=1.0,
  ax/.style={->, draw=engblack, thick},
  tick/.style={draw=engblack!60, thin},
  curve/.style={draw=engblue, very thick},
  up/.style={draw=engblue, very thick},
  down/.style={draw=engred, very thick},
  ann/.style={font=\footnotesize, align=center},
]

% --- Left: dead-band plus saturation (single-valued) ---
\begin{scope}[xshift=0cm]
  \draw[ax] (-3.2, 0) -- (3.2, 0) node[below right, ann] {input $u$};
  \draw[ax] (0, -2.2) -- (0, 2.2) node[left, ann] {output $y$};
  % dead-band [-0.8, 0.8]; linear region to +/-2.0; saturate at +/-1.6
  \draw[curve] (-3.0, -1.6) -- (-2.0, -1.6); % low saturation
  \draw[curve] (-2.0, -1.6) -- (-0.8, 0);    % rising through dead-band edge
  \draw[curve] (-0.8, 0) -- (0.8, 0);        % dead band (zero output)
  \draw[curve] (0.8, 0) -- (2.0, 1.6);       % rising
  \draw[curve] (2.0, 1.6) -- (3.0, 1.6);     % high saturation
  \draw[tick] (-0.8, -0.05) -- (-0.8, 0.05);
  \draw[tick] (0.8, -0.05) -- (0.8, 0.05);
  \node[ann, below, yshift=-2pt] at (0, -0.15) {dead band};
  \node[ann, anchor=south west, fill=white, inner sep=1pt] at (2.0, 1.6)
    {saturation};
  \node[ann, anchor=north] at (0, -2.5)
    {(a) dead band $+$ saturation\\(a function: one $y$ per $u$)};
\end{scope}

% --- Right: hysteresis loop (not single-valued) ---
\begin{scope}[xshift=8.5cm]
  \draw[ax] (-3.2, 0) -- (3.2, 0) node[below right, ann] {input $u$};
  \draw[ax] (0, -2.2) -- (0, 2.2) node[left, ann] {output $y$};
  % rising branch shifted right; falling branch shifted left
  \draw[up] (-3.0, -1.6) -- (-2.0, -1.6) -- (1.2, 1.6) -- (3.0, 1.6);
  \draw[down] (3.0, 1.6) -- (2.0, 1.6) -- (-1.2, -1.6) -- (-3.0, -1.6);
  % arrows on branches
  \draw[up, -{Stealth[length=4pt]}] (-0.6, 0.0) -- (-0.4, 0.1);
  \draw[down, -{Stealth[length=4pt]}] (0.6, 0.0) -- (0.4, -0.1);
  \node[ann, text=engblue, anchor=south east] at (1.0, 1.2)
    {increasing $u$};
  \node[ann, text=engred, anchor=north west] at (-1.0, -1.2)
    {decreasing $u$};
  \node[ann, anchor=north] at (0, -2.5)
    {(b) hysteresis loop\\(not a function: $y$ depends on history)};
\end{scope}

\end{tikzpicture}%
}%
\caption[Dead band, saturation, and hysteresis]{Two nonlinear
input-output characteristics. Panel (a): a dead band (a window of
input near zero that produces no output, common in valves, gear
trains with backlash, and threshold comparators) followed by a linear
region and saturation at the rails. The characteristic is still a
single-valued function, one output for each input, so the chapter's
machinery applies directly. Panel (b): a hysteresis loop, where the
output on the rising branch differs from the output on the falling
branch at the same input value. A magnetic material's $B$-$H$ curve,
a thermostat with a switching differential, and a mechanism with
static friction all show this shape. Hysteresis is not a function of
the instantaneous input: the output depends on the direction and
history of travel, so it requires a state variable beyond the input
to predict. Recognising that a characteristic has crossed from
function to history-dependent map is the first step in modelling it
correctly.}
\label{fig:vol02:ch02:hysteresis-deadband}
\end{figure}


\Cref{fig:vol02:ch02:hysteresis-deadband} sets a dead-band
characteristic beside a \engterm{hysteresis} loop, and the contrast
is the point. A hysteresis loop assigns a different output to the
same input depending on whether the input is rising or falling: a
magnetic material's $B$-$H$ curve retains magnetisation after the
field is removed, a thermostat with a switching differential turns on
at one temperature and off at a lower one, a mechanism with static
friction sticks before it slips. The rising branch and the falling
branch cross the same input value at two different output values, so
the relation fails the single-output test that defines a function.
Hysteresis is not a function of the instantaneous input. The output
depends on the direction and history of travel, which means a faithful
model carries a \engterm{state variable} beyond the input, and the
input-output description becomes a small dynamical system rather than
a static curve.

\paragraph{Worked example: a thermostat's switching differential.}
A room thermostat is set to $20\,\si{\degreeCelsius}$ with a
switching differential of $1\,\si{\degreeCelsius}$. The heater turns
on when the temperature falls to $19.5\,\si{\degreeCelsius}$ and off
when it rises to $20.5\,\si{\degreeCelsius}$. At exactly
$20.0\,\si{\degreeCelsius}$ the heater state is undetermined by the
temperature alone: it is on if the room is cooling through that point
on its way down from a recent off-switch, and off if the room is
warming through that point on its way up from a recent on-switch. The
output (heater on or off) is a function of the input (temperature)
plus one bit of state (the last switch direction). The differential
is engineering, not defect: a zero-differential thermostat would
chatter, switching thousands of times an hour as sensor noise crossed
the single setpoint. The hysteresis loop's width is the design margin
that buys switching stability at the cost of a steady-state
temperature band. Recognising that the characteristic has crossed
from function to history-dependent map is the first modelling
decision, and it changes the tool from algebra to state machine.

\begin{mastery}
The boundary between a function and a hysteresis map is the boundary
between a memoryless and a stateful system, and it recurs across
engineering. A resistor is memoryless: its current is a function of
its present voltage. A capacitor is stateful: its current depends on
the history of its voltage through the charge it has accumulated. The
function model of Section~2.1 describes the memoryless case exactly
and the stateful case only instantaneously, frozen at one point in
its history. When a characteristic shows hysteresis, the engineer's
move is to find the hidden state, the variable whose value
distinguishes the rising branch from the falling one. For the
thermostat it is the last switch direction; for the magnetic material
it is the remanent magnetisation; for the sticking mechanism it is
whether static or kinetic friction is acting. Naming the state turns
a multivalued relation back into a function on an enlarged domain
that includes the state, which is the move Volume IX makes formally
when it writes a system as a state-space model $\dot{x} = f(x, u)$,
$y = g(x, u)$ with the output $y$ a genuine function of state $x$ and
input $u$ together.
\end{mastery}

\section{Exponentials and logarithms}\pathtag{core}

The \engterm{exponential function} with base $a > 0$, $a \neq 1$, is

\[
f(x) = a^{x},
\]

defined for all real $x$. Two bases dominate engineering: $a = e
\approx 2.71828$, the natural exponential, used wherever a quantity's
rate of change is proportional to its current value; and $a = 10$,
used wherever a quantity ranges over many orders of magnitude.

The exponential function has the defining property
$a^{x + y} = a^{x} \cdot a^{y}$. Combined with $a^{0} = 1$, this rule
means the function turns addition into multiplication: a unit step in
$x$ multiplies the output by a fixed factor $a$. That property makes
exponentials the right model for compound interest, radioactive
decay, charging and discharging capacitors, bacterial growth in
unrestricted media, drug clearance, RC and RL circuit transients,
chemical reaction rates with Arrhenius temperature dependence, and
sound and signal attenuation through homogeneous media.

The \engterm{logarithm} of base $a$ is the inverse of the
exponential of base $a$:

\[
y = \log_{a}(x) \iff a^{y} = x,
\]

defined for $x > 0$. The natural logarithm $\ln(x) = \log_{e}(x)$
and the common logarithm $\log_{10}(x)$ are the two the engineer
uses without thinking. Logarithms turn multiplication into addition,
which is what makes them useful for decades of magnitude.

The change-of-base rule,

\[
\log_{a}(x) = \frac{\log_{b}(x)}{\log_{b}(a)},
\]

lets the reader move between bases without losing accuracy. The
constants the reader should know on sight: $\ln(10) \approx 2.303$,
$\log_{10}(e) \approx 0.4343$, $\log_{10}(2) \approx 0.301$,
$\log_{10}(3) \approx 0.477$. Two further numbers that recur:
$2^{10} = 1024 \approx 10^{3}$, so ten doublings is roughly three
decades; and $\ln(2) \approx 0.693$, the half-life-to-time-constant
conversion.

\subsection{Working in base $e$ and base $10$}

The two bases serve different engineering purposes.

Base $e$ is the natural language of growth and decay. If a quantity
$N(t)$ satisfies $\dd N / \dd t = k N$ (a relation Volume II Chapter 5
treats formally), the solution is $N(t) = N_{0} e^{kt}$. The constant
$k$ has units of inverse time and is the \engterm{rate constant}; its
reciprocal $\tau = 1/k$ is the \engterm{time constant}, the time over
which $N$ changes by a factor $e \approx 2.718$. Engineers report
exponential decays in time constants because $\tau$ has a direct
physical meaning: after one time constant the response is at $1/e
\approx 37\,\%$ of its initial value, after three at $5\,\%$, after
five at $0.7\,\%$. The half-life $t_{1/2} = \tau \ln 2 \approx 0.693
\tau$ is the alternative report; both are useful, and the engineer
should be able to convert in one line.

Base $10$ is the natural language of scale comparison. The
\engterm{decade} is a factor of ten; an instrument that operates over
six decades runs from, say, $10^{-3}$ to $10^{3}$ in some unit. The
\engterm{decibel} (Volume I Chapter 2) is one tenth of a base-$10$
log ratio. The pH scale is base $10$. The Richter and moment
magnitude scales for earthquakes are base $10$ in energy. The
astronomical magnitude scale is base $10$ raised to a fractional
power. Wherever an engineering quantity ranges over orders of
magnitude, the engineer's first move is to take the base-$10$
logarithm and look at the exponent.

\paragraph{Worked example: a doubling cadence as an exponential.}
A class of integrated circuits has been observed to double its
transistor count roughly every two years. Treating the count as a
continuous exponential $N(t) = N_{0} \cdot 2^{t/T_{2}}$ with doubling
time $T_{2} = 2\,\text{yr}$, the engineer converts to base $e$ and
base $10$ as needed. The continuous rate constant is $k = \ln 2/T_{2}
= 0.693/2 = 0.347\,\text{yr}^{-1}$, so $N(t) = N_{0} e^{0.347 t}$. On
a base-$10$ semilog plot the slope is $k \log_{10} e = 0.347 \cdot
0.4343 = 0.151$ decades per year, which says the count crosses a
factor of ten every $1/0.151 \approx 6.6\,\text{yr}$, consistent
with the rule that ten doublings ($2^{10} \approx 10^{3}$) is three
decades and therefore $3 \times 6.6 \approx 20\,\text{yr}$ for three
decades, or one decade every $6.6\,\text{yr}$. The estimate doubles
as a sanity check: a quantity that doubles every two years multiplies
by about $32$ in a decade, and $\log_{10} 32 \approx 1.5$ decades in
ten years matches the $0.151 \times 10$ the slope predicts. The
exponential model holds only while the underlying mechanism persists;
the historical doubling cadence in semiconductor scaling, current as
of the mid-2020s, has slowed as device dimensions approach atomic
limits, which is the validity-domain caveat the failure section will
generalise.

\paragraph{Worked example: combining quantities across decades.}
A logarithm turns a product into a sum, which is what makes
order-of-magnitude arithmetic fast. To estimate the energy released
by a magnitude-7 earthquake relative to a magnitude-5, the engineer
uses that each unit of moment magnitude corresponds to a factor of
$10^{1.5} \approx 31.6$ in radiated energy. Two units of magnitude is
$10^{3} = 1000$ times the energy, computed by adding exponents
($1.5 + 1.5 = 3$) rather than multiplying the factors ($31.6 \times
31.6 \approx 1000$). The same addition rule converts a chain of gains
in a signal path: a $+20\,\si{\deci\bel}$ amplifier followed by a
$-6\,\si{\deci\bel}$ attenuator and a $+14\,\si{\deci\bel}$ stage
gives a total of $20 - 6 + 14 = 28\,\si{\deci\bel}$ by addition,
where the linear-scale calculation would multiply $10 \times 0.5
\times 5 = 25$ in power ratio. Adding logarithms is the working form
of multiplying magnitudes, and it is why engineers reach for the
decibel and the decade whenever quantities span ranges the eye cannot
hold linearly.

\subsection{Semilog and log-log axes}

A function plotted on linear axes shows changes by addition. A
function plotted on logarithmic axes shows changes by multiplication.
Two cases recur, and they are not interchangeable.

% Figure: an exponential decay plotted on linear axes (left) and on
% semilog axes (right). The same data, the same physics; the semilog
% version reads the time constant as a slope and the initial value
% as the intercept.
\begin{figure}[htbp]
\centering
\resizebox{\linewidth}{!}{%
\begin{tikzpicture}[scale=1.0,
  ax/.style={->, draw=engblack, thick},
  tick/.style={draw=engblack!60, thin},
  curve/.style={draw=engred, thick},
  fitline/.style={draw=engblue, thick, dashed},
  pt/.style={circle, fill=engred, inner sep=1.2pt},
  ann/.style={font=\footnotesize, align=center},
]

% --- Left: linear ---
\begin{scope}[xshift=0cm]
  \draw[ax] (0, 0) -- (5.8, 0) node[below right, ann] {$t$ (s)};
  \draw[ax] (0, 0) -- (0, 4.2) node[left, ann] {$V$ (V)};
  \foreach \x in {1, 2, 3, 4, 5} {
    \draw[tick] (\x, -0.05) -- (\x, 0.05);
    \node[ann, below, yshift=-2pt] at (\x, 0) {\x};
  }
  \foreach \y in {1, 2, 3, 4} {
    \draw[tick] (-0.05, \y) -- (0.05, \y);
    \node[ann, left, xshift=-2pt] at (0, \y) {\y};
  }
  % V(t) = 4 e^{-t/1.5}, scaled to fit
  \draw[curve, domain=0:5.5, samples=40, smooth]
    plot (\x, {4*exp(-\x/1.5)});
  % Sampled points
  \foreach \x/\y in {0/4.0, 0.5/2.87, 1.0/2.05, 1.5/1.47,
    2.0/1.05, 2.5/0.75, 3.0/0.54, 3.5/0.39, 4.0/0.28, 4.5/0.20, 5.0/0.14} {
    \node[pt] at (\x, \y) {};
  }
  \node[ann, anchor=south] at (3.0, 3.5) {(a) linear axes};
\end{scope}

% --- Right: semilog (vertical log) ---
\begin{scope}[xshift=8cm]
  \draw[ax] (0, 0) -- (5.8, 0) node[below right, ann] {$t$ (s)};
  \draw[ax] (0, -0.5) -- (0, 4.2) node[left, ann] {$\log_{10} V$};
  \foreach \x in {1, 2, 3, 4, 5} {
    \draw[tick] (\x, -0.05) -- (\x, 0.05);
    \node[ann, below, yshift=-2pt] at (\x, 0) {\x};
  }
  % Major decades; log V scaled so 1 unit = one decade
  \foreach \y/\yl in {0/0, 1/1, 2/2, 3/3, 4/4} {
    \draw[tick] (-0.05, \y) -- (0.05, \y);
    \node[ann, left, xshift=-2pt] at (0, \y) {$10^{\yl}$};
  }
  % Plot log10(4 e^{-t/1.5}) shifted vertically by +2 so it stays
  % visible: y_axis = 2 + log10(4) + (-t/1.5)*log10(e)
  % = 2.602 - 0.2895 * t
  \draw[fitline, domain=0:5.5, samples=2]
    plot (\x, {2.602 - 0.2895*\x});
  \foreach \x in {0, 0.5, 1.0, 1.5, 2.0, 2.5, 3.0, 3.5, 4.0, 4.5, 5.0} {
    \pgfmathsetmacro{\yy}{2.602 - 0.2895*\x}
    \node[pt] at (\x, \yy) {};
  }
  \node[ann, anchor=south] at (3.0, 3.5) {(b) semilog axes};
  \node[ann, anchor=west, draw=engblue, fill=white, inner sep=2pt]
    at (1.4, 3.0)
    {slope $= -\log_{10}(e)/\tau$\\$\tau$ from slope};
\end{scope}

\end{tikzpicture}%
}%
\caption[An exponential decay on linear and semilog axes]{The same
exponential decay $V(t) = V_{0} e^{-t/\tau}$ with $V_{0} =
4\,\si{\volt}$ and $\tau = 1.5\,\si{\second}$, plotted on linear
axes (left) and on a vertical-log axis (right). On linear axes the
curve is the recognisable convex decay that the eye reads as ``fast
at first, then slow''. On semilog axes the same data is a straight
line, with intercept $\log_{10} V_{0}$ and slope $-\log_{10}(e)/\tau
\approx -0.290\,\si{\per\second}$. The engineer who fits a straight
line to the semilog data recovers $\tau$ from the slope in one step
($\tau = -\log_{10}(e) / \text{slope}$) and $V_{0}$ from the
intercept in another. A straight-line fit on the linear panel
recovers nothing.}
\label{fig:vol02:ch02:semilog-decay}
\end{figure}


\Cref{fig:vol02:ch02:semilog-decay} shows the same exponential decay
plotted on linear and on semilog axes. The two are the same data;
the semilog version converts the curve into a straight line, from
which the time constant reads off the slope and the initial value
reads off the intercept. A pair of plots like this is one of the
most reused diagnostics in working engineering, because almost every
first-order process the engineer meets, RC discharge, isotope decay,
drug clearance, thermal cooling, capacitor leakage, has the same
shape.

A \engterm{semilog plot} has one logarithmic axis (almost always the
vertical) and one linear axis. An exponential function $y = A e^{kx}$
becomes a straight line on a semilog plot:

\[
\log_{10} y = \log_{10} A + (k \log_{10} e) x,
\]

with slope $k \log_{10} e \approx 0.4343 k$ and intercept $\log_{10}
A$. A radioactive decay, an RC discharge, a population in unrestricted
growth, or a thermistor's resistance against temperature all read as
straight lines on semilog paper. The engineer who sees a straight
line on a semilog plot can extract the time constant from the slope
and the initial value from the intercept in two lines.

A \engterm{log-log plot} has both axes logarithmic. A power-law
function $y = A x^{n}$ becomes a straight line on a log-log plot:

\[
\log_{10} y = \log_{10} A + n \log_{10} x,
\]

with slope $n$ (the exponent of the power law) and intercept
$\log_{10} A$. The drag coefficient versus Reynolds number, a fatigue
$S$-$N$ curve, the Manning equation for open-channel flow, and the
linewidth of a transistor's $1/f$ noise all read as straight lines
(at least in regime) on log-log paper. The engineer who fits a power
law on log-log axes is reading exponents directly off the slope.

The discipline. Before reading a curve on log axes, ask three
questions. Is the data exponential, in which case semilog is
correct? Is the data a power law, in which case log-log is correct?
Or is the curve plotted on log axes only because the data spans
decades, in which case the slope on the log plot has no special
meaning and a fit must be motivated by physics, not by visual
straightness?

\paragraph{Worked example: the Arrhenius transform.}
A negative-temperature-coefficient (NTC) thermistor's resistance
follows the Steinhart approximation
$R(T) = R_{0} \exp[B(1/T - 1/T_{0})]$ with $T$ in kelvin and $B$ a
material constant \cite{v2c02-paper:steinhart-hart-1968}. The
function is neither exponential in temperature (because the
exponent depends on $1/T$, not $T$) nor a power law in temperature.
On linear axes the data is a smooth declining curve from which $B$
is hard to extract by eye. On \engterm{Arrhenius axes}, plotting
$\ln(R/R_{0})$ on the vertical and $1/T$ on the horizontal, the
function becomes a straight line of slope $B$ and intercept zero
through the point $(1/T_{0}, 0)$.
% Figure: a thermistor calibration. Two panels: resistance vs.
% temperature on linear axes (left) and on Arrhenius-style axes (ln
% R vs 1/T) with a straight-line fit (right). The slope of the fit
% gives the B-constant of the thermistor model.
\begin{figure}[htbp]
\centering
\resizebox{\linewidth}{!}{%
\begin{tikzpicture}[scale=1.0,
  ax/.style={->, draw=engblack, thick},
  tick/.style={draw=engblack!60, thin},
  curve/.style={draw=engred, thick},
  fitline/.style={draw=engblue, thick, dashed},
  pt/.style={circle, fill=engred, inner sep=1.2pt},
  ann/.style={font=\footnotesize, align=center},
]

% --- Left: linear axes, R vs T (in degrees C) ---
\begin{scope}[xshift=0cm]
  \draw[ax] (0, 0) -- (6.0, 0) node[below right, ann] {$T$ ($\degC$)};
  \draw[ax] (0, 0) -- (0, 4.2) node[left, ann] {$R$ (k$\Omega$)};
  \foreach \xc/\xl in {0/0, 1/10, 2/20, 3/30, 4/50, 5/80} {
    \draw[tick] (\xc, -0.05) -- (\xc, 0.05);
    \node[ann, below, yshift=-2pt] at (\xc, 0) {\xl};
  }
  \foreach \y/\yl in {1/10, 2/30, 3/100, 4/300} {
    \draw[tick] (-0.05, \y) -- (0.05, \y);
    \node[ann, left, xshift=-2pt] at (0, \y) {\yl};
  }
  % Data points (R, T): T=0, 10, 20, 25, 30, 50, 80
  % R = 10 k * exp(B*(1/T_K - 1/298))   with B = 3800 K
  % Plotted at scaled coordinates
  \foreach \x/\y in {0/3.8, 1/3.3, 2/2.7, 2.4/2.45, 3/2.05, 4/1.25, 5/0.65} {
    \node[pt] at (\x, \y) {};
  }
  \draw[curve, thick] plot[smooth] coordinates
    {(0, 3.8) (1, 3.3) (2, 2.7) (2.4, 2.45) (3, 2.05) (4, 1.25) (5, 0.65)};
  \node[ann, anchor=south] at (3.5, 3.6) {(a) $R$ vs $T$, linear};
\end{scope}

% --- Right: Arrhenius axes, ln(R/R_0) vs 1/T ---
\begin{scope}[xshift=8cm]
  \draw[ax] (0, 0) -- (5.5, 0)
    node[below right, ann] {$1000 / T$ (K$^{-1}$)};
  \draw[ax] (0, -1) -- (0, 4.0) node[left, ann] {$\ln(R/R_{25})$};
  \foreach \xc/\xl in {0.5/2.8, 1.5/3.0, 2.5/3.2, 3.5/3.4, 4.5/3.6} {
    \draw[tick] (\xc, -0.05) -- (\xc, 0.05);
    \node[ann, below, yshift=-2pt] at (\xc, 0) {\xl};
  }
  \foreach \y/\yl in {-1/-1, 0/0, 1/1, 2/2, 3/3} {
    \draw[tick] (-0.05, \y) -- (0.05, \y);
    \node[ann, left, xshift=-2pt] at (0, \y) {\yl};
  }
  % Linear in 1/T: B*(1/T - 1/T_ref) ~ slope * (x - x_ref)
  % Fit slope ~ 3.8 K, so ln(R/R_25) ~ slope*(1000/T - 1000/298)
  % Plot 7 points on a line of slope 3.8 with offsets
  \foreach \x/\y in {0.5/-0.5, 1.5/-0.1, 2.5/0.6, 3.0/1.0, 3.5/1.5, 4.5/2.5} {
    \node[pt] at (\x, \y) {};
  }
  \draw[fitline] (0.3, -0.65) -- (4.7, 2.65);
  \node[ann, anchor=west, draw=engblue, fill=white, inner sep=2pt]
    at (0.6, 3.2) {fit: $\ln(R/R_{25}) = B(1/T - 1/T_{25})$\\
    slope $\approx B \approx 3800$ K};
  \node[ann, anchor=south] at (3.5, 3.6) {(b) Arrhenius axes};
\end{scope}

\end{tikzpicture}%
}%
\caption[Thermistor calibration: linear vs Arrhenius axes]{A
negative-temperature-coefficient (NTC) thermistor's resistance falls
with temperature according to the Steinhart approximation
$R(T) = R_{25} \exp\bigl[B(1/T - 1/T_{25})\bigr]$, with $T$ in
kelvin, $R_{25}$ the resistance at $25\,\degC$, and $B$ a
material constant. On linear axes (a) the data is a smooth decreasing
curve from which $B$ is hard to read. On Arrhenius axes (b),
plotting $\ln(R/R_{25})$ against $1/T$, the same data is a straight
line whose slope is exactly $B$. The example here uses a typical
$10\,\si{\kilo\ohm}$ NTC bead at $25\,\degC$ with $B = 3800\,\si
{\kelvin}$, calibrated across $0\,\degC$ to $80\,\degC$. The
chapter's design exercise (Section~2.7) asks the reader to set up
this same plot; the figure is a worked instance of the answer.}
\label{fig:vol02:ch02:thermistor-fit}
\end{figure}

\Cref{fig:vol02:ch02:thermistor-fit} shows a typical
$10\,\si{\kilo\ohm}$ NTC bead with $B = 3800\,\si{\kelvin}$
calibrated across $0\,\degC$ to $80\,\degC$ on both axis pairs. On
the right panel the data falls on a straight line; the slope
recovers $B$ in one step, and a residual plot (not shown here)
reveals departures from the two-parameter Steinhart model. The
Design exercise on thermistor calibration asks the reader to set
up exactly this plot from the data file
\texttt{data/thermistor\_calibration.csv}; the worked instance in
the figure is the answer.

\paragraph{Worked example: identifying a power law from a dataset.}
A pump manufacturer's bench log records the power $P$ a centrifugal
pump draws against its rotational speed $N$ over a speed range from
$600$ to $3000\,\text{rpm}$, five decades-worth of points across a
factor of five in speed. Plotted on linear axes the data is a steep
rising curve from which no parameter reads cleanly. The model is to
be identified, not assumed. Taking $\log_{10} P$ against $\log_{10}
N$ straightens the data into a line, which says the relationship is a
power law $P = A N^{n}$. The slope read off the log-log line is the
exponent $n$. For this pump the slope comes out close to $3$: a
doubling of speed multiplies the power draw by $2^{3} = 8$. The
exponent is the diagnostic. The affinity laws for centrifugal pumps
predict $P \propto N^{3}$ from dimensional reasoning (flow scales as
$N$, head as $N^{2}$, and power as their product times efficiency),
so a measured slope of $3$ confirms the pump obeys the affinity laws
across the tested range, and a measured slope that departed from $3$
would flag a regime where the laws break, cavitation at the high end
or laminar losses at the low end. The intercept $\log_{10} A$ fixes
the prefactor once the exponent is known, $A = 10^{\text{intercept}}$.
The reading order is exponent from slope, prefactor from intercept,
physics from the value of the exponent. A log-log plot turns a
two-parameter fit into two straight-line readings and turns the
exponent into a statement about the mechanism.

\begin{mastery}
The elementary families of this chapter do not exhaust the functions
an engineer meets, and the working vocabulary extends to a short list
of \engterm{special functions} that arise as solutions to the
canonical differential equations of physics and as the limits of
common processes. The error function $\operatorname{erf}(x) =
(2/\sqrt{\pi}) \int_{0}^{x} e^{-t^{2}}\,\dd t$ is the running integral
of the Gaussian and describes diffusion fronts, the spread of a
heat pulse, and the cumulative of the normal distribution (Volume I
Chapter 6). The gamma function $\Gamma(z)$ extends the factorial to
non-integer and complex argument and sits inside the chi-squared,
gamma, and beta distributions. The Bessel functions $J_{n}(x)$ are
the radial modes of a vibrating drumhead, a cylindrical waveguide,
and a circular optical aperture (Volume IV). The Lambert $W$ of the
previous box inverts $x e^{x}$. None of these has an expression in
elementary functions, yet each is as concrete as the sine once the
engineer knows its defining equation, its asymptotic behaviour, and
the library call that evaluates it. The working stance is the same one
the chapter takes toward the elementary functions: carry the shape,
the limiting behaviour, and the domain, and reach for the numerical
implementation when a value is needed. A special function is a named
shape with a known graph and a tabulated evaluator, not a wall.
\end{mastery}

\begin{estimation}
\textbf{How many real digits in a curve fit?}

\smallskip
\textit{Estimate, before reading on.}

\smallskip
A junior engineer fits the exponential decay $V = V_{0} e^{-t/\tau}$
to fifty data points spread over five time constants, with the
voltage measurement accurate to $\pm 1\,\%$ at each sample. They
report the recovered $\tau$ as $1.5124\,\si{\second}$, citing four
significant figures. How many of those four digits are real?

The signal range covers $V_{0}/e^{5} \approx 0.0067 V_{0}$ to
$V_{0}$, so the data spans roughly $\log_{10}(e^{5}) \approx 2.2$
decades. The error in the fitted slope on a semilog plot scales as
$\sigma_{\text{slope}} \approx \sigma_{y}/(\sqrt{n}\,\Delta x)$,
where $\sigma_{y}$ is the noise in $\ln y$ (which is the relative
noise in $y$, here $0.01$), $\Delta x$ is the range in time, and
$\sqrt{n}$ is the $\sqrt{50}$ improvement from averaging. With
$\Delta x = 5\tau$,
\[
  \sigma_{\text{slope}}/\text{slope}
  \approx \frac{0.01}{\sqrt{50} \cdot 5\tau \cdot (1/\tau)}
  = \frac{0.01}{\sqrt{50} \cdot 5}
  \approx 3 \times 10^{-4}.
\]
So $\tau$ is known to roughly three parts per ten thousand, that is,
the third decimal place in $1.5124$ is real and the fourth is
optimistic. A working report would state $\tau = 1.512 \pm 0.001\,
\si{\second}$ or $\tau = 1.51\,\si{\second}$ and reserve the fourth
digit for a careful uncertainty analysis.

The estimate ignores the bias that linear-on-log fitting introduces
when the noise is additive rather than multiplicative; that bias
sits in the third decimal place too, so the conclusion is unchanged
for the multiplicative-noise case but should be redone if the noise
model is wrong. The script
\texttt{code/semilog\_fit.py} runs the two fits on a synthetic
dataset and prints both recovered values; running it ten times with
different random seeds gives an empirical sigma that should match
this estimate to within a factor of two. A fit reported to more
significant figures than the data supports is a standing source of
false precision in published engineering work.
\end{estimation}

\section{Trigonometric functions}\pathtag{standard}

The trigonometric functions $\sin$, $\cos$, and $\tan$ encode the
geometry of rotation. We treat them in Chapter 3 as the engineer's
tool for oscillation, phase, and complex exponentials; here we
collect the function-and-graph facts the rest of the volume uses.

The \engterm{sine} and \engterm{cosine} are bounded periodic
functions defined on all of $\R$, with range $[-1, 1]$ and period
$2\pi$ in radians (or $360^{\circ}$ in degrees). They differ by a
quarter-period shift: $\sin(x) = \cos(x - \pi/2)$. The Pythagorean
identity $\sin^{2}(x) + \cos^{2}(x) = 1$ is the algebraic statement
that they describe a point on the unit circle. The reader who carries
that identity has the lever for most trigonometric simplification.

The \engterm{tangent} $\tan(x) = \sin(x)/\cos(x)$ is a rational
function of the other two. It is unbounded, with vertical asymptotes
where $\cos(x) = 0$ (at $x = \pi/2 + n\pi$ for integer $n$). It
records slope: the slope of the line from the origin to a point on
the unit circle at angle $x$ is $\tan(x)$. Engineers meet $\tan$ in
slope calculations, in the small-angle approximation $\tan(x) \approx
x$ (good to $1\,\%$ for $x \lesssim 10^{\circ}$, treated in Volume II
Chapter 4), and in surveying where elevation differences come from
horizontal distance times $\tan$ of the angle of inclination.

\subsection{Phase, frequency, and wavelength}

The general sinusoid is

\[
y(t) = A \sin(\omega t + \varphi),
\]

with \engterm{amplitude} $A$, \engterm{angular frequency} $\omega$
in $\si{\radian\per\second}$, and \engterm{phase} $\varphi$ in
radians. The ordinary frequency $f = \omega/(2\pi)$ is in $\si
{\hertz}$; the period $T = 1/f$ in $\si{\second}$. Two waves with the
same $\omega$ but different $\varphi$ are said to differ in phase;
phase difference is the lever Chapter 3 uses to combine sinusoids
algebraically through complex exponentials.

In space rather than time, a travelling wave reads $y(x, t) = A
\sin(kx - \omega t + \varphi)$, where $k = 2\pi/\lambda$ is the
\engterm{wavenumber} in $\si{\radian\per\meter}$ and $\lambda$ is
the \engterm{wavelength} in $\si{\meter}$. The phase velocity is
$v_{p} = \omega/k = f \lambda$. Volume IV will use these forms for
electromagnetic and acoustic waves.

\paragraph{Worked example: a surveyor's tower height.}
A surveyor stands $50.0\,\si{\meter}$ horizontally from the base of
a transmission tower and measures the angle of elevation to the
top as $42.3^{\circ}$ with a digital theodolite. The eye height of
the theodolite above ground is $1.55\,\si{\meter}$. The tower
height above ground is
$h = 1.55\,\si{\meter} + 50.0\,\si{\meter} \cdot \tan(42.3^{\circ})
= 1.55 + 50.0 \cdot 0.9083 = 1.55 + 45.41 = 46.96\,\si{\meter}$.
The error in the recovered height is dominated by the angle
measurement; with theodolite resolution $\pm 0.05^{\circ} \approx
\pm 8.7 \times 10^{-4}\,\si{\radian}$ and horizontal distance
$d$, the derivative $\dd h/\dd \theta = d \sec^{2}\theta = 50.0
\cdot 1.83 = 91.4\,\si{\meter\per\radian}$ gives
$\sigma_{h} = 91.4 \cdot 8.7 \times 10^{-4} \approx 0.08\,\si
{\meter}$. The recovered height $46.96 \pm 0.08\,\si{\meter}$
matches the precision the theodolite was bought for. The point of
the example is the discipline of the chain: the engineer reports
$h$ to two decimals because the input chain supports two decimals,
not four. The arc-tangent branch trap, however, is real elsewhere:
a tower behind the surveyor produces the same $\tan$ value as one
in front (the principal branch loses the sign of the horizontal
distance), and $\mathrm{atan2}(\Delta y, \Delta x)$ is the
disambiguating function.

\subsection{Combining sinusoids of the same frequency}

Two sinusoids of the same frequency add to a single sinusoid of that
frequency. The identity is the workhorse of phasor analysis, and it
deserves to be stated and used before Chapter 3 builds the complex
machinery around it. For $a \cos(\omega t) + b \sin(\omega t)$, the
sum is
\[
  a \cos(\omega t) + b \sin(\omega t)
  = R \cos(\omega t - \varphi),
\]
with amplitude $R = \sqrt{a^{2} + b^{2}}$ and phase $\varphi =
\mathrm{atan2}(b, a)$. The amplitude is the Pythagorean combination
of the two components, and the phase is the angle whose tangent is
$b/a$ with the quadrant resolved by the signs of $a$ and $b$. The
result says that the in-phase and quadrature parts of a signal
combine like the legs of a right triangle, with $R$ the hypotenuse.

\paragraph{Worked example: two-source interference.}
Two loudspeakers driven at the same frequency reach a listener with
pressures $p_{1}(t) = 3 \cos(\omega t)$ and $p_{2}(t) = 4
\sin(\omega t)$ in pascals. The combined pressure is $p(t) = 3
\cos(\omega t) + 4 \sin(\omega t) = R \cos(\omega t - \varphi)$ with
$R = \sqrt{3^{2} + 4^{2}} = 5\,\si{\pascal}$ and $\varphi =
\mathrm{atan2}(4, 3) \approx 0.927\,\si{\radian} \approx
53.1^{\circ}$. The two sources combine to a single tone of amplitude
$5\,\si{\pascal}$, larger than either alone because the
quarter-period phase offset between a cosine and a sine puts them
partly in step. Had the second source been $-4 \sin(\omega t)$
instead, the amplitude would still be $5\,\si{\pascal}$ but the phase
would flip sign, which is the algebra of constructive and destructive
interference that Volume IV develops for waves. The amplitude never
exceeds $|3| + |4| = 7\,\si{\pascal}$ (full constructive) nor falls
below $|4| - |3| = 1\,\si{\pascal}$ for these magnitudes; the actual
$5\,\si{\pascal}$ sits between, set by the phase relationship.

\section{Inverse functions}\pathtag{standard}

A function $f: A \to B$ has an \engterm{inverse} $f^{-1}: B \to A$
if and only if $f$ is bijective. The defining property is
$f^{-1}(f(x)) = x$ for $x \in A$ and $f(f^{-1}(y)) = y$ for $y \in B$.
On a graph, the curves of $f$ and $f^{-1}$ are mirror images across
the line $y = x$.

% Figure: a function and its inverse as mirror images across y = x.
% The exponential e^x and the natural logarithm ln x, reflected
% across the diagonal. The point is that inversion is reflection in
% the line y = x, and the domain/range swap under reflection.
\begin{figure}[htbp]
\centering
\begin{tikzpicture}[scale=0.95,
  ax/.style={->, draw=engblack, thick},
  tick/.style={draw=engblack!60, thin},
  expo/.style={draw=engred, thick},
  loga/.style={draw=engblue, thick},
  mirror/.style={draw=enggray, thin, dashed},
  ann/.style={font=\footnotesize, align=left},
]
% --- axes ---
\draw[ax] (-3.3, 0) -- (4.3, 0) node[below right, ann] {$x$};
\draw[ax] (0, -3.3) -- (0, 4.3) node[left, ann] {$y$};
\foreach \x in {-3, -2, -1, 1, 2, 3, 4} {
  \draw[tick] (\x, -0.05) -- (\x, 0.05);
  \node[ann, below, yshift=-2pt] at (\x, 0) {$\x$};
}
\foreach \y in {-3, -2, -1, 1, 2, 3, 4} {
  \draw[tick] (-0.05, \y) -- (0.05, \y);
  \node[ann, left, xshift=-2pt] at (0, \y) {$\y$};
}
% mirror line y = x
\draw[mirror] (-3.2, -3.2) -- (4.2, 4.2);
\node[ann, anchor=south west, color=enggray] at (3.4, 3.4) {$y=x$};
% exponential y = e^x
\draw[expo, domain=-3:1.45, samples=50, smooth] plot (\x, {exp(\x)});
\node[ann, anchor=west, draw=engred, fill=white, inner sep=2pt]
  at (1.1, 3.6) {$y = e^{x}$};
% natural log y = ln x (reflection of the above)
\draw[loga, domain=0.05:4.2, samples=60, smooth] plot (\x, {ln(\x)});
\node[ann, anchor=north west, draw=engblue, fill=white, inner sep=2pt]
  at (3.4, 1.0) {$y = \ln x$};
% paired points (0,1) and (1,0)
\node[circle, fill=engblack, inner sep=1.0pt] at (0, 1) {};
\node[circle, fill=engblack, inner sep=1.0pt] at (1, 0) {};
\draw[mirror] (0,1) -- (1,0);
\end{tikzpicture}
\caption[A function and its inverse mirror across $y = x$]{The
exponential $y = e^{x}$ (red) and its inverse the natural logarithm
$y = \ln x$ (blue) are mirror images across the diagonal $y = x$
(grey dashed). The reflection swaps the roles of input and output: a
point $(0, 1)$ on the exponential maps to the point $(1, 0)$ on the
logarithm, and the short grey segment shows the reflection. The
domain of the exponential (all of $\R$) becomes the range of the
logarithm, and the range of the exponential ($(0, \infty)$) becomes
the domain of the logarithm. Reflection across $y = x$ is what
inversion looks like, and the domain/range swap is what it means.}
\label{fig:vol02:ch02:inverse-mirror}
\end{figure}


\Cref{fig:vol02:ch02:inverse-mirror} shows the exponential and the
logarithm as mirror images across the diagonal. The reflection makes
the domain/range swap visible: the part of the plane the exponential
reaches becomes the part the logarithm departs from. Reading the
mirror picture is a fast check on whether a proposed inverse is
right; a curve and its claimed inverse that are not reflections of
each other across $y = x$ are not an inverse pair.

The standard pairings the reader must know on sight:

\begin{itemize}[leftmargin=*]
  \item Squaring on $[0, \infty)$ and the principal square root.
  \item The exponential $e^{x}$ on $\R$ and the natural logarithm
    $\ln$ on $(0, \infty)$.
  \item The base-$10$ exponential and the common logarithm.
  \item Sine on $[-\pi/2, \pi/2]$ and the principal $\arcsin$ on
    $[-1, 1]$.
  \item Cosine on $[0, \pi]$ and the principal $\arccos$ on $[-1, 1]$.
  \item Tangent on $(-\pi/2, \pi/2)$ and the principal $\arctan$ on
    $\R$.
\end{itemize}

The trigonometric inverses require domain restriction because the
trigonometric functions are not bijective on their full domain. The
engineer's habit: when an inverse trigonometric function is reported,
state which branch is meant. A surveyor's $\arctan$ on a calculator
returns a value in $(-\pi/2, \pi/2)$, which is wrong by $\pi$ for
half of the bearings the surveyor cares about; a $\mathrm{atan2}(y,
x)$ function with both arguments returns the correct quadrant. The
distinction has cost a published paper its bearings before now.

\subsection{The inverse of a composition}

If $f$ and $g$ are bijective, the inverse of their composition is the
composition of their inverses in reversed order:

\[
(g \circ f)^{-1} = f^{-1} \circ g^{-1}.
\]

The order matters because composition is not commutative. The
engineering picture: undoing a chain of measurements (sensor,
amplifier, ADC, software calibration) requires reversing each step
and reversing the order of reversal. Volume I Chapter 3's calibration
chain is the same idea seen through metrology eyes.

\paragraph{Worked example: inverting a calibration curve.}
A flow meter's calibration is $Q = a \sqrt{\Delta p} + b$, with flow
$Q$ in $\si{\liter\per\minute}$, differential pressure $\Delta p$ in
$\si{\kilo\pascal}$, slope $a = 12.0\,\si{\liter\per\minute}/\sqrt{
\si{\kilo\pascal}}$, and offset $b = 2.0\,\si{\liter\per\minute}$.
The forward function maps pressure to flow; the field engineer needs
the inverse, mapping a measured flow back to the pressure the
transmitter reports. Solving $Q = a \sqrt{\Delta p} + b$ for $\Delta
p$: subtract $b$, divide by $a$, square. The inverse is $\Delta p =
((Q - b)/a)^{2}$, valid for $Q \geq b$ because the square root in the
forward map returns only non-negative values. The domain restriction
$Q \geq b$ is the engineering content of the inverse: a reported flow
below the offset $b$ cannot have come from this calibration and
signals a fault, a zero-drift, or operation outside the characterised
range. The inverse exists only on the image of the forward map, and
the boundary of that image is a diagnostic.

\paragraph{Worked example: composing and inverting a sensor chain.}
A two-stage temperature channel composes a thermistor with a
linearising amplifier. The thermistor maps temperature to resistance
by $R = g(T) = R_{0} e^{B(1/T - 1/T_{0})}$ (Section~2.3), and the
amplifier maps resistance to output voltage by an affine rule
$V = h(R) = \alpha R + \beta$. The full forward characteristic is the
composition $V = (h \circ g)(T) = \alpha R_{0} e^{B(1/T - 1/T_{0})} +
\beta$, read right-to-left: apply $g$ to get resistance, then $h$ to
get voltage. The field task is the inverse, recovering temperature
from a measured voltage. By the composition-inverse rule the inverse
is $T = (h \circ g)^{-1}(V) = (g^{-1} \circ h^{-1})(V)$, undone in
reversed order. First undo the amplifier, $R = h^{-1}(V) = (V -
\beta)/\alpha$; then undo the thermistor, $T = g^{-1}(R) =
\bigl(1/T_{0} + (1/B)\ln(R/R_{0})\bigr)^{-1}$. Composing the two
inverses gives the calibration formula a logger applies to every
reading. The order trap is concrete here: applying $g^{-1}$ to the
voltage (treating a voltage as a resistance) returns a temperature
that is finite, plausible, and wrong, with no symptom in the
arithmetic. Each inverse exists because each stage is monotone on the
operating range, the thermistor strictly decreasing in $T$ and the
amplifier strictly increasing in $R$ (assuming $\alpha > 0$), so the
composition is strictly monotone and the channel is invertible end to
end.

\begin{mastery}
A function has a single-valued inverse on an interval if and only if
it is strictly monotone there, and strict monotonicity is the
condition the engineer actually checks, because it is local and
verifiable where global bijectivity is neither. For a differentiable
function the working test is the sign of the derivative: a derivative
that keeps one sign across the interval certifies strict monotonicity
and so invertibility. A derivative that changes sign marks a turning
point, and at a turning point the function folds back on itself and
loses its inverse, which is exactly why the square root requires the
restriction to $[0, \infty)$ (where $x^{2}$ is increasing) and why
$\arcsin$ requires the restriction to $[-\pi/2, \pi/2]$ (where $\sin$
is increasing). The engineering consequence is sharp: a sensor whose
characteristic has a turning point inside its operating range cannot
be inverted to recover the measurand near that point, because two
distinct inputs produce the same output and the reading is ambiguous.
A thermocouple operated across the temperature where its
Seebeck-coefficient curve doubles back, or a flow meter operated
across a reversal, returns a value that does not name a unique state.
The remedy is to restrict the operating range to a monotone segment,
which is the design rule that keeps every well-designed sensor on one
side of its turning points.
\end{mastery}

\begin{mastery}
Not every useful function has a closed-form inverse. The map $y = x
e^{x}$ has an inverse, the Lambert $W$ function, with no expression in
elementary functions; it arises in delay-differential equations, in
the solution of $x = a^{x}$, and in the kinetics of certain
saturating processes. When a closed-form inverse is unavailable the
engineer inverts numerically: a bisection or Newton solve on
$f(x) - y = 0$ recovers $x$ from $y$ to any required precision,
provided $f$ is monotone on the relevant interval so the inverse is
single-valued. The script \texttt{code/bisection\_root.py} is the
general tool, and the monotonicity check is the precondition that
makes the numerical inverse well-defined. The existence of an inverse
is a property of the function (bijectivity on the domain of
interest); the means of computing it (formula or numerical solve) is
a separate, practical question.
\end{mastery}

\section{Reading a graph: shape, slope, asymptote, scale}\pathtag{core}

A graph is a compressed claim about a function. The engineer's
first job is to decompress it correctly. The reader will see hundreds
of plots in a working week: in datasheets, in textbooks, in journal
papers, in standards documents, in vendor whitepapers, in screening
software, in the dashboards of running systems. Reading a plot is
not optional. Reading it wrong, repeatedly, is a chronic source of
engineering error that no amount of correct calculation downstream
will recover.

We name the discipline. Before reading the curve, the engineer reads
the frame: the axes, the units, the scale, the origin, the legend,
and the source. Only then does the engineer read the data.

\subsection{Axes: linear, log, or something else}

The first question on any plot is the type of each axis. Linear axes
mark equal differences with equal spacing. Logarithmic axes mark
equal ratios with equal spacing: the distance from $1$ to $10$ equals
the distance from $10$ to $100$. A log axis is identifiable by gridlines
spaced at $1, 2, 3, \ldots, 9, 10, 20, 30, \ldots$, by tick labels at
$10^{n}$, or by an explicit annotation. A plot whose vertical axis
runs from $0$ to $10^{6}$ in five equal divisions is linear; the same
range labelled $1, 10, 100, 1000, 10\,000, 100\,000, 10^{6}$ at equal
divisions is logarithmic.

% Figure: how a logarithmic axis is gridded, and how the same data
% point reads differently against a linear and a log axis. The left
% panel shows the uneven minor gridlines of a single decade; the
% right panel contrasts a linear and a log vertical axis carrying the
% same three data values.
\begin{figure}[htbp]
\centering
\resizebox{\linewidth}{!}{%
\begin{tikzpicture}[scale=1.0,
  ax/.style={->, draw=engblack, thick},
  major/.style={draw=engblack, thick},
  minor/.style={draw=engblack!45, thin},
  pt/.style={circle, fill=engred, inner sep=1.4pt},
  ann/.style={font=\footnotesize, align=center},
]
% --- Left: anatomy of one decade on a log axis ---
\begin{scope}[xshift=0cm]
  \draw[ax] (0, 0) -- (0, 6.2) node[left, ann] {log axis};
  % decade from 10^0 to 10^1: minor lines at log10(2..9)
  \foreach \v in {1,2,3,4,5,6,7,8,9} {
    \pgfmathsetmacro{\yy}{4*ln(\v)/ln(10)}
    \draw[minor] (-0.08, \yy) -- (0.4, \yy);
    \node[ann, right, xshift=2pt] at (0.4, \yy) {\v};
  }
  % major at 10^0, 10^1, and partial 10^2
  \draw[major] (-0.15, 0) -- (0.5, 0);
  \node[ann, left, xshift=-2pt] at (-0.15, 0) {$10^{0}$};
  \draw[major] (-0.15, 4) -- (0.5, 4);
  \node[ann, left, xshift=-2pt] at (-0.15, 4) {$10^{1}$};
  \foreach \v in {2,3,4,5,6,7,8,9} {
    \pgfmathsetmacro{\yy}{4 + 4*ln(\v)/ln(10)}
    \ifdim \yy pt < 6.0pt
      \draw[minor] (-0.08, \yy) -- (0.4, \yy);
    \fi
  }
  \node[ann, anchor=south] at (0.6, 6.0) {(a) one decade,};
  \node[ann, anchor=south] at (0.6, 5.6) {minor gridlines};
\end{scope}

% --- Right: same data on linear vs log vertical axis ---
\begin{scope}[xshift=6cm]
  % linear axis
  \draw[ax] (0, 0) -- (0, 6.2) node[left, ann] {linear};
  \foreach \v/\yy in {0/0, 200/1.2, 400/2.4, 600/3.6, 800/4.8, 1000/6.0} {
    \draw[minor] (-0.08, \yy) -- (0.08, \yy);
    \node[ann, left, xshift=-2pt] at (-0.08, \yy) {\v};
  }
  % three data points at 5, 50, 950 on linear scale (scaled /1000*6)
  \foreach \v in {5, 50, 950} {
    \pgfmathsetmacro{\yy}{\v/1000*6}
    \node[pt] at (0, \yy) {};
  }
  \node[ann, anchor=south] at (0.3, 6.0) {(b) linear};

  % log axis
  \draw[ax] (4, 0) -- (4, 6.2) node[left, ann] {log};
  \foreach \e/\yy in {0/0, 1/2, 2/4, 3/6} {
    \draw[minor] (3.92, \yy) -- (4.08, \yy);
    \node[ann, left, xshift=-2pt] at (3.92, \yy) {$10^{\e}$};
  }
  % same three data points 5, 50, 950 on log scale (log10(v)*2)
  \foreach \v in {5, 50, 950} {
    \pgfmathsetmacro{\yy}{2*ln(\v)/ln(10)}
    \node[pt] at (4, \yy) {};
  }
  \node[ann, anchor=south] at (4.3, 6.0) {(c) log};
\end{scope}
\end{tikzpicture}%
}%
\caption[Anatomy of a log axis and the same data on two scales]{Panel
(a) is the anatomy of a logarithmic axis: within each decade the
minor gridlines sit at $2, 3, \ldots, 9$, crowded toward the top of
the decade because the gaps shrink as $\log_{10}$ of consecutive
integers. The signature of a log axis is this uneven crowding, not
the even spacing of a linear axis. Panels (b) and (c) carry the same
three data values, $5$, $50$, and $950$, on a linear axis and a log
axis. On the linear axis the two small values pile up
indistinguishably near the origin and only the largest is readable;
on the log axis all three separate cleanly. The choice of axis
decides which structure the eye can see.}
\label{fig:vol02:ch02:log-axis-reading}
\end{figure}


\Cref{fig:vol02:ch02:log-axis-reading} draws the signature of a log
axis and shows the same three data values on a linear and a log
scale. The crowded minor gridlines within each decade are the visual
tell that distinguishes a log axis from a linear one at a glance.

The mistake the engineer must not make is treating one as the other.
A linear extrapolation off a log axis is a power-law extrapolation in
the underlying quantity; the slope of a straight line on a log-log
plot is an exponent, not a rate. The discipline is to translate
between the two before reading.

\paragraph{Worked example: reading a value off a log axis.}
A datasheet plots leakage current on a vertical log axis with major
gridlines labelled $1\,\si{\nano\ampere}$, $10\,\si{\nano\ampere}$,
$100\,\si{\nano\ampere}$, and a data point sitting two-fifths of the
way up the decade between $10$ and $100\,\si{\nano\ampere}$. The
linear-axis instinct reads ``two-fifths of the way from $10$ to
$100$'' as $10 + 0.4 \cdot 90 = 46\,\si{\nano\ampere}$, and it is
wrong. On a log axis, two-fifths of the way up a decade is
$10^{0.4} \approx 2.5$ times the bottom of the decade, so the value
is $10 \cdot 2.5 = 25\,\si{\nano\ampere}$. The error is nearly a
factor of two and grows with position in the decade: a point
four-fifths of the way up reads as $10^{0.8} \approx 6.3$ times the
decade bottom, that is, $63\,\si{\nano\ampere}$, not the
$82\,\si{\nano\ampere}$ a linear instinct would assign. The rule is
to convert the fractional height to a power of ten before reading,
because position on a log axis is the logarithm of the value, not the
value.

A handful of less common axis types do appear. A \engterm{probability
axis} (used on Q-Q and probability plots) compresses the central
region and stretches the tails so that a normal distribution plots as
a straight line. A \engterm{reciprocal axis} (used in Arrhenius
plots) is linear in $1/T$ and lets activation energies read as
slopes. A \engterm{Smith chart} maps complex impedances onto a
circular grid. Each comes with its own discipline; the only universal
rule is to identify the axis transformation before reading the curve.

\subsection{Slope}

The slope of a curve at a point is the rate at which the vertical
quantity changes per unit change in the horizontal quantity. On
linear axes, slope reads as $\Delta y / \Delta x$ in whatever units
the axes carry. On a semilog plot (vertical log), slope reads as
$\Delta (\log_{10} y) / \Delta x$; the underlying rate is the slope
divided by $\log_{10} e$. On a log-log plot, slope is the exponent of
the power law that fits the data locally.

A flat region indicates that the quantity is insensitive to the
input over that range. A steep region indicates a regime where small
changes in the input produce large changes in the output, where the
engineer's calibration needs to be best, and where extrapolation is
most dangerous. The engineer who reads slope before fitting a curve
recognises the working regions before applying the formula.

\subsection{Asymptotes and limits at infinity}

An \engterm{asymptote} is a line that the curve approaches without
reaching. Horizontal asymptotes record the limit of the function as
$x \to \pm \infty$: a sensor's saturation level, a population's
carrying capacity, a binding curve's plateau, an enzyme's
$V_{\max}$. Vertical asymptotes record the inputs at which the
function blows up: a resonance frequency, a phase boundary, a pole
of a transfer function, a critical loading.

The engineer reads asymptotes for what they imply about the
underlying physics. A horizontal asymptote at zero suggests
exponential decay; one at a non-zero value suggests a saturation
mechanism; a vertical asymptote suggests a singularity in the model.
Asymptotes also tell the engineer which simplifications are safe
where: a $\sinh(x)/x \to 1$ as $x \to 0$ result lets the engineer
replace the function with $1$ near zero with controlled error, and a
function that approaches an asymptote slowly may be replaced by the
asymptote at the cost of the rate of approach.

\subsection{Scale, origin, and what the eye is being asked to compress}

A graph compresses or amplifies a quantity by its choice of scale.
The engineer reads scale honestly. We collect the patterns the eye is
most easily misled by.

A \engterm{suppressed origin} starts an axis at a value other than
zero. A bar chart whose vertical axis runs from $99\,\%$ to $100\,\%$
visually exaggerates a tiny difference; the same data on an axis
running from $0$ to $100\,\%$ shows the difference as the small thing
it is. Suppression is sometimes legitimate (resolving fine variation
in a quantity that operates near a saturation level) and sometimes
deceptive (selling the size of a marginal effect). The reader's
discipline: name the suppression and re-imagine the plot with the
origin at zero.

A \engterm{broken axis} (a wiggle or pair of slashes through the
axis) marks a discontinuity in scale. Broken axes are honest only
when the break is visible, the segments use the same units, and the
text discusses the choice. An invisible break is a trap.

A \engterm{truncated range} cuts off the data outside a chosen
window. Truncation is engineering when the window is the operating
range and decoration when the window is a selection of flattering
points. The reader's discipline: ask what is outside the plotted
range and why.

\engterm{Units omitted from axis labels} is one of the most common
defects in engineering writing. ``Power'' is not a unit; the reader
needs to know whether the axis runs in watts, kilowatts, or
horsepower. ``Time'' is not a unit; the reader needs seconds,
minutes, hours, days, years. A plot with a numerical axis and no unit
is unreadable, and the only fix is rejection.

\engterm{Asymmetric scaling on a symmetric quantity} flatters one
direction and exaggerates the other. A plot of an error signal whose
positive vertical extent is twice the negative extent makes positive
errors look more important than negative errors of equal size. The
reader's discipline: check that a quantity that is symmetric about
zero is plotted symmetrically.

\engterm{Unstated reference levels} are the cousin of suppressed
origins. A decibel scale needs its reference level (the $20\,\si
{\micro\pascal}$ for sound, the impedance basis for voltage, the full
scale for an ADC); a temperature in degrees needs to say degrees of
what; a percent needs to say percent of what. The graph that omits
the reference level can be misread by a factor of any number at all.

\engterm{Missing uncertainty} is the silent defect. A measured curve
without error bars is a claim that the curve is exactly known. Almost
no measured curve is. The engineer who plots without error bars is
asking the reader to believe the data more than the engineer should.

\subsection{A working catalogue of axis transforms}

Engineering plots draw on a small repertoire of axis transforms
that straighten common functional forms. The catalogue, ordered by
the functional form each straightens, is short enough to memorise
and saves the reader the effort of rederiving the right axes from
first principles for the tenth or hundredth time.

\begin{itemize}[leftmargin=*]
  \item \textbf{Linear--linear.} For functions that are themselves
    linear, or for narrow operating windows where the variation is
    visually first-order. Slope reads as the derivative in the
    units of the axes. Default for measurement work where the
    dynamic range is under one decade.
  \item \textbf{Semilog ($\log_{10} y$ vs $x$).} Straightens
    exponential growth and decay. Slope reads as $k \log_{10} e
    \approx 0.4343 k$ for $y = A e^{kx}$. Default for processes
    governed by $\dd y/\dd x = k y$, that is, radioactive decay,
    population growth in unrestricted media, RC discharge, drug
    clearance, capacitor leakage.
  \item \textbf{Log-log ($\log_{10} y$ vs $\log_{10} x$).}
    Straightens power laws. Slope is the exponent $n$ for $y = A
    x^{n}$. Default for fluid drag versus Reynolds number, $S$-$N$
    fatigue curves, $1/f$ noise, Manning's open-channel flow,
    allometric scaling in biology, Zipf's law in linguistics.
  \item \textbf{Arrhenius ($\ln R$ vs $1/T$).} Straightens
    activation-energy processes. Slope is $-E_{a}/k_{B}$ for
    $R = R_{0} \exp(-E_{a}/k_{B} T)$. Default for chemical
    reaction rates, semiconductor carrier-recombination lifetimes,
    thermistor calibration (Steinhart approximation
    \cite{v2c02-paper:steinhart-hart-1968}), Arrhenius diffusion.
  \item \textbf{Bode ($20 \log_{10} |H|$ vs $\log_{10} f$).}
    Straightens transfer functions into their asymptotic
    straight-line approximations. Slope reads as $\pm 20 n\,\si
    {\deci\bel}$ per decade for a power-law exponent $n$. Default
    for filter gain plots, control-system frequency response,
    audio-equipment specification.
  \item \textbf{Probability axis ($z$-quantile vs $y$).} Straightens
    a normal distribution into a straight line on a Q-Q plot.
    Slope and intercept read as standard deviation and mean.
    Default for residual normality checks, regulatory yield
    plots, reliability Weibull plots (with a related transform).
  \item \textbf{Logit ($\log[y/(1 - y)]$ vs $x$).} Straightens a
    sigmoidal binding or activation curve. Slope reads as the Hill
    coefficient or the activation steepness; intercept reads as
    the half-activation point. Default for dose-response, receptor
    binding, ion-channel activation, logistic regression.
  \item \textbf{Smith chart (impedance on a circular grid).}
    Reads complex impedance as a single point and the reflection
    coefficient as a chord; mismatch as a circular arc. Default
    for RF and microwave engineering.
\end{itemize}

The transform repertoire is the engineer's diagnostic toolbox for
identifying the functional form of unfamiliar data. The procedure:
plot the data on each transform that the dimensional analysis or
the physics suggests; pick the one on which the data falls closest
to a straight line; verify with a residual subplot. The reader who
runs this procedure on every unfamiliar dataset eliminates the
common error of fitting the wrong functional form.

\subsection{The frame, then the data}

We collect the discipline as a procedure the reader should run on
every plot.

\begin{enumerate}[leftmargin=*]
  \item Identify the axes. What quantity, in what unit, on what scale
    (linear, log, reciprocal, probability)?
  \item Identify the origin. Where does each axis start? Is it
    suppressed? Is the suppression legitimate for the question being
    asked?
  \item Identify the range. Is the data the engineer needs inside or
    outside the plotted window?
  \item Identify scale breaks, truncations, and asymmetries. Are they
    annotated?
  \item Identify the data type. Is the curve a measurement, a fit,
    a model, or a schematic? Are uncertainties shown?
  \item Identify the source. Where does the plot come from? Who
    measured? When? On what instrument?
  \item Only then read the data. Slope, shape, asymptote, comparison
    to a reference curve, comparison to the engineer's prior estimate.
\end{enumerate}

The procedure is short. Run it for two minutes on every plot and the
rate of plot-induced engineering error drops by an order of magnitude.

\begin{mastery}
A plot that fails the frame check is not always wrong; sometimes it
is a working sketch made by an engineer who has the frame in their
head and the reader does not. The remedy is to ask. A plot that
fails the frame check in a published document is a different matter:
the writer has the burden of supplying the frame, and the reader has
the right to refuse the data until the frame is supplied. Polya's
discipline of restating a problem before solving it
\cite{gen:polya1945} translates here: restate the plot as words and
units before reading the curve.
\end{mastery}

\section{Case study: characterising a load cell from data to
model to use}\pathtag{standard}

The chapter's vocabulary, domain and image, family and parameters,
composition and inversion, residual and asymptote, all converge on a
single recurring engineering task: taking a device whose input-output
behaviour is known only through measurements and producing a function
the engineer can compute with. We work the task end to end on a
strain-gauge load cell, the sensor that turns force into voltage in
every electronic scale, crane overload monitor, and materials-test
machine. The case study is one problem carried from raw data through
model selection to field use, and the discipline it shows transfers
to any input-output characterisation the reader will face.

\subsection{The device and the raw measurement}

A load cell is a metal body with strain gauges bonded to it in a
Wheatstone bridge. Applied force deforms the body, the deformation
strains the gauges, the strain changes their resistance, and the
bridge converts the resistance change to a small voltage proportional
to the excitation voltage that drives it. The output is reported in
millivolts per volt of excitation (\si{\milli\volt}/\si{\volt}) so
that it does not depend on the supply, a normalisation the engineer
should recognise as dividing out one stage of the composition before
the data is even recorded.

The characterisation experiment loads the cell with a series of
calibrated dead weights from zero to its rated capacity of
$500\,\si{\kilogram}$ and records the bridge output at each step. The
raw data is a column of $(F, V)$ pairs: applied force in kilograms
against output in millivolts per volt. The engineer's first move is
the frame check of Section~2.6, applied to the engineer's own data
before anyone else sees it. The input is force, in kilograms, over
the domain $[0, 500]\,\si{\kilogram}$ set by the rated capacity. The
output is bridge voltage, in millivolts per volt, with an image to be
determined from the data. The measurement is a single ascending load
sequence, so any hysteresis between loading and unloading is not yet
characterised, a gap the engineer records now and closes later.

% Figure: load-cell characterisation case study.
% Left panel: measured output (mV/V) vs applied load with a straight-line
% fit through it; the eye sees a good straight line.
% Right panel: the residual of the straight-line fit vs load, revealing
% a systematic bow (the nonlinearity the straight line hides).
\begin{figure}[htbp]
\centering
\resizebox{\linewidth}{!}{%
\begin{tikzpicture}[scale=1.0,
  ax/.style={->, draw=engblack, thick},
  tick/.style={draw=engblack!60, thin},
  fitline/.style={draw=engblue, thick},
  pt/.style={circle, fill=engred, inner sep=1.1pt},
  zero/.style={draw=engblack!55, thin, dashed},
  ann/.style={font=\footnotesize, align=center},
]

% --- Left: output vs load with straight-line fit ---
\begin{scope}[xshift=0cm]
  \draw[ax] (0, 0) -- (6.4, 0) node[below right, ann] {load $F$ (kg)};
  \draw[ax] (0, 0) -- (0, 4.4) node[left, ann] {output (mV/V)};
  \foreach \x/\xl in {1/100, 2/200, 3/300, 4/400, 5/500} {
    \draw[tick] (\x, -0.05) -- (\x, 0.05);
    \node[ann, below, yshift=-2pt] at (\x, 0) {\xl};
  }
  \foreach \y/\yl in {1/0.8, 2/1.6, 3/2.4, 4/3.2} {
    \draw[tick] (-0.05, \y) -- (0.05, \y);
    \node[ann, left, xshift=-2pt] at (0, \y) {\yl};
  }
  % straight-line fit: output = 0.78 * F/100 (units: 1 chart-unit = 0.8 mV/V)
  \draw[fitline, domain=0:5.4, samples=2] plot (\x, {0.765*\x});
  % data: slightly bowed (parabolic) about the line
  \foreach \x/\y in {0/0.0, 0.5/0.40, 1/0.79, 1.5/1.17, 2/1.55,
    2.5/1.93, 3/2.31, 3.5/2.68, 4/3.05, 4.5/3.42, 5/3.80} {
    \node[pt] at (\x, \y) {};
  }
  \node[ann, anchor=north west, fill=white, inner sep=1pt] at (0.2, 4.2)
    {(a) output vs load,\\straight-line fit};
\end{scope}

% --- Right: residual of straight-line fit (bowed) ---
\begin{scope}[xshift=8.5cm]
  \draw[ax] (0, -1.6) -- (6.4, -1.6) node[below right, ann] {load $F$ (kg)};
  \draw[ax] (0, -2.4) -- (0, 1.6) node[left, ann]
    {residual (\% FS)};
  \foreach \x/\xl in {1/100, 2/200, 3/300, 4/400, 5/500} {
    \draw[tick] (\x, -1.65) -- (\x, -1.55);
    \node[ann, below, yshift=-2pt] at (\x, -1.6) {\xl};
  }
  \foreach \y/\yl in {-1/-0.5, 0/0, 1/0.5} {
    \draw[tick] (-0.05, \y) -- (0.05, \y);
    \node[ann, left, xshift=-2pt] at (0, \y) {\yl};
  }
  \draw[zero] (0, 0) -- (5.4, 0);
  % residual traces a downward bow: r(F) ~ 0.5 - (F-2.5)^2 * scale (sign:
  % straight line over-reads in the middle, so residual dips)
  \foreach \x/\y in {0/0.0, 0.5/-0.55, 1/-0.92, 1.5/-1.10, 2/-1.08,
    2.5/-0.95, 3/-0.78, 3.5/-0.58, 4/-0.30, 4.5/0.10, 5/0.55} {
    \node[pt] at (\x, \y) {};
  }
  \draw[fitline, domain=0:5.0, samples=40, smooth]
    plot (\x, {-0.205*(\x)*(5-\x) + 0.14*\x});
  \node[ann, anchor=north west, fill=white, inner sep=1pt] at (0.2, 1.4)
    {(b) residual reveals\\a systematic bow};
\end{scope}

\end{tikzpicture}%
}%
\caption[A load cell characterised from data]{Characterising a
strain-gauge load cell. Panel (a) plots the measured bridge output in
millivolts per volt of excitation against applied calibration load,
with a least-squares straight line through the eleven points. To the
eye the fit is excellent; a straight-line model with one slope and
one offset looks complete. Panel (b) plots the residual of that fit,
the measured output minus the fitted line, as a percentage of full
scale. The residual is not random scatter about zero; it traces a
systematic bow, the signature of a curvature the straight line
cannot represent. The bow peaks near mid-range at roughly $0.5\,\%$
of full scale, which for a precision cell is many times the
specified linearity. The residual plot is the diagnostic that
promotes the model from a one-parameter line to a two-parameter
quadratic; reading panel (a) alone would have certified a model that
is wrong where it matters.}
\label{fig:vol02:ch02:loadcell-characterisation}
\end{figure}


\subsection{First model: the straight line, and why it is not enough}

A load cell is sold as a linear device, so the first family to try is
the affine function $V = a F + b$, with sensitivity $a$ (the slope,
in \si{\milli\volt}/\si{\volt} per kilogram) and offset $b$ (the
zero-load output, in \si{\milli\volt}/\si{\volt}). A least-squares
straight line through the eleven points recovers the two parameters
in one fit. \Cref{fig:vol02:ch02:loadcell-characterisation}, panel
(a), shows the result, and to the eye the fit is excellent: the
points sit on the line, the slope is well determined, and the offset
reads off the intercept. An engineer who stopped here would report a
sensitivity, an offset, and a clean linear model, and would be wrong
in a way the scatter plot conceals.

The diagnostic that catches the error is the residual plot, panel
(b), which plots the measured output minus the fitted line against
load. If the straight line were the right family, the residual would
be random scatter about zero, the size of the measurement noise. It
is not. The residual traces a systematic bow, negative across the
mid-range and returning toward zero at the ends, peaking near
mid-scale at roughly $0.5\,\%$ of full scale. A systematic shape in
the residual is the signature of a model that has the wrong
functional form, not merely noisy parameters. The bow is the part of
the device's behaviour the straight line cannot represent, and for a
precision cell whose specified nonlinearity might be $0.05\,\%$ of
full scale, a $0.5\,\%$ residual is ten times the budget. The
residual plot promotes the model; reading the scatter plot alone
would have certified a model that is wrong exactly where the device
is used most, in the middle of its range.

\subsection{Second model: adding the quadratic term}

The bow in the residual is a parabola, so the family that absorbs it
is the quadratic $V = a F + b + c F^{2}$, with the new coefficient
$c$ carrying the nonlinearity. Fitting the quadratic to the same data
flattens the residual into random scatter at the noise level, which
is the test that the family is now adequate: a residual with no
remaining structure says the model has captured everything the data
contains beyond noise. The coefficient $c$ is small (the device is
nearly linear) but not negligible, and its sign and size are the
device's nonlinearity specification, the number a datasheet reports
as ``nonlinearity, percent of full scale.'' The engineer now has a
three-parameter model whose residual is honest, and the three
parameters each carry meaning: $a$ is the sensitivity, $b$ the zero
offset, $c$ the curvature.

The estimation habit guards the third parameter. Before trusting $c$,
the engineer asks whether the data supports it. The quadratic term
contributes at most $c F^{2}_{\max}$ at full scale, and if that
contribution is smaller than the scatter in the data, the fit is
reading curvature out of noise and $c$ is not real. Here the bow
peaks at $0.5\,\%$ of full scale while the point-to-point scatter is
a few hundredths of a percent, so the quadratic term is an order of
magnitude above the noise and the third parameter is earned. The same
arithmetic that justifies $c$ would forbid a fourth, cubic, term
whose contribution fell below the scatter: a model carries exactly as
many parameters as the data can resolve above noise, and no more.

\subsection{Using the model: inversion and its domain}

The fitted function maps force to voltage, but the field device does
the opposite: it reads a voltage and must report the force. The
engineer needs the inverse. For the linear model the inverse is
immediate, $F = (V - b)/a$, an affine map with the reciprocal slope.
For the quadratic model the inverse comes from the quadratic formula
applied to $c F^{2} + a F + (b - V) = 0$, taking the root that lies
in the operating range:
\[
  F = \frac{-a + \sqrt{a^{2} - 4c(b - V)}}{2c}.
\]
The inverse exists and is single-valued only where the forward
function is monotone, and the quadratic $V(F)$ is monotone on
$[0, 500]\,\si{\kilogram}$ precisely because the vertex of the
parabola lies outside the operating range. An engineer who chose a
quadratic whose turning point fell inside the range would have built
a sensor that returns the same voltage for two different forces, an
ambiguity no amount of downstream calibration can repair. The
monotonicity check of Section~2.5 is therefore a design constraint on
the characterisation, not a mere algebraic nicety: the model must be
invertible across the whole operating domain, and the engineer
verifies it before the model ships.

The image of the forward map sets the diagnostic boundary. The valid
output range is $[V(0), V(500)]$, the voltages the cell produces
across its rated load. A field reading outside that band cannot have
come from a load in the rated range, and the inverse formula's
discriminant going negative is the algebraic alarm: a voltage so far
out of range that no real force produces it signals a broken bridge,
a disconnected excitation, or a load beyond capacity. The boundary of
the image is the self-test the calibration carries for free.

\subsection{Closing the hysteresis gap}

The single ascending load sequence left one gap: the cell's behaviour
on unloading. Repeating the experiment with a descending load
sequence and overlaying the two curves reveals whether the cell
shows hysteresis, the loading and unloading curves separating into a
loop of the kind Section~2.2 drew. A good load cell holds its
hysteresis below a small fraction of full scale; a cell with
significant hysteresis is not a function of the instantaneous load
and needs either a state-dependent model or a mechanical fix. The
characterisation is complete when the engineer can state four
numbers, sensitivity, offset, nonlinearity, and hysteresis, each
read off a feature of the data, and can compute force from voltage
across the rated range with a known and bounded error. One device,
one dataset, the full vocabulary of the chapter, and a model the
engineer can defend.

\section{Failure: log-scale traps and linear extrapolation}\pathtag{core}

The chapter's failure mode is the misreading of a plot that mixes
linear intuition with logarithmic reality, or that extrapolates
beyond the data the curve was fit on. The pattern is universal across
domains.

\subsection{Linear extrapolation of an exponential process}

A bacterial culture, a viral epidemic, a runaway chemical reaction,
a charging capacitor in the wrong feedback loop, a positive-feedback
oscillation: each is an exponential growth process whose first few
data points, plotted on linear axes, look like a gentle upward curve
that the eye reads as nearly linear. The engineer who extrapolates
that gentle curve linearly is wrong, almost always by orders of
magnitude.

The mechanism is simple. An exponential $y = A e^{kt}$ at small $kt$
has $y \approx A(1 + kt + (kt)^{2}/2 + \cdots)$. At $kt < 0.3$ the
quadratic and higher terms are small and the curve is visually
linear. By $kt = 1$ the value has tripled; by $kt = 3$ the value has
multiplied by twenty; by $kt = 6$ by four hundred. The visual linear
extrapolation from the first regime undershoots the third by factors
that grow without bound.

The remedy is procedural. When the underlying physics is exponential,
the plot should be on semilog axes, where the linear-looking curve
is the actual exponential and a straight-line extrapolation is the
correct one. The engineer who prepares an exponential plot on linear
axes is asking the reader to be fooled.

% Figure: the linear-extrapolation trap on an exponential process.
% Early exponential data looks nearly linear; a linear fit through it
% and an exponential fit through it agree in the fitted window and
% diverge by orders of magnitude outside it.
\begin{figure}[htbp]
\centering
\begin{tikzpicture}[scale=1.0,
  ax/.style={->, draw=engblack, thick},
  tick/.style={draw=engblack!60, thin},
  expo/.style={draw=engred, thick},
  linf/.style={draw=engblue, thick, dashed},
  fitwin/.style={draw=engamber, thin},
  pt/.style={circle, fill=engblack, inner sep=1.2pt},
  ann/.style={font=\footnotesize, align=left},
]
% --- axes ---
\draw[ax] (0, 0) -- (6.5, 0) node[below right, ann] {$t$};
\draw[ax] (0, 0) -- (0, 5.3) node[left, ann] {$y$};
\foreach \x in {1, 2, 3, 4, 5, 6} {
  \draw[tick] (\x, -0.05) -- (\x, 0.05);
  \node[ann, below, yshift=-2pt] at (\x, 0) {$\x$};
}
% fitted window 0..2 shaded by a bracket
\draw[fitwin] (0, -0.25) -- (2, -0.25);
\draw[fitwin] (0, -0.18) -- (0, -0.32);
\draw[fitwin] (2, -0.18) -- (2, -0.32);
\node[ann, anchor=north] at (1, -0.35) {fitted window};
% exponential y = 0.3 e^{0.8 t}, scaled to fit
\draw[expo, domain=0:3.05, samples=50, smooth]
  plot (\x, {0.3*exp(0.8*\x)});
\node[ann, anchor=south east, draw=engred, fill=white, inner sep=2pt]
  at (3.0, 4.9) {true: $y = 0.3 e^{0.8t}$};
% data points in window
\foreach \x in {0, 0.5, 1.0, 1.5, 2.0} {
  \pgfmathsetmacro{\yy}{0.3*exp(0.8*\x)}
  \node[pt] at (\x, \yy) {};
}
% linear fit through the window: slope from (0,0.3) to (2, 1.49)
% slope = (1.49-0.3)/2 = 0.595; line y = 0.3 + 0.595 t
\draw[linf, domain=0:6.3, samples=2]
  plot (\x, {0.3 + 0.595*\x});
\node[ann, anchor=west, draw=engblue, fill=white, inner sep=2pt]
  at (4.2, 2.7) {linear extrapolation};
\end{tikzpicture}
\caption[The linear-extrapolation trap on an exponential]{An
exponential process $y = 0.3 e^{0.8 t}$ (red) sampled over a fitted
window $t \in [0, 2]$ (black points). Across the window the curve is
nearly straight, and a linear fit (blue dashed) tracks it closely; an
analyst who saw only the windowed data would not distinguish the two.
Outside the window the curves diverge without bound: at $t = 3$ the
exponential is already climbing past the top of the frame while the
linear extrapolation predicts a modest value. The visual closeness
inside the fitted window is exactly what makes the linear
extrapolation seductive and wrong, and the remedy is to plot the
process on semilog axes where the exponential is the straight line
and the extrapolation is honest.}
\label{fig:vol02:ch02:extrapolation-trap}
\end{figure}


\Cref{fig:vol02:ch02:extrapolation-trap} draws the trap. Inside the
fitted window the exponential and a linear fit through it are
visually indistinguishable; outside the window they diverge without
bound. The seduction is the closeness inside the window, and the cost
is paid outside it.

\begin{estimation}
\textbf{The cost of linear extrapolation of an epidemic.}

\smallskip
\textit{Estimate, before reading on.}

\smallskip
An outbreak doubles every three days. On day $0$ there are $100$
cases; over the first week the case count is plotted on linear axes
and a planner fits a straight line through the early points,
projecting the load on a hospital with $4000$ beds. Estimate the day
on which the linear projection says the beds fill, and the day on
which the true exponential says they fill.

The true count is $N(t) = 100 \cdot 2^{t/3}$. A linear fit through
the first week catches the average slope over $t \in [0, 7]$:
$N(7) = 100 \cdot 2^{7/3} \approx 100 \cdot 5.04 = 504$, so the
fitted slope is about $(504 - 100)/7 \approx 58$ cases per day. The
linear projection reaches $4000$ beds at $t \approx (4000 - 100)/58
\approx 67\,\si{\day}$. The true exponential reaches $4000$ when
$2^{t/3} = 40$, that is, $t = 3 \log_{2}(40) = 3 \cdot 5.32 =
16\,\si{\day}$. The linear extrapolation is off by a factor of four
in the time to crisis, predicting two months of headroom where the
true process gives a little over two weeks. The error is not a small
miscalibration; it is the difference between an orderly response and
an overrun, and it comes entirely from extrapolating an exponential
process with a linear rule. The corrective is the semilog plot, on
which the doubling reads as a constant slope and the fill date reads
off the straight line without the factor-of-four error. The estimate
ignores the eventual saturation of the epidemic curve (the logistic
turnover Volume II Chapter 7 treats), which bounds the true count but
only after the beds have already filled.
\end{estimation}

\subsection{Power-law extrapolation off the log-log line}

A power law $y = A x^{n}$ on log-log axes is a straight line of slope
$n$. The engineer who fits a power law on log-log axes and
extrapolates the straight line outside the fitted range can be wrong
in two ways. The first is that the power law's exponent changes with
regime: a viscous-drag relation that scales as $v^{1}$ at low Reynolds
becomes a pressure-drag relation that scales as $v^{2}$ at high
Reynolds, and the transition is invisible on a log-log plot fitted
only in one regime. The second is that the power law is approximate:
the real curve has multiplicative corrections that read as small on
the log plot but compound across decades.

The remedy is to plot the data over the widest range available, fit
in the regime of interest, and extrapolate only after stating the
regime boundaries. A working engineer treats the slope of a log-log
fit as a working number for the fitted regime, not as a universal
exponent.

\subsection{Suppressed origins and selection bias: Challenger,
1986}

The pre-launch decision teleconference for Space Shuttle mission
STS-51-L on the night of 27--28 January 1986 included a chart prepared
by Morton Thiokol engineers showing the count and severity of O-ring
erosion incidents on prior Shuttle flights. The chart plotted only
the flights on which O-ring damage had been observed, sorted along
the horizontal axis. The flights on which no damage had been
observed (a larger set, including launches at higher temperatures
where the O-ring was performing as designed) were not plotted. Visual
inspection of the chart showed no clear trend with temperature; the
launch decision proceeded \cite{acc:nasa-challenger}.

The Rogers Commission report identified the launch decision and the
data presentation as a critical failure point: ``the engineers' technical
concerns had been correctly recognised at the time but were not
preserved through the decision chain''
\cite{acc:nasa-challenger}. Replotting the same data with all
twenty-four prior flights on the horizontal axis, sorted by ambient
temperature, shows the trend the night-before chart concealed: every
launch below about $18\,\degC$ had shown some O-ring distress, no
launch above about $24\,\degC$ had. The launch ambient temperature on
28 January 1986 was approximately $-2 \,\degC$, well below any prior
launch. The vehicle broke up at $T + 73 \,\si{\second}$ after the
right solid rocket booster's aft field joint failed to seal at low
ambient temperature.

% Figure: the Challenger pre-launch chart in two versions. The top
% panel shows the night-before chart selecting only flights on which
% O-ring damage occurred. The bottom panel replots the same data
% over all 24 prior flights, sorted by ambient temperature, with the
% STS-51-L launch ambient marked.
\begin{figure}[htbp]
\centering
\resizebox{\linewidth}{!}{%
\begin{tikzpicture}[scale=1.0,
  ax/.style={->, draw=engblack, thick},
  tick/.style={draw=engblack!60, thin},
  damaged/.style={circle, fill=engred, draw=engred, inner sep=1.6pt},
  cleanpt/.style={circle, fill=white, draw=engblack, thick, inner sep=1.5pt},
  alarm/.style={draw=engred, dashed, thick},
  ann/.style={font=\footnotesize, align=left},
  caption/.style={font=\small\bfseries, align=left},
]

% --- Top panel: the night-before chart (incidents only) ---
\begin{scope}[yshift=5.5cm]
  \draw[ax] (0, 0) -- (8.3, 0)
    node[below right, ann] {flight (sorted as presented)};
  \draw[ax] (0, 0) -- (0, 3.3) node[left, ann] {damage index};
  \foreach \y in {1, 2, 3} {
    \draw[tick] (-0.05, \y) -- (0.05, \y);
    \node[ann, left, xshift=-2pt] at (0, \y) {\y};
  }
  % Seven incident flights placed across the horizontal axis,
  % sorted by the order Morton Thiokol presented; damage index
  % values from the engineering chart's tally.
  \foreach \x/\y/\flt in
    {0.8/2/{STS-2}, 1.8/1/{41-B}, 2.8/2/{41-C}, 3.8/3/{51-C},
     4.8/2/{61-A}, 5.8/1/{61-C}, 6.8/2/{41-D}} {
    \node[damaged] at (\x, \y) {};
    \node[ann, below, yshift=-3pt] at (\x, 0) {\flt};
  }
  \node[caption] at (4.0, 3.2)
    {(a) chart as presented: incidents only};
  \node[ann, anchor=west, fill=white, inner sep=2pt]
    at (5.1, 2.6)
    {no apparent\\trend};
\end{scope}

% --- Bottom panel: corrected chart (all 24 flights vs temperature) ---
\begin{scope}[yshift=0cm]
  \draw[ax] (0, 0) -- (10.2, 0)
    node[below right, ann] {launch ambient temperature ($\degC$)};
  \draw[ax] (0, 0) -- (0, 3.3) node[left, ann] {damage index};
  \foreach \xc/\xl in {1/-5, 2/0, 3.5/5, 5/10, 6.5/15, 8/20, 9.5/25} {
    \draw[tick] (\xc, -0.05) -- (\xc, 0.05);
    \node[ann, below, yshift=-2pt] at (\xc, 0) {\xl};
  }
  \foreach \y in {1, 2, 3} {
    \draw[tick] (-0.05, \y) -- (0.05, \y);
    \node[ann, left, xshift=-2pt] at (0, \y) {\y};
  }
  % Incident flights at cool temperatures
  \node[damaged] at (4.4, 3) {};
  \node[ann, anchor=west, xshift=2pt] at (4.4, 3) {51-C (12$\degC$)};
  \node[damaged] at (6.0, 2) {};
  \node[damaged] at (6.2, 2) {};
  \node[damaged] at (6.4, 1) {};
  \node[damaged] at (6.6, 1) {};
  \node[damaged] at (6.8, 2) {};
  \node[damaged] at (7.0, 2) {};
  % Clean flights at warm temperatures
  \node[cleanpt] at (7.6, 0.25) {};
  \node[cleanpt] at (7.8, 0.25) {};
  \node[cleanpt] at (8.0, 0.25) {};
  \node[cleanpt] at (8.2, 0.25) {};
  \node[cleanpt] at (8.4, 0.25) {};
  \node[cleanpt] at (8.6, 0.25) {};
  \node[cleanpt] at (8.8, 0.25) {};
  \node[cleanpt] at (9.0, 0.25) {};
  \node[cleanpt] at (9.2, 0.25) {};
  \node[cleanpt] at (9.4, 0.25) {};
  % STS-51-L predicted ambient ~ -2 C: mark as vertical line
  \draw[alarm] (1.6, 0) -- (1.6, 3.0);
  \node[ann, anchor=south, draw=engred, fill=white, inner sep=2pt]
    at (1.6, 3.2)
    {STS-51-L launch\\ambient $\approx -2\,\degC$};
  \node[caption] at (5.0, 3.2)
    {(b) corrected: all 24 flights vs temperature};
  \node[ann, anchor=south west, fill=white, inner sep=2pt]
    at (3.6, -1.4)
    {filled red: O-ring distress observed.\\
     open black: no O-ring distress observed.\\
     replotted from
     \cite{acc:nasa-challenger}, vol.~I.};
\end{scope}

\end{tikzpicture}%
}%
\caption[The Challenger pre-launch chart: incidents only vs all
flights]{The pre-launch teleconference for STS-51-L on 27--28
January 1986 was shown a chart in the form of panel (a): only the
flights on which O-ring damage had been observed appeared, sorted
along the horizontal axis without temperature. The eye reads no
trend. Panel (b) is the same data replotted with all twenty-four
prior Shuttle flights present, sorted by launch ambient temperature.
The trend is unmistakable: every observed instance of O-ring distress
sits at $18\,\degC$ or below; no launch above $24\,\degC$ has shown
distress. The vertical dashed line marks the predicted ambient at
launch on 28 January 1986, $\approx -2\,\degC$, well to the left of
any prior data. The defect in (a) is selection on the dependent
variable: a plot of incidents that omits the non-incidents is
unreadable for the question of whether incidents correlate with the
explanatory variable. The corrected plot (b) is one of the most
widely reproduced examples in the engineering-communication
literature; the chart-preparation lesson is narrow and clean. The
broader organisational, material, and procedural lessons of the
Challenger loss are developed in their proper depth in Volume X
chapters~9 and 11.}
\label{fig:vol02:ch02:challenger-chart}
\end{figure}


\Cref{fig:vol02:ch02:challenger-chart} sets the two versions of the
chart side by side. Panel (a) is the form actually presented during
the night-before teleconference, with only the seven flights that
had shown O-ring damage on the horizontal axis. The eye reads no
trend. Panel (b) is the same data replotted with all twenty-four
prior Shuttle flights present, sorted by launch ambient temperature;
the trend is unmistakable, and the predicted launch ambient for
STS-51-L sits to the left of any prior datum.

The data file
\texttt{data/challenger\_oring\_flights.csv} carries the twenty-four
prior flights with their launch dates, ambient temperatures, and
damage indices, sourced from the Rogers Commission report; the
Failure-analysis exercises invite the reader to reproduce the
corrected chart from this file. The defect Tufte's \emph{Visual
Display of Quantitative Information}
\cite{v2c02-text:tufte2001} names ``selection on the dependent
variable'' is the one the chart presents, and the corrected version
appears in Tufte's later treatment of the case as a paradigm of
how the same data, plotted differently, supports opposite
decisions.

The graph-reading lesson is narrow and clean. A plot of incidents
against an explanatory variable that omits the non-incidents is
unreadable for the question of whether the incidents correlate with
the variable. Selection on the dependent variable is one of the
oldest defects in empirical engineering, and it lives in chart
preparation as much as in study design. The engineer who would have
caught the defect is the engineer who runs the seven-step frame check
on every plot before reading the data.

The Challenger case is properly studied across multiple chapters of
this work for its material, organisational, and ethical mechanisms;
Volume X chapters 9 and 11 develop the launch-decision teleconference
and the engineering recommendation reversal in their proper depth.
Here we use it only for the chart-preparation defect.

\subsection{Truncation: Mars Climate Orbiter and Air France 447,
revisited}

Two cases briefly revisited from Volume I. The Mars Climate Orbiter
spacecraft was lost on 23 September 1999 because impulse data crossed
a software interface in pound-seconds while the receiving program
expected newton-seconds, a factor of $4.45$ in the trajectory
correction \cite{acc:nasa-mco-mib}. The display of the trajectory
data, had it been read on a plot whose units were stated, would have
shown the offset. Truncation of axis labels (units omitted) was part
of the failure path.

Air France Flight 447 was lost on 1 June 2009 in cruise over the
Atlantic after pitot probes iced and gave inconsistent airspeed
indications. The flight-deck display of airspeed crossed an interface
between sensors and pilots without a scale annotation that
distinguished a genuine airspeed loss from a sensor failure
\cite{acc:bea-af447}. The graph-reading discipline at issue is the
same: a quantity displayed without provenance, scale, and uncertainty
cannot be read safely under time pressure.

\subsection{Cross-domain: the same defects across industries}

The chart-preparation defects of Section~2.6.4 are not specific to
aerospace. The pattern recurs across industries, and the working
engineer who carries the vocabulary can identify the defect on
unfamiliar plots from outside their domain. A short tour, by
defect.

\engterm{Suppressed origins} live in financial dashboards (a stock
chart with a vertical axis from \$$172.50$ to \$$173.50$ makes a
$0.6\,\%$ move look like a cliff), in election-night displays (a
bar chart from $48\,\%$ to $52\,\%$ vote share visually amplifies
a tied race), and in news graphics generally. The engineering
discipline is the same as the political one: state the suppression
and re-imagine the chart with the origin at zero before forming a
judgement on the size of the effect.

\engterm{Units omitted from axis labels} live in datasheet curves
for almost every electronic component (``output current'' without
specifying milliamps or microamps), in manufacturing-yield reports
(``defect rate'' without specifying per million, per thousand, or
per hundred), and in benchmark performance plots (``operations''
without specifying integer or floating, single-thread or aggregate).
The Mars Climate Orbiter case is the spaceflight cousin
\cite{acc:nasa-mco-mib}; the everyday cousin is on the datasheet
of every $\si{\mega\ohm}$ resistor on the engineer's desk.

\engterm{Selection on the dependent variable} lives in clinical
trial outcome reporting (publishing only the positive trials and
filing the negative ones), in mutual-fund performance brochures
(the funds that survived to print are not a random sample of the
funds that were launched), and in any retrospective study that
samples cases by outcome rather than by population. The Challenger
chart is the engineering version of a problem that biostatisticians
named selection bias and that publication-bias reformers named
the file-drawer problem; the discipline of plotting the
non-incidents alongside the incidents is the same in every field.

\engterm{Missing uncertainty} lives in single-instrument readouts
(a fuel gauge that reads $40\,\%$ does not say $\pm 5\,\%$), in
sensor-network displays (each individual reading has its accuracy
spec but the aggregate display does not), and in machine-learning
benchmark plots (a single accuracy number on a held-out set
without confidence bounds is the silent default of much of the
published machine-learning literature). The discipline is the
same: a measurement reported without its uncertainty is a claim
that the measurement is exactly known.

\engterm{Asymmetric scaling on a symmetric quantity} lives in
audio-equipment plots (a noise floor stated in $\si{\deci\bel}$ on
a scale that runs further below zero than above gives the
impression that the noise floor is more important than the
positive headroom), in error-signal plots (positive errors and
negative errors of equal magnitude treated differently by a
non-symmetric chart layout), and in financial volatility plots
(implied volatility for downside and upside not plotted on the
same scale). The discipline is to plot symmetric quantities
symmetrically and to note any deliberate departures from symmetry
in the caption.

The point of the cross-domain tour is not that engineers should
read financial newspapers; it is that the chart-preparation
vocabulary the chapter develops applies to every plot the engineer
will see across a career, and the same defect lives behind many
different industrial errors. The reader who carries the vocabulary
of Section~2.6.4 has a transferable diagnostic tool.

\begin{principle}[Read the frame before reading the curve]
Every plot crossed at an engineering interface, whether on paper, on
screen, in a slide deck, in a control-room display, or in an
operating-system dashboard, carries its axes, units, scales,
origins, ranges, and uncertainties explicitly. The reader has the
right to refuse a curve until the frame is supplied. The writer has
the obligation to supply it. The discipline is the same one Volume I
named for units at every interface, applied here to graphical
representation.
\end{principle}

The principle has the same status as the unit-carrying principle
from Volume I Chapter 2: it is a habit the reader will use in every
chapter that follows, and the rate of plot-induced engineering error
drops sharply when it becomes automatic.

\section{Functions as input-output: the bridge forward}\pathtag{standard}

A function read as input-output is the central object of the rest of
the volume and of much of engineering. A system takes an input
signal, applies a rule, and produces an output signal; the rule is a
function, and the engineer's whole working description of the system
is the description of that function. The vocabulary of this chapter,
domain, range, composition, asymptote, scaling, is the vocabulary in
which input-output behaviour is written.

The simplest input-output function is a constant gain, $y = K x$:
double the input, double the output. A gain is a linear function
through the origin, and its graph is a line of slope $K$. Most real
systems are not pure gains. A sensor saturates (the piecewise
saturator of Section~2.2), so its input-output function is linear
near the origin and flat at the rails. An amplifier has a frequency
response (the rational transfer function of the low-pass worked
example), so its gain is a function of the input frequency, read off
the Bode plot. A control loop composes the plant function with the
controller function (the composition of Section~2.1), so its overall
behaviour is the composed map.

\engterm{Transfer function} is the name for the input-output rule of
a linear time-invariant system, written as a rational function
$H(s) = N(s)/D(s)$ in a complex frequency variable $s$. The chapter
has already met its parts: the zeros of $N$ are the input frequencies
the system rejects, the poles (the roots of $D$) are the input
frequencies at which the system resonates or goes unstable, and the
ratio of leading coefficients sets the high-frequency gain. Chapter 3
builds the complex-exponential machinery that makes the transfer
function computable; Volume IX uses it to design controllers. The
function-and-graph reading discipline of this chapter is the
foundation both rest on.

\paragraph{Worked example: an input-output map read three ways.}
A loudspeaker driver has, over its working band, an input-output
relation that an engineer can describe as a function of three
different variables, and each description answers a different
question. As a function of input voltage at a fixed frequency, the
cone displacement is approximately linear at small signal and
compresses at large signal, a saturating piecewise function whose
flat region warns of clipping. As a function of frequency at fixed
voltage, the displacement is the driver's frequency response, a
rational function with a resonance peak (a pole pair) at the driver's
natural frequency and a roll-off above it. As a function of time at a
fixed input tone, the displacement is a sinusoid whose amplitude and
phase are read off the frequency-response function at that tone. One
physical device, three input-output functions, three reading
disciplines, all drawn from this chapter. The engineer chooses the
description that answers the question at hand, and reads the chosen
function with the frame-check habit the chapter trains.

\section{Project}\pathtag{core}

\begin{project}[Read ten plots from the engineering literature]
The reader selects ten plots from primary engineering literature
and produces a diagnostic write-up identifying what each plot's axes
hide, plus a 1500-word reflection on the patterns that recur.

\textbf{Track}. Analysis only. No physical work; no instruments. The
reader works with published plots and a pencil.
\textbf{Hazard class}: none.
\textbf{Tools}. Ten plots from primary engineering literature
selected by the reader from textbooks, journal articles, standards
documents, or industry reports. Paper, pencil. Optionally a plotting
tool to reproduce one plot at a corrected scale.

\textbf{Procedure}. For each of the ten plots, run the seven-step
frame-check procedure of Section~2.6.5. For each plot, the reader
records:

\begin{enumerate}[leftmargin=*]
  \item Source citation, with year and page or figure number.
  \item Each axis: quantity, unit, scale type (linear, log,
    reciprocal, probability, other), and origin location.
  \item Range plotted, and any data outside the plotted range that
    the reader can identify from the surrounding text.
  \item Scale breaks, truncations, asymmetries, and whether they are
    annotated.
  \item Whether the data is a measurement, a fit, a model, or a
    schematic; whether uncertainties are shown.
  \item One specific defect or trap the plot contains, named with
    the vocabulary of Section~2.6.4 (suppressed origin, units omitted,
    truncated range, asymmetric scaling, unstated reference, missing
    uncertainty, scale break, selection on the plotted variable).
  \item One sentence on how the plot would be read differently if the
    defect were corrected.
\end{enumerate}

The reader is encouraged to choose plots from a range of sources:
mechanical or electrical datasheets; civil-engineering codes (loads
on a structure); biology, chemistry, or pharmacology figures; control
or signal plots from a process industry; an industrial dashboard
screenshot; an academic paper from outside the reader's primary
domain. Variety teaches the patterns better than ten plots from a
single subdiscipline.

\textbf{Deliverable}. A 10-plot diagnostic write-up of approximately
two to three pages of structured tables or notes (one page per
plot), plus a 1500-word reflection addressing:

\begin{itemize}[leftmargin=*]
  \item Which defect recurred most often across the ten plots.
  \item Which defect was hardest to detect, and which was easiest.
  \item Which plot the reader would correct first, and how.
  \item One plot from the reader's own work or notebooks that
    contains a defect from the same vocabulary, identified honestly.
  \item One graph-reading habit the reader has acquired across this
    chapter that they did not have before.
\end{itemize}

\textbf{Time}. Twenty minutes per plot for the diagnostic, plus two
to three hours for the reflection. Total approximately five to seven
hours.

\textbf{Phases}. The project runs in four phases.
\emph{Selection} (30 minutes): the reader assembles ten candidate
plots from at least four distinct sources, with at most three from
any one source. A working selection mixes a textbook figure, a
journal paper figure, a standards-document figure, a manufacturer
datasheet plot, and a screenshot from an operating dashboard.
\emph{Diagnostic pass} (3--4 hours): twenty minutes per plot
following the seven-step frame check. Each plot's record is a
half-page table or a tight paragraph.
\emph{Cross-plot reflection} (2 hours): the 1500-word reflection
collecting patterns across the ten plots, the answer to the prompt
questions, and a single honest self-criticism from the reader's own
work.
\emph{Single redraw} (1 hour, optional): the reader picks one of the
ten plots and redraws it with the defect corrected, in any plotting
tool they prefer; the redraw counts as evidence the diagnostic was
honest.

\textbf{Deliverables}. Three files. (1) A diagnostic write-up of
ten plots in the seven-item structure above, one page per plot.
(2) The 1500-word reflection. (3) Optionally, the redraw of one
plot at corrected scale, with a paragraph stating the choice and
the rationale.

\textbf{Success criteria}. The diagnostic is complete when every
plot has all seven frame-check items filled, every defect is named
with the vocabulary of Section~2.6.4, and the reflection answers
all five prompt questions with a specific example for each. The
redraw, when present, shows the same data with the defect
corrected and reads cleanly against a stranger's eye.

\textbf{Common failure modes}. Three predictable ones.
\emph{Selection bias on the diagnostic itself}: the reader picks
ten plots they already know are defective and reports finding
defects, learning nothing the reader did not already believe. The
fix is to include at least two plots the reader expects to pass the
diagnostic cleanly, and to record honestly when they do.
\emph{Verbose write-up}: the reader writes a paragraph per item
when a sentence will do, the diagnostic balloons to thirty pages,
and the reflection runs out of breath. The fix is a tight tabular
format with one row per item.
\emph{Reflection without specifics}: the reader writes a generic
essay on the importance of careful graph design that could have
been written without doing the diagnostic. The fix is to require
that every claim in the reflection cite a specific plot from the
diagnostic pass, by source and page number.
\end{project}

\section{Exercises}

The exercises move from calculation and derivation through estimation
and simulation to design, diagnosis, failure analysis, and judgment.
Calculation and derivation problems have full solutions; open-ended
problems have rubrics or reference write-ups. Exercises are tagged
with their reader-path tier.

\subsection*{Calculation}\pathtag{core}

\begin{exercise}
For $f(x) = (x - 2)(x + 3)(x - 5)$, state the roots, the sign of
$f$ in each of the four intervals they define, and the polynomial's
degree and leading coefficient. Sketch the polynomial without
computing intermediate values.

\paragraph{Solution sketch.}
The roots are $x = -3, 2, 5$. Expanding the leading terms gives
$f(x) = x^{3} + \cdots$, so the degree is $3$ and the leading
coefficient is $+1$. The sign in each interval is set by the sign
of each factor. On $(-\infty, -3)$ all three factors are negative,
so $f < 0$. On $(-3, 2)$ the first two factors have opposite
signs ($x - 2 < 0$, $x + 3 > 0$, $x - 5 < 0$): two negatives, one
positive, $f > 0$. On $(2, 5)$ we have $x - 2 > 0$, $x + 3 > 0$,
$x - 5 < 0$, one negative, $f < 0$. On $(5, \infty)$ all three are
positive, $f > 0$. The cubic enters at $-\infty$, crosses the
$x$-axis at $x = -3$ going up, reaches a local maximum somewhere
in $(-3, 2)$, descends through the root at $x = 2$, reaches a
local minimum in $(2, 5)$, then ascends through the root at $x = 5$
and runs to $+\infty$. The sketch needs only the four interval
signs and the three roots; no intermediate values are required.
\end{exercise}

\begin{exercise}
For the rational function $r(x) = (2x^{2} - 3x - 2)/(x^{2} - 4)$,
identify the zeros, the poles, the horizontal asymptote, and any
removable singularities. State whether the function changes sign
across each pole.

\paragraph{Solution sketch.}
Factor numerator and denominator. The numerator factors as
$2x^{2} - 3x - 2 = (2x + 1)(x - 2)$, zeros at $x = -1/2$ and $x =
2$. The denominator factors as $x^{2} - 4 = (x - 2)(x + 2)$, zeros
at $x = 2$ and $x = -2$. The common factor $(x - 2)$ cancels:
$r(x) = (2x + 1)/(x + 2)$ for $x \neq 2$. The point $x = 2$ is a
\emph{removable} singularity (a hole in the graph at $y = 5/4$).
The remaining pole at $x = -2$ is simple. The remaining zero is at
$x = -1/2$. The horizontal asymptote, from the leading-term ratio
of the original $r$ or equivalently from the simplified form
$(2x + 1)/(x + 2)$, is $y = 2$ (the ratio of leading coefficients).
At the pole $x = -2$, the sign of $r$ changes across the pole (a
simple pole always changes sign), going to $-\infty$ on one side and
$+\infty$ on the other: as $x \to -2^{+}$, $(2x + 1) \to -3 < 0$
and $(x + 2) \to 0^{+}$, so $r \to -\infty$; as $x \to -2^{-}$,
$(x + 2) \to 0^{-}$ and $r \to +\infty$.
\end{exercise}

\begin{exercise}
Convert each of the following between $\ln$ and $\log_{10}$:
(a) $\ln(1000)$, (b) $\log_{10}(e^{5})$, (c) $\log_{10}(2)$ to four
decimal places, (d) $\ln(2)$ to four decimal places. Use the
change-of-base rule and the constants in the chapter.

\paragraph{Solution sketch.}
The change-of-base rule reads $\log_{a}(x) = \ln(x)/\ln(a)$ and
$\log_{10}(x) = \ln(x)/\ln(10)$ with $\ln(10) \approx 2.3026$.
(a) $\ln(1000) = \ln(10^{3}) = 3 \ln(10) \approx 3 \cdot 2.3026
= 6.9078$.
(b) $\log_{10}(e^{5}) = 5 \log_{10}(e) \approx 5 \cdot 0.4343
= 2.1715$, exactly $5/\ln(10)$.
(c) $\log_{10}(2) \approx 0.3010$ (the chapter quotes $0.301$ to
three decimals; the fourth decimal is $0$).
(d) $\ln(2) = \log_{10}(2) \cdot \ln(10) \approx 0.30103 \cdot
2.30259 \approx 0.6931$. Cross-check: $\ln(2)$ is the half-life
to time-constant conversion the chapter cites as $\approx 0.693$.
\end{exercise}

\begin{exercise}
A capacitor charges through a resistor with time constant $\tau =
2.0 \,\si{\milli\second}$. State (a) the half-life of the charging
process, (b) the time at which the capacitor reaches $99\,\%$ of its
asymptotic voltage, (c) the time at which it reaches $99.9\,\%$.
Compare the half-life to the time constant.

\paragraph{Solution sketch.}
The charging is $V(t) = V_{\infty}(1 - e^{-t/\tau})$, so the
fraction unreached is $e^{-t/\tau}$.
(a) Half-life of the deficit: $1/2 = e^{-t/\tau}$ gives $t = \tau
\ln 2 \approx 2.0 \,\si{\milli\second} \cdot 0.693 = 1.39 \,\si
{\milli\second}$.
(b) $V/V_{\infty} = 0.99$ means deficit $0.01 = e^{-t/\tau}$, so
$t = \tau \ln(100) = 2.0 \,\si{\milli\second} \cdot 4.605 = 9.21 \,
\si{\milli\second}$, that is, about $4.6 \tau$.
(c) $V/V_{\infty} = 0.999$ means $t = \tau \ln(1000) = 2.0 \,\si
{\milli\second} \cdot 6.908 = 13.82 \,\si{\milli\second}$, about
$6.9 \tau$.
The half-life $\tau \ln 2 \approx 0.693 \tau$ is shorter than the
time constant: by one time constant the deficit has fallen to
$1/e \approx 37\,\%$, but by one half-life it has only fallen to
$50\,\%$. The half-life is the more honest report for ``how long
until the signal has decayed by a factor of two''; the time
constant is the more honest report for ``the natural timescale of
the exponential.''
\end{exercise}

\begin{exercise}
A bacterial population doubles every $20$ minutes under unrestricted
growth conditions. Starting from $10^{4}$ cells, state the
population at $1, 4, 8, 12$ hours. State the rate constant $k$ in
$\si{\per\hour}$, and the time at which the population reaches
$10^{12}$ cells, the order of magnitude at which growth saturation
typically becomes important in laboratory cultures.

\paragraph{Solution sketch.}
Doubling time $T_{2} = 20\,\si{\minute} = 1/3 \,\si{\hour}$. The
rate constant from $T_{2} = \ln(2)/k$ is $k = \ln(2)/T_{2} =
0.693 / (1/3) = 2.08\,\si{\per\hour}$.
Hourly population: $N(t) = 10^{4} \cdot 2^{t/T_{2}} = 10^{4} \cdot
2^{3t/\text{hour}}$. At $t = 1\,\si{\hour}$, $N = 10^{4} \cdot 2^{3}
= 8 \times 10^{4}$. At $4\,\si{\hour}$, $N = 10^{4} \cdot 2^{12}
\approx 4.1 \times 10^{7}$. At $8\,\si{\hour}$, $N = 10^{4} \cdot
2^{24} \approx 1.7 \times 10^{11}$. At $12\,\si{\hour}$, $N = 10^{4}
\cdot 2^{36} \approx 6.9 \times 10^{14}$.
For $N = 10^{12}$: $2^{3t/\text{hour}} = 10^{8}$, that is,
$3t/\text{hour} \cdot \log_{10}(2) = 8$, so $t = 8 / (3 \cdot
0.301) \cdot \si{\hour} = 8.86 \,\si{\hour}$.
The growth saturation at $\sim 10^{12}$ cells per millilitre is
roughly the order of magnitude at which oxygen, nutrient, or
quorum-sensing factors become limiting in a typical laboratory
shake flask. The exponential model from this exercise is the
small-$N$ asymptote of the logistic curve that Volume II Chapter 7
develops in its proper depth.
\end{exercise}

\subsection*{Derivation}\pathtag{standard}

\begin{exercise}
Show that if $f$ and $g$ are bijective functions with $f: A \to B$
and $g: B \to C$, then $(g \circ f)^{-1} = f^{-1} \circ g^{-1}$.
State the domain of each function appearing in the derivation.

\paragraph{Solution sketch.}
Because $f$ and $g$ are bijective, both $f^{-1}: B \to A$ and
$g^{-1}: C \to B$ exist as functions on their full codomains. The
composition $g \circ f: A \to C$ is bijective (the composition of
bijections is a bijection), so $(g \circ f)^{-1}: C \to A$ exists.
To check that $(g \circ f)^{-1} = f^{-1} \circ g^{-1}$, compose
either way and verify the identity. From the right:
$(g \circ f) \circ (f^{-1} \circ g^{-1}) = g \circ (f \circ
f^{-1}) \circ g^{-1} = g \circ \mathrm{id}_{B} \circ g^{-1} =
g \circ g^{-1} = \mathrm{id}_{C}$. From the left:
$(f^{-1} \circ g^{-1}) \circ (g \circ f) = f^{-1} \circ (g^{-1}
\circ g) \circ f = f^{-1} \circ \mathrm{id}_{B} \circ f =
f^{-1} \circ f = \mathrm{id}_{A}$. Both checks pass, so
$f^{-1} \circ g^{-1}$ is the two-sided inverse of $g \circ f$,
which is unique. The domain of $f^{-1} \circ g^{-1}$ is $C$
(read right-to-left: first $g^{-1}$ takes $C$ to $B$, then $f^{-1}$
takes $B$ to $A$). The order-reversal is the algebra of undoing a
chain of operations: to undo first $f$ then $g$, undo $g$ first.
\end{exercise}

\begin{exercise}
Derive the slope-and-intercept form of the semilog plot for the
exponential function $y = A e^{kx}$. Show explicitly how the slope
on a base-10 semilog plot relates to the rate constant $k$ in the
underlying exponential.

\paragraph{Solution sketch.}
Take base-10 log of both sides: $\log_{10} y = \log_{10}(A e^{kx})
= \log_{10} A + \log_{10}(e^{kx}) = \log_{10} A + kx \log_{10} e$.
On a semilog plot with $\log_{10} y$ on the vertical and $x$ on
the horizontal, this is a straight line with slope $m = k
\log_{10} e \approx 0.4343 k$ and intercept $b = \log_{10} A$.
Conversely, the engineer extracts the rate constant from a fitted
slope by $k = m / \log_{10} e = m \ln(10) \approx 2.303 m$, and
the prefactor from $A = 10^{b}$. The conversion factor $\log_{10}
e$ is the constant the reader should know on sight, and it is the
reason logarithm tables of the nineteenth century quoted base-10
values: a slope-and-intercept fit in any natural base translates
to the universal base-10 reading by multiplication by $\log_{10}
e$.
\end{exercise}

\begin{exercise}
Derive the slope-and-intercept form of the log-log plot for the
power-law function $y = A x^{n}$. Show that the slope is the exponent
$n$ regardless of the base of the logarithm and that only the
intercept changes between base $e$ and base $10$.

\paragraph{Solution sketch.}
Take log of both sides in base $b$ (any positive base $\neq 1$):
$\log_{b} y = \log_{b}(A x^{n}) = \log_{b} A + n \log_{b} x$.
On a log-log plot in base $b$, with $\log_{b} y$ on the vertical
and $\log_{b} x$ on the horizontal, this is a straight line with
slope $n$ and intercept $\log_{b} A$. The slope is the exponent
$n$ regardless of the base of the logarithm: changing base
multiplies both axes by the same constant ($1/\log_{e} b$ or
similar), so the ratio of axis units is preserved. The intercept
changes because $\log_{e} A \neq \log_{10} A$ in general; the
two are related by $\log_{10} A = \log_{e} A / \ln 10$. The
practical consequence: a power-law exponent read off a log-log
slope is the same whether the engineer uses base $e$, base $10$,
or natural log; only the prefactor $A$ has to be back-converted to
the original scale.
\end{exercise}

\begin{exercise}
Show that any rational function $r(x) = p(x)/q(x)$ with $\deg p <
\deg q$ approaches zero as $|x| \to \infty$, and that one with $\deg
p = \deg q$ approaches the ratio of leading coefficients. State the
behaviour when $\deg p > \deg q$ and identify the form of the
asymptote when $\deg p = \deg q + 1$.

\paragraph{Solution sketch.}
Write $p(x) = a_{m} x^{m} + \cdots + a_{0}$ and $q(x) = b_{n} x^{n}
+ \cdots + b_{0}$. Factor out the leading term of each:
$r(x) = \frac{a_{m} x^{m}(1 + \mathcal{O}(1/x))}{b_{n} x^{n}(1 +
\mathcal{O}(1/x))} = \frac{a_{m}}{b_{n}} x^{m - n}(1 +
\mathcal{O}(1/x))$ as $|x| \to \infty$. If $m < n$, the exponent
$m - n < 0$ and the factor $x^{m-n} \to 0$, so $r \to 0$. If $m = n$,
$r \to a_{m}/b_{n}$, a finite horizontal asymptote. If $m > n$, the
factor $x^{m-n}$ grows without bound and $r$ has no horizontal
asymptote. For $m = n + 1$, polynomial long division gives
$r(x) = (a_{m}/b_{n}) x + c_{0} + (\text{lower-order terms})/q(x)$,
where the remainder is a rational function of degree-difference at
least $1$. The first two terms are the oblique asymptote
$y = (a_{m}/b_{n}) x + c_{0}$, and the remainder tends to zero at
infinity. The engineering picture: the leading-term ratio sets the
asymptote, and the lower-order terms are corrections that vanish
at large $|x|$.
\end{exercise}

\begin{exercise}
Show that the function $f(x) = \tanh(x)$ is bijective from $\R$ to
$(-1, 1)$, and derive the closed-form expression for its inverse
$\mathrm{arctanh}(y) = \tfrac{1}{2} \ln\bigl((1+y)/(1-y)\bigr)$.
State the domain of the inverse.

\paragraph{Solution sketch.}
$\tanh(x) = (e^{x} - e^{-x})/(e^{x} + e^{-x})$ is defined on all
of $\R$. Its derivative is $\operatorname{sech}^{2}(x) = 1/\cosh^{2}
(x) > 0$, so $\tanh$ is strictly increasing and hence injective.
As $x \to +\infty$, $\tanh(x) \to +1$; as $x \to -\infty$,
$\tanh(x) \to -1$; and $\tanh(0) = 0$. By continuity and
monotonicity, $\tanh$ surjects $\R$ onto $(-1, 1)$. So $\tanh$ is
a bijection $\R \to (-1, 1)$. To invert, set $y = \tanh(x)$ and
solve for $x$. Let $u = e^{x}$ so $e^{-x} = 1/u$:
$y = (u - 1/u)/(u + 1/u) = (u^{2} - 1)/(u^{2} + 1)$. Rearrange:
$y(u^{2} + 1) = u^{2} - 1$, $u^{2}(1 - y) = 1 + y$, $u^{2} =
(1 + y)/(1 - y)$, so $u = \sqrt{(1 + y)/(1 - y)}$ (positive root,
since $u = e^{x} > 0$). Then $x = \ln u = \tfrac{1}{2} \ln[(1 + y)/
(1 - y)]$. The expression is defined for $y \in (-1, 1)$; the
endpoints $y = \pm 1$ are excluded because the argument of the
log goes to $0$ or $\infty$. Domain of $\mathrm{arctanh}$: $(-1, 1)$.
The function $\mathrm{arctanh}$ appears in special relativity
(rapidity), in particle-physics (the Fisher transformation for
correlation coefficients), and in signal processing (companding
laws).
\end{exercise}

\subsection*{Estimation}\pathtag{core}

\begin{exercise}
Estimate the order of magnitude of the number of grains of rice in a
$1 \,\si{\kilogram}$ bag, then state the result on a base-10
logarithmic scale ($\log_{10}$ of the number). State the assumption
on the mass of one grain and check the result against a published
value if you can find one.

\paragraph{Solution sketch.}
A single long-grain rice grain has mass roughly $25\,\si{\milli
\gram}$ (varies between $20$ and $30\,\si{\milli\gram}$ across
common varieties). One kilogram is $10^{6}\,\si{\milli\gram}$, so
$N \approx 10^{6} / 25 = 4 \times 10^{4}$ grains. On a base-10 log
scale, $\log_{10} N \approx 4.6$. Published values from USDA and
FAO rice-quality datasheets (current as of 2026) put $N$ between
$4.0 \times 10^{4}$ and $5.0 \times 10^{4}$ for long-grain
varieties, consistent with the estimate. The estimate's uncertainty
sits in the grain-mass figure, which a kitchen scale can refine to
the third significant figure by weighing a counted hundred grains.
The order of magnitude, $\log_{10} N \approx 4.6$, is the working
result: the bag contains tens of thousands of grains, not thousands
and not hundreds of thousands.
\end{exercise}

\begin{exercise}
Estimate the time constant for a cup of coffee, initially at
$80 \,\degC$ in a $20 \,\degC$ room, to cool to $40 \,\degC$. State
your assumption on the cooling rate constant in $\si{\per\hour}$ and
state the half-life of the temperature difference. Sketch the
expected curve on linear axes and on semilog axes (vertical
axis: $T - T_{\text{room}}$).

\paragraph{Solution sketch.}
The temperature difference $\Delta T = T - T_{\text{room}}$ obeys
Newton's cooling law $\Delta T(t) = \Delta T_{0} e^{-t/\tau}$, with
$\Delta T_{0} = 60\,\si{\kelvin}$ initially. Cooling to $40\,\degC$
is $\Delta T = 20\,\si{\kelvin}$, that is, $\Delta T/\Delta T_{0}
= 1/3$. Solving $1/3 = e^{-t/\tau}$ gives $t = \tau \ln 3 \approx
1.099\,\tau$. A typical ceramic mug in an unsealed room has time
constant $\tau \approx 0.5\,\si{\hour}$ (assumption: convective
cooling dominates; corresponding rate constant $k = 1/\tau =
2\,\si{\per\hour}$), so cooling from $80\,\degC$ to $40\,\degC$
takes roughly $0.55\,\si{\hour}$ or $33\,\si{\minute}$. The
half-life of the temperature difference is $\tau \ln 2 \approx
0.35\,\si{\hour} \approx 21\,\si{\minute}$. On linear axes the
curve $\Delta T(t)$ is the standard convex decay. On semilog
axes (vertical: $\log_{10} \Delta T$, horizontal: $t$), the curve
is a straight line of slope $-\log_{10}(e)/\tau = -0.4343 / 0.5 =
-0.87\,\si{\per\hour}$ and intercept $\log_{10}(60) \approx 1.78$.
\end{exercise}

\begin{exercise}
Estimate the slope on a log-log plot of the drag force on a sphere
versus the sphere's velocity in three regimes: deep Stokes
($\mathrm{Re} \ll 1$), Newton (constant $C_{D}$, $10^{3} <
\mathrm{Re} < 3 \times 10^{5}$), and supercritical
(post-drag-crisis). State the slope in each regime; sketch the full
curve. The chapter does not derive the curves; the reader uses
qualitative dimensional reasoning from Volume I Chapter 2.

\paragraph{Solution sketch.}
The drag force is $F = \tfrac{1}{2} C_{D} \rho v^{2} A$ in general.
The slope on the log-log plot is determined by how $C_{D}$ depends
on $v$ in each regime.
Deep Stokes ($\mathrm{Re} \ll 1$): $C_{D} \propto 1/\mathrm{Re}
\propto 1/v$, so $F \propto v^{2} \cdot (1/v) = v^{1}$. Slope $= 1$.
Newton (intermediate $\mathrm{Re}$): $C_{D} \approx 0.4$ to $0.5$,
constant. $F \propto v^{2}$. Slope $= 2$.
Supercritical (post-drag-crisis at $\mathrm{Re} \approx 3 \times
10^{5}$): $C_{D}$ drops abruptly from about $0.4$ to about $0.1$
as the boundary layer transitions to turbulence; the drag actually
decreases over a narrow range before continuing to rise. Slope is
negative across the transition window, then returns to roughly $2$.
The sketch: a straight line of slope $1$ across many decades at low
$v$, transitioning through a knee to slope $2$ across the Newton
regime, with a sharp dip and recovery at the drag crisis. The
qualitative shape is the working diagram every wind-tunnel engineer
carries; the quantitative details depend on the geometry (Volume
III Chapter 12).
\end{exercise}

\begin{exercise}
A drug is cleared from the bloodstream with first-order kinetics and
half-life $4.0 \,\si{\hour}$. Estimate the time at which the drug
concentration falls below $5\,\%$ of its initial value, and the time
at which it falls below $1\,\%$. Plot the expected concentration on
linear axes and on semilog axes; identify which axis lets the reader
read the half-life directly.

\paragraph{Solution sketch.}
First-order clearance: $C(t)/C_{0} = 2^{-t/T_{1/2}}$ with
$T_{1/2} = 4.0\,\si{\hour}$.
$5\,\%$ remaining: $0.05 = 2^{-t/4}$, $-t/4 = \log_{2}(0.05) =
\log_{10}(0.05)/\log_{10}(2) = -1.30 / 0.301 = -4.32$, so $t =
17.3\,\si{\hour}$. Equivalently, $0.05 = 1/20$, so $t = 4 \cdot
\log_{2}(20) = 4 \cdot 4.32 = 17.3\,\si{\hour}$.
$1\,\%$ remaining: $0.01 = 2^{-t/4}$, $t = 4 \cdot \log_{2}(100)
= 4 \cdot 6.64 = 26.6\,\si{\hour}$.
On linear axes the curve is the standard convex decay, and the
half-life reads as the time at which the curve crosses $0.5\,C_{0}$,
but only if the eye can see it precisely; for tail values it is
unreadable. On semilog axes the curve is a straight line of slope
$-\log_{10}(2)/T_{1/2} = -0.301/4 = -0.0753\,\si{\per\hour}$, and
the half-life reads off the slope by $T_{1/2} = -\log_{10}(2)/
\text{slope}$. The semilog axis lets the reader extract the
half-life with no ambiguity. This is why pharmacokinetic curves are
plotted semilog in every PK/PD textbook.
\end{exercise}

\begin{exercise}
Estimate the dynamic range, in decades, of (a) a kitchen scale that
reads from $1 \,\si{\gram}$ to $5 \,\si{\kilogram}$, (b) a digital
multimeter that reads from $1 \,\si{\micro\volt}$ to $1000 \,\si
{\volt}$, (c) the human auditory range from $20 \,\si{\hertz}$ to
$20 \,\si{\kilo\hertz}$. State which of the three is most naturally
plotted on log axes.

\paragraph{Solution sketch.}
(a) $5\,\si{\kilo\gram}/1\,\si{\gram} = 5000$, so $\log_{10}(5000)
\approx 3.7$ decades.
(b) $1000\,\si{\volt}/1\,\si{\micro\volt} = 10^{9}$, so $9$
decades.
(c) $20{,}000/20 = 1000$, so $3$ decades.
The multimeter is the obvious case for log axes: nine decades on a
linear scale would put most of the dynamic range invisible at the
bottom. The auditory range, three decades, is borderline; it is
plotted on log axes by convention because human pitch perception
is roughly logarithmic (an octave is a factor of two, and the ear
hears equal-octave intervals as equally spaced). The kitchen scale,
three and a half decades, is usually plotted on linear axes
because the working range over a single weighing is narrow (a
typical recipe-scale operation spans one decade or less). The
working rule: linear axes when the user's measurements span less
than two decades, log axes when they span more than three, judgment
call in between.
\end{exercise}

\subsection*{Simulation}\pathtag{mastery}

\begin{exercise}
Generate $1000$ values of $x$ logarithmically spaced between $10^{-2}$
and $10^{2}$, and plot $y = x^{2}$ on (a) linear-linear axes, (b)
log-log axes, (c) semilog (log-$y$) axes. State the slope of the
curve on the log-log axes and verify it is the exponent in the
power law. State the qualitative shape on the semilog axes and
explain why a power law does not appear straight there.

\paragraph{Solution sketch.}
On linear-linear axes the curve is a parabola opening upward; the
range from $x = 10^{-2}$ to $x = 10^{2}$ runs $y$ from $10^{-4}$ to
$10^{4}$, eight decades; almost all of the curve sits invisibly
flat against the $x$-axis until $x$ approaches $10^{2}$. The plot
is unreadable across the full range.
On log-log axes the curve is a straight line of slope $2$,
intercept $\log_{10}(1) = 0$. Verifying: $\log_{10}(y) =
\log_{10}(x^{2}) = 2 \log_{10}(x)$, so the slope is exactly $2$,
the exponent in the power law.
On semilog (log-$y$, linear-$x$) axes, $\log_{10}(y) = 2 \log_{10}
(x)$ depends on $\log_{10}(x)$, not on $x$ linearly; the curve is
the function $\log_{10}(y) = 2 \log_{10}(x)$ vs $x$, which is the
logarithm of $x$ scaled by $2$, a logarithmic shape. A power law
appears straight only on log-log axes; an exponential appears
straight only on semilog (vertical log) axes. The two are not
interchangeable.
\end{exercise}

\begin{exercise}
Generate noisy exponential data $y = A e^{kt} + \varepsilon$ with
$A = 1$, $k = 0.5$, and Gaussian noise of standard deviation
$\sigma = 0.1 A e^{kt}$ (i.e., $10\,\%$ relative noise). Sample at
$50$ points equally spaced in $t \in [0, 5]$. Fit a straight line
to the data on linear axes and a straight line to $\log y$ on
linear axes; state the recovered $k$ from each fit and identify
which fit is closer to the true value. Comment on why.

\paragraph{Solution sketch.}
The script \texttt{code/semilog\_fit.py} implements both fits and
prints the recovered values. A straight line on linear axes fits
the wrong functional form (linear, not exponential); the recovered
``slope'' has no relation to $k$ except in the small-$t$ Taylor
expansion. A straight line on $\log y$ versus $t$ fits the correct
functional form, since $\log y = \log A + k t$ exactly when the
noise is multiplicative log-normal. With $10\,\%$ relative noise
treated as multiplicative, the log-linear fit recovers $k$ to
within roughly $1$ to $2\,\%$ of the true value across typical
random seeds (true $k = 0.500$; recovered $k$ between $0.495$ and
$0.510$). The linear fit recovers a number that has no direct
physical meaning. The log-linear fit is the right choice for an
exponential under multiplicative noise; a weighted nonlinear least
squares is the right choice when the noise is additive Gaussian.
The distinction is the working content of the second estimation
block in Section~2.3: a fit is only as honest as its noise model.
\end{exercise}

\begin{exercise}
The Hodgkin--Huxley sodium-channel activation curve is approximately
sigmoidal in membrane voltage. Generate a sigmoidal curve $y =
1/(1 + e^{-k(v - v_{0})})$ for $k = 0.2 \,\si{\per\milli\volt}$ and
$v_{0} = -50 \,\si{\milli\volt}$, sampled across $v \in [-100, 0]$
$\si{\milli\volt}$. Plot it on linear axes; plot $\log_{10}(y/(1-y))$
versus $v$ on the same range. State the slope of the second plot and
verify it equals $k \log_{10} e$. Identify the plot on which the
parameter $k$ reads most cleanly.

\paragraph{Solution sketch.}
The logit transform $\log_{10}(y/(1 - y))$ converts the sigmoid
into a straight line. Starting from $y = 1/(1 + e^{-k(v - v_{0})})$,
$1 - y = e^{-k(v - v_{0})}/(1 + e^{-k(v - v_{0})})$, so
$y/(1 - y) = e^{k(v - v_{0})}$. Taking base-10 log gives
$\log_{10}(y/(1 - y)) = k(v - v_{0}) \log_{10} e = 0.4343\, k\,
(v - v_{0})$. The slope on a plot of $\log_{10}(y/(1 - y))$
against $v$ is $k \log_{10} e \approx 0.0869\,\si{\per\milli\volt}$
for $k = 0.2\,\si{\per\milli\volt}$, and the line passes through
zero at $v = v_{0} = -50\,\si{\milli\volt}$. On the linear-axes
sigmoid plot, $k$ can be estimated only through the slope at the
midpoint (which equals $k/4$), so a small error in locating the
midpoint propagates into a large error in $k$. The logit plot lets
the reader fit a straight line over the full data range and recover
$k$ cleanly. The logit transform is the standard tool in receptor
binding, dose-response, and ion-channel activation analysis;
Hill-plot analysis is its multiparameter cousin.
\end{exercise}

\subsection*{Design}\pathtag{standard}

\begin{exercise}
Design a plot for a thermistor's resistance-versus-temperature
characteristic that lets a reader extract both the room-temperature
resistance and the $B$ constant of the thermistor model from the
plot's intercept and slope. State the axis types and the
transformation applied to the resistance.

\paragraph{Discussion.}
The Steinhart approximation $R = R_{0} \exp[B(1/T - 1/T_{0})]$
straightens on the Arrhenius axes: vertical $\ln(R/R_{0})$ (or
equivalently $\ln R$), horizontal $1/T$ in $\si{\per\kelvin}$.
Slope on these axes equals $B$ (in $\si{\kelvin}$); intercept at
$1/T = 1/T_{0}$ (taken as $1/(298.15\,\si{\kelvin})$ for a
$25\,\degC$ reference) is $\ln R_{0}$. \Cref{fig:vol02:ch02:
thermistor-fit} is the worked instance. The reader's plot should
include data spanning at least one decade of resistance, axis
labels with units, the fitted line with its slope ($B$) annotated,
the intercept reading ($\ln R_{0}$), and a residual subplot
underneath showing departures from the two-parameter Steinhart
model (the Steinhart polynomial adds a $C \ln^{3} R$ term that
captures the residual). A rubric for grading: full marks for axes,
units, fit, and intercept; partial marks for missing the residual
plot or the explicit $B$ readout; full marks deducted for a
linear-axes plot that requires the reader to estimate $B$ from a
midpoint slope.
\end{exercise}

\begin{exercise}
Design a plot for the speed-up of a parallel computation against the
number of processors that lets the reader identify whether Amdahl's
law (with a fixed serial fraction) or Gustafson's law (with a fixed
parallel work) better describes the data. State the axis types,
the parameters that read off slope and intercept, and the regimes in
which each model dominates.

\paragraph{Discussion.}
Amdahl's speedup is $S(N) = 1/(s + (1 - s)/N)$ where $s$ is the
serial fraction and $N$ the processor count
\cite{v2c02-paper:amdahl1967}. As $N \to \infty$, $S \to 1/s$, a
ceiling. Gustafson's speedup is $S(N) = N - s(N - 1)$, linear in
$N$ \cite{v2c02-paper:gustafson1988}.
Plot design: horizontal axis $N$ on log scale (covering at least
two decades from $N = 1$ to $N = 100$ or beyond); vertical axis
$S(N)$ on a linear scale. The plot includes (i) the Amdahl curve
for the inferred serial fraction (a saturating curve approaching
$1/s$ asymptotically), (ii) the Gustafson curve (a straight line
of slope $1 - s$ and intercept $s$), and (iii) the measured speedup
data with error bars. The two models predict different shapes at
large $N$: Amdahl saturates, Gustafson keeps rising linearly. The
regime in which each dominates: Amdahl is right when the problem
size is fixed (strong scaling); Gustafson is right when the problem
size grows with $N$ (weak scaling). The reader's plot should label
which regime the measurement targets and state $s$ in both
parameterisations. A working high-performance-computing benchmark
uses both laws side by side as Amdahl-ceiling and Gustafson-line
overlays; the convergence to one or the other diagnoses the
measurement regime.
\end{exercise}

\begin{exercise}
Design a Bode-style plot (gain in $\si{\deci\bel}$ versus frequency,
log-frequency axis) of a first-order low-pass filter with cutoff
$f_{c} = 1 \,\si{\kilo\hertz}$. State the gain at $f \ll f_{c}$, at
$f = f_{c}$, and at $f \gg f_{c}$. State the asymptotic slope above
the cutoff in $\si{\deci\bel}$ per decade, and verify that the
slope corresponds to a power-law exponent of $-1$.

\paragraph{Discussion.}
Magnitude response: $|H(f)| = 1/\sqrt{1 + (f/f_{c})^{2}}$. In
decibels, $G(f) = 20 \log_{10}|H| = -10 \log_{10}(1 + (f/f_{c})^{2})$.
At $f \ll f_{c}$: $G \to 0\,\si{\deci\bel}$ (the passband).
At $f = f_{c}$: $G = -10 \log_{10}(2) \approx -3\,\si{\deci\bel}$.
At $f \gg f_{c}$: $G \approx -20 \log_{10}(f/f_{c}) = -20
\log_{10}(f) + 20 \log_{10}(f_{c})$, a straight line of slope
$-20\,\si{\deci\bel}/\text{decade}$ in the Bode plot.
Slope as a power-law exponent: $-20\,\si{\deci\bel}/\text{decade}$
means a factor of $10$ in frequency gives a factor of $10$ in
$|H|$, that is, $|H| \propto f^{-1}$, exponent $-1$ in the power
law. The plot shows three regimes: a flat passband at $0\,\si
{\deci\bel}$, a knee at $f_{c}$ at $-3\,\si{\deci\bel}$, and an
asymptotic roll-off of $-20\,\si{\deci\bel}/\text{decade}$. A
working filter plot annotates the asymptotes as straight lines
that intersect at $f_{c}$; the Bode-asymptote approximation is the
shortcut every filter designer uses for hand sketches.
\end{exercise}

\begin{exercise}
Design a presentation plot for an experimental error analysis whose
quantity ranges from $0$ to $100\,\%$ but whose dynamic range of
interest is between $99.0\,\%$ and $99.99\,\%$. State the axis
type, the suppression of the origin, the legend, and the annotation
that prevents the reader from misreading the plot as a representation
of the full $0$ to $100\,\%$ range.

\paragraph{Discussion.}
Plot the \emph{complement} on a log axis: vertical axis is the
error fraction $1 - p$ (the ``defect'' or ``leakage'') on a
$\log_{10}$ axis, with tick marks at $10^{-1}, 10^{-2}, 10^{-3},
10^{-4}$. The complement transform spreads the four decades of
interest across the full vertical extent of the plot. The legend
states: ``vertical axis: error fraction $1 - p$ on log scale;
range corresponds to yield $90\,\%$ at top to $99.99\,\%$ at
bottom.'' Annotate the suppression explicitly: a corner banner
reads ``axis suppressed: $p < 90\,\%$ not shown.'' The plot title
states the operating range, for example, ``Yield in
high-reliability operating regime (90\,\% to 99.99\,\%).'' This is
the convention used in semiconductor yield plots, fault-tree
analysis, and reliability-engineering Weibull plots: the engineer
plots the complement on a log axis because the engineering
relevance lives at the small end. A reader who fails to read the
annotation and the legend may misread the plot, but the
discipline of stating the suppression keeps the writer's hands
clean.
\end{exercise}

\subsection*{Diagnosis and reverse engineering}\pathtag{core}

\begin{exercise}
A plot in a draft engineering report shows ``Throughput (units)''
versus ``Time'' with no other annotation. The throughput axis runs
from $0$ to $1000$, the time axis from $0$ to $100$. The curve
ascends linearly from $(0, 100)$ to $(100, 900)$. Identify (a) at
least three frame defects in the plot, (b) the questions you must
answer before trusting the curve, and (c) the minimum annotation a
revision should add.

\paragraph{Solution sketch.}
(a) Three or more frame defects: (1) ``Throughput (units)'' carries
no unit; the reader cannot tell whether the axis runs in items per
second, items per hour, kilograms per minute, or revenue per day.
(2) ``Time'' carries no unit; the reader cannot tell seconds,
minutes, hours, days. (3) The curve is a measurement, a fit, a
projection, or a schematic, and the plot does not say which.
(4) The plot has no source, no instrument, no measurement
condition. (5) The plot has no uncertainty; the reader is asked to
believe the data exactly.
(b) Questions to answer before trusting the curve: what is the
quantity being throughputted, in what unit per what time? Is the
curve a measurement of a single experimental run or an average
over many? What was the operating condition (load, supply, input
signal)? Was the system in steady state at $t = 0$? Why does the
curve start at $100$ rather than $0$ at $t = 0$, and what is the
significance of that intercept?
(c) Minimum revision annotation: axis labels with units; the
nature of the curve (measurement, fit, projection); the source
(instrument, dataset, model); error bars or a shaded uncertainty
band; the operating condition in a caption.
\end{exercise}

\begin{exercise}
A draft figure plots a measured fatigue $S$-$N$ curve on linear axes
with stress on the vertical and number of cycles to failure on the
horizontal. The data span $10^{3}$ to $10^{8}$ cycles. State why the
plot is unreadable, propose the correct axes, and identify the
quantities that should read off slope and intercept under the
proposed transformation.

\paragraph{Solution sketch.}
The horizontal axis spans five decades; on linear axes, the data
between $10^{3}$ and $10^{7}$ cycles is invisibly compressed
against the left edge, and only the $10^{8}$ region is visible.
The plot is unreadable across the range of engineering interest.
Standard practice: plot $\log_{10}(N)$ horizontally (or use a log
axis directly) and $\log_{10}(S)$ vertically, on a log-log plot.
The Basquin form $S = S' (2N)^{b}$ (the working
high-cycle-fatigue model) becomes a straight line on log-log axes
with slope $b$ (the fatigue strength exponent, typically between
$-0.05$ and $-0.15$ for engineering metals) and intercept
$\log_{10}(S')$ at $N = 1/2$ (the fatigue strength coefficient).
The reader extracts $b$ from the slope and $S'$ from the intercept
in two readings. The endurance limit, if present, appears as a
horizontal asymptote at large $N$ and is read directly from the
plot's right edge. A working $S$-$N$ plot has axis units (MPa
vertical, dimensionless count horizontal), the fitted Basquin line
with $b$ annotated, error bars on each data point (each from a
batch of test specimens), and a note on the load ratio $R$ at which
the curve applies. The unreadable linear-axes version is one of
the most common defects in fatigue-life reports drafted by
junior engineers.
\end{exercise}

\begin{exercise}
A plot of $y = A x^{n}$ on log-log axes shows a clean straight line
through five decades of $x$. A junior engineer concludes the
relationship $y = A x^{n}$ holds for all $x$. State the two
extrapolation traps the conclusion ignores, and propose two checks
the engineer should run before extending the fit beyond the plotted
range.

\paragraph{Solution sketch.}
The two traps. (1) Regime change: a power law that holds in one
physical regime may give way to a different exponent at higher or
lower $x$. The viscous-to-pressure drag transition on a sphere
($n$ goes from $1$ to $2$ at $\mathrm{Re} \sim 10$) is a textbook
example. (2) Multiplicative corrections: a power law $y = A x^{n}
(1 + a/x + b/x^{2} + \cdots)$ has the leading-order exponent $n$
but corrections that read as small on the log plot and compound
over decades. Both traps appear together in fluid drag, in
electrical resistivity at low temperature, in fatigue life across
crack-initiation and crack-growth regimes.
Two checks. (1) Plot the data over the widest range available, not
just the fitted regime; look for departures from the straight line
at the endpoints. If both endpoints curve away in the same
direction, the leading-order exponent has changed. (2) Compute the
residuals $\log_{10}(y_{\text{data}}) - \log_{10}(y_{\text{fit}})$
and plot them against $\log_{10}(x)$. A trend in the residuals
across decades indicates a systematic correction the power law
misses. If the residual plot is flat and random, the power law is
honest within the fitted range; if it curves, the extrapolation is
not safe. Both checks are cheap; neither is optional for an
extrapolation an engineer will commit to in writing.
\end{exercise}

\subsection*{Failure analysis}\pathtag{core}

\begin{exercise}
The pre-launch chart presented during the night-before teleconference
for Space Shuttle mission STS-51-L plotted O-ring incidents only on
the flights on which damage had been observed. Argue, in 200--300
words, why this presentation made the temperature dependence of
O-ring failure invisible, and how a corrected chart would have
displayed the data. Cite \cite{acc:nasa-challenger}. State whether
the failure is a chart-preparation failure, an
organisational-decision failure, both, or something else, and which
classification produces the most useful corrective action.

\paragraph{Discussion.}
A passing response treats the chart-preparation defect (selection
on the dependent variable) and the organisational-decision context
as distinct but linked. Strong responses identify that the
incidents-only chart is statistically incoherent for any question
about whether incidents correlate with temperature: the
non-incidents are exactly the data that establish the relationship,
and omitting them makes the correlation invisible
\cite{acc:nasa-challenger}. The corrected chart plots all twenty-
four prior flights' damage index against launch ambient temperature
(\Cref{fig:vol02:ch02:challenger-chart}, panel b). Classification:
both a chart-preparation failure and an organisational-decision
failure. The chart-preparation failure is the proximate cause; the
organisational-decision failure is the structural cause that let
the chart-preparation failure stand. Useful corrective action:
chart-preparation discipline (the seven-step frame check)
addresses the proximate cause and is teachable to every engineer
on a finite timescale. Organisational reform addresses the
structural cause and is much harder to enact reliably. Both are
necessary; chart-preparation discipline is the one this chapter
trains.
\end{exercise}

\begin{exercise}
Construct a working sketch (axes labelled, scale stated, origin
located) of the corrected Challenger pre-launch chart, plotting
ambient temperature on the horizontal axis and O-ring damage
severity on the vertical axis, including the prior twenty-four
flights. State the temperature at which the night-before launch
ambient sat on this chart, and identify the visual cue the corrected
chart would have given the launch teleconference.

\paragraph{Solution sketch.}
Horizontal axis: launch ambient temperature, in $\si{\degreeCelsius}$,
running from $-5\,\si{\degreeCelsius}$ to $30\,\si{\degreeCelsius}$
linearly. Vertical axis: damage index (count of erosion/blow-by
incidents at field and nozzle joints, integer scale $0$ to $3$).
Origin at $(-5\,\si{\degreeCelsius}, 0)$. Data: twenty-four points
read from
\texttt{data/challenger\_oring\_flights.csv}; STS-1 at $(18.9, 0)$,
51-C at $(11.7, 3)$, and so on. The STS-51-L launch ambient was
approximately $-2\,\si{\degreeCelsius}$, well to the left of any
prior datum (the next coldest was 51-C at $11.7\,\si{\degreeCelsius}$).
The visual cue: the cluster of distress points at the cool end and
the cluster of clean points at the warm end give the eye an
unmistakable temperature trend; the predicted launch sits in a
temperature regime with no prior data. \Cref{fig:vol02:ch02:
challenger-chart} panel (b) is the worked instance. The reader's
sketch should at minimum mark the STS-51-L launch ambient with a
labelled vertical line and annotate the cluster of damaged flights
below $\sim 18\,\si{\degreeCelsius}$.
\end{exercise}

\begin{exercise}
A control-room display of a process plant's reactor pressure runs
from $0$ to a full-scale value, with the alarm threshold at $80\,\%$
of full scale. The display shows a curve that has approached the
alarm threshold over the last twenty minutes on linear-time axes.
State the frame defects that would prevent an operator from reading
the rate of approach correctly under time pressure, and propose two
display modifications that would address them.

\paragraph{Discussion.}
Frame defects relevant to time-pressure reading. (1) The vertical
axis covers $0$ to full scale; the operating window of concern is
the top $20\,\%$, so most of the curve sits in a region the
operator does not need to resolve and the slope near the threshold
is visually compressed. (2) Linear-time axes give equal visual
weight to old data and recent data; under time pressure the
operator wants the recent data emphasised. (3) The alarm threshold
is shown as a horizontal line but the projected time-to-alarm,
which is the operationally critical quantity, is not displayed.
(4) There is no annotation of the rate of approach (slope).
Two modifications. (1) Add a second panel showing only the top
$30\,\%$ of the pressure range, with the alarm threshold marked
and the slope (rate of approach) annotated; this is the
\emph{magnified-margin display} convention used in process control.
(2) Compute and display a projected time-to-alarm based on the
linear extrapolation of the recent slope: a clock countdown or a
text annotation ``alarm in approximately X minutes at current
rate.'' Optional third modification: use a logarithmic-time axis
that emphasises recent data over old; this is uncommon in process
control but standard in financial trading dashboards.
\end{exercise}

\subsection*{Judgment}\pathtag{standard}

\begin{exercise}
Argue, in one paragraph, whether the working engineer should default
to log-log axes or to linear axes when first plotting unfamiliar
data. State a specific example from your domain in which each
default would mislead, and propose a procedure that defends against
both errors.

\paragraph{Discussion.}
The honest answer is neither default is universally safe. A passing
response identifies the dynamic-range rule: linear when the data
spans less than one decade, log when it spans more than three,
judgment-call (and worth trying both) in between. Examples that
should appear in a passing response. Linear default misleads when:
plotting bacterial growth across decades on linear axes (the early
exponential phase compresses to invisibility); plotting a Bode plot
on linear-frequency axes (the cutoff knee is invisible). Log
default misleads when: plotting a quantity that fluctuates around
zero, since the log axis cannot represent the zero crossings; or
plotting a quantity whose engineering relevance is the absolute
size of a small effect against a large baseline, where the log
compression makes the effect look unimportant. The defensive
procedure: plot the data on both linear and log axes before any
fitting; pick the axis that lets the eye see the structure most
honestly; state the axis choice and its rationale in the figure
caption. Two minutes of double-plotting at the start of an
analysis saves hours of misreading later.
\end{exercise}

\begin{exercise}
A senior colleague tells the reader that ``error bars are decoration;
the data is the data.'' Write a one-paragraph reply that agrees where
agreement is honest and disagrees where disagreement is honest, with
a specific example from a working engineering domain.

\paragraph{Discussion.}
A passing response distinguishes the cases. Where the colleague is
right: a single measurement made by a calibrated instrument with a
well-characterised uncertainty does not need error bars repeated
on every plot, because the uncertainty is the instrument's stated
specification and a footnote suffices; over-decorating a plot with
identical-size error bars at every point can clutter the
presentation. Where the colleague is wrong: any measurement of a
quantity that varies between repeats, depends on operating
conditions, or carries systematic uncertainty larger than the
instrument's stated precision needs error bars because the data
alone misrepresents the engineer's knowledge of the quantity.
A specific working example: a fatigue-life $S$-$N$ curve drawn
from a batch of test specimens has individual specimen scatter that
can span an order of magnitude in cycles to failure at the same
stress level; the median curve without scatter bars is an actively
misleading representation of the data, and a design that uses the
median curve as if it were exact will fail at the rate of the
scatter the median hid. Error bars are the engineering record of
what the data does not say.
\end{exercise}
